% La lunette astronomique — Physique 1ère TI — Cours signé OnBuch+
% Programme officiel MINESEC. Compiler avec : tectonic N10-lunette-astronomique.tex
\newcommand{\DOCMATIERE}{Physique}
\newcommand{\DOCNIVEAU}{1ère TI}
\newcommand{\DOCMODULE}{Module 3 : Optique}
\newcommand{\DOCLECON}{N10}
\newcommand{\DOCTITRE}{La lunette astronomique}
% OnBuch+ — préambule « cours signés » ENRICHI (compilé avec tectonic / XeLaTeX)
% Système visuel « Pop » d'OnBuch+ : fond crème (jamais blanc), bordures encre
% épaisses, ombres « dures » décalées, titres Archivo Black en capitales.
% Version MATHÉMATIQUES : graphes (pgfplots/tikz), boîtes pédagogiques
% (définition, propriété, méthode, attention, le savais-tu, exercice résolu,
% exercice, TP, bilan), page de garde et signature OnBuch+ ancrée en bas.
%
% Macros attendues dans main.tex AVANT \input{preamble} :
%   \DOCMATIERE  \DOCNIVEAU  \DOCMODULE  \DOCLECON (numéro)  \DOCTITRE
\documentclass[11pt]{article}

\usepackage{fontspec}
\usepackage[french]{babel}
\usepackage[a4paper,margin=2cm,top=3.4cm,bottom=2.8cm,headheight=1.4cm,footskip=1.3cm]{geometry}
\usepackage{amsmath,amssymb,amsfonts}
\usepackage{mathtools} % \xmapsto, \coloneqq... souvent utilisés par les modèles
\usepackage{cancel} % simplifications barrées (bilans de forces)
% \abs{z} (module, valeur absolue) et \norm{u} (norme), utilisés par les modèles
\providecommand{\abs}{}\let\abs\relax\DeclarePairedDelimiter{\abs}{\lvert}{\rvert}
\providecommand{\norm}{}\let\norm\relax\DeclarePairedDelimiter{\norm}{\lVert}{\rVert}
\usepackage[table]{xcolor} % [table] : fournit aussi \rowcolors (alternance de lignes)
\usepackage{pagecolor}
\usepackage{tikz}
\usetikzlibrary{arrows.meta,positioning,calc,shapes.geometric,shapes.misc,shapes.arrows,decorations.pathmorphing,decorations.pathreplacing,decorations.markings,patterns,shadows,mindmap,trees,fit,backgrounds,matrix,intersections,angles,quotes}
\usepackage{pgfplots}
\usepackage[version=4]{mhchem} % équations nucléaires \ce{^{235}_{92}U}
\usepackage[european,siunitx]{circuitikz} % schémas électriques
\pgfplotsset{compat=1.18}
\usepgfplotslibrary{fillbetween,groupplots} % aires entre courbes (calcul intégral)
\usepackage{enumitem}
\usepackage{fancyhdr}
% [most] charge toutes les bibliothèques tcolorbox (listings, minted, raster,
% documentation...) et gonfle inutilement la mémoire TeX sur un document long
% (~300 boîtes) : on ne charge que ce qui est réellement utilisé.
\usepackage[skins,breakable]{tcolorbox}
\usepackage{graphicx}
\usepackage{colortbl}
\usepackage{booktabs}
\usepackage{tabularx}
\usepackage{diagbox} % case d'en-tête diagonale des tableaux à double entrée
% Colonnes extensibles centrée / gauche / droite (C, L, R), souvent utilisées par les modèles
\newcolumntype{C}{>{\centering\arraybackslash}X}
\newcolumntype{L}{>{\raggedright\arraybackslash}X}
\newcolumntype{R}{>{\raggedleft\arraybackslash}X}
\usepackage{array}
\usepackage{multicol}
\usepackage{siunitx}
% parse-numbers=false : accepte aussi \SI{2,5 \times 10^{-3}}{...}
\sisetup{locale=FR,output-decimal-marker={,},group-minimum-digits=4,parse-numbers=false}
\DeclareSIUnit\ohm{\ensuremath{\symup{\Omega}}} % U+2126 absent de la police
\DeclareSIUnit\division{div} % réglages d'oscilloscope (ms/div, V/div)
\DeclareSIUnit\tr{tr} % tours (vitesse de rotation, ex. \SI{1500}{\tr\per\minute})
\DeclareSIUnit\ch{ch} % cheval-vapeur (puissance moteur, unité usuelle hors SI)
\DeclareSIUnit\franccfa{FCFA} % franc CFA (coûts, contexte camerounais)
\DeclareSIUnit\dioptre{\ensuremath{\delta}} % vergence d'une lentille (dioptrie)
\usepackage{qrcode}
% \degree : commande fréquemment utilisée par les modèles pour un angle ou une
% température en dehors de \SI{}{} (non fournie par les paquets chargés).
\providecommand{\degree}{^{\circ}}
% \qed (fin de démonstration), utilisé par les modèles sans amsthm
\providecommand{\qed}{\hfill$\square$}
% \barre : double barre des valeurs interdites dans un tableau de variations (tkz-tab)
\providecommand{\barre}{\ensuremath{\|}}
% environnement proof (amsthm non chargé) utilisé par les modèles
\ifdefined\proof\else\newenvironment{proof}[1][Démonstration]{\par\noindent\textit{#1.}\ }{\qed\par}\fi
\usepackage{lastpage}
\usepackage{hyperref}
\hypersetup{hidelinks}

% ============================================================
% Palette « Pop » OnBuch+ — mêmes hex que la classe P (plus_ui.dart)
% ============================================================
\definecolor{poporange}{HTML}{F59321}
\definecolor{popdark}{HTML}{E07A0C}
\definecolor{popcream}{HTML}{FFF6E8}
\definecolor{popink}{HTML}{1C1714}
\definecolor{poppurple}{HTML}{7A5AE0}
\definecolor{popgreen}{HTML}{2E9E6A}
\definecolor{poppink}{HTML}{E0457B}
\definecolor{popblue}{HTML}{2F7DE1}
\definecolor{popgold}{HTML}{C98A12}
\definecolor{popmuted}{HTML}{8A7F76}
\definecolor{popline}{HTML}{EADFD2}
\definecolor{popgray}{HTML}{8A7F76}
\colorlet{popgrey}{popgray}
% Alias tolérants (le modèle nomme parfois les couleurs différemment)
\colorlet{popwhite}{white}
\colorlet{popblack}{popink}
\colorlet{popred}{poppink}
\colorlet{popyellow}{popgold}
% Teintes claires pour fonds de tableaux / zones
\colorlet{popblueL}{popblue!10!white}
\colorlet{poporangeL}{poporange!14!white}
\colorlet{popgreenL}{popgreen!12!white}
\colorlet{poppurpleL}{poppurple!12!white}
\colorlet{poppinkL}{poppink!12!white}
\colorlet{popgoldL}{popgold!14!white}

\pagecolor{popcream}

% ============================================================
% Typographie
% ============================================================
\defaultfontfeatures{Ligatures=TeX}
\setmainfont{PlusJakartaSans-Medium.ttf}[Path=../../../fonts/,BoldFont=PlusJakartaSans-Bold.ttf]
% Maths en sans-serif (Fira Math) avec chiffres et lettres droites de Plus
% Jakarta, pour que les formules se fondent dans le texte.
\usepackage[math-style=ISO,bold-style=ISO]{unicode-math}
\setmathfont{FiraMath-Regular.otf}[Path=../../../fonts/]
\setmathfont{PlusJakartaSans-Medium.ttf}[Path=../../../fonts/,range={up/num,up/{Latin,latin},\mathup}]
\usepackage{icomma} % virgule décimale sans espace en mode math
% Caractères mathématiques Unicode absents des polices de texte (Plus Jakarta,
% Archivo Black) : rendus par leur équivalent mathématique, en texte comme en
% formule — sinon ils s'affichent en carré vide (ex. « Arithmétique dans ℤ »).
\usepackage{newunicodechar}
\newunicodechar{ℝ}{\ensuremath{\mathbb{R}}}
\newunicodechar{ℤ}{\ensuremath{\mathbb{Z}}}
\newunicodechar{ℂ}{\ensuremath{\mathbb{C}}}
\newunicodechar{ℕ}{\ensuremath{\mathbb{N}}}
\newunicodechar{ℚ}{\ensuremath{\mathbb{Q}}}
\newunicodechar{𝔻}{\ensuremath{\mathbb{D}}}
\newunicodechar{α}{\ensuremath{\alpha}}
\newunicodechar{β}{\ensuremath{\beta}}
\newunicodechar{γ}{\ensuremath{\gamma}}
\newunicodechar{δ}{\ensuremath{\delta}}
\newunicodechar{ε}{\ensuremath{\varepsilon}}
\newunicodechar{θ}{\ensuremath{\theta}}
\newunicodechar{λ}{\ensuremath{\lambda}}
\newunicodechar{μ}{\ensuremath{\mu}}
\newunicodechar{σ}{\ensuremath{\sigma}}
\newunicodechar{τ}{\ensuremath{\tau}}
\newunicodechar{φ}{\ensuremath{\varphi}}
\newunicodechar{ψ}{\ensuremath{\psi}}
\newunicodechar{ω}{\ensuremath{\omega}}
\newunicodechar{Σ}{\ensuremath{\Sigma}}
% majuscules produites par \MakeUppercase dans les titres (α-aminés -> Α)
\newunicodechar{Α}{\ensuremath{\alpha}}
\newunicodechar{Β}{\ensuremath{\beta}}
\newunicodechar{Γ}{\ensuremath{\Gamma}}
\newunicodechar{Δ}{\ensuremath{\Delta}}
\newunicodechar{Ε}{\ensuremath{\varepsilon}}
\newunicodechar{Θ}{\ensuremath{\Theta}}
\newunicodechar{Λ}{\ensuremath{\Lambda}}
\newunicodechar{Μ}{\ensuremath{\mu}}
\newunicodechar{Τ}{\ensuremath{\tau}}
\newunicodechar{Φ}{\ensuremath{\Phi}}
\newunicodechar{Ψ}{\ensuremath{\Psi}}
\newunicodechar{Ω}{\ensuremath{\Omega}}
\newunicodechar{↦}{\ensuremath{\mapsto}}
\newunicodechar{∈}{\ensuremath{\in}}
\newunicodechar{∉}{\ensuremath{\notin}}
\newunicodechar{∧}{\ensuremath{\wedge}}
\newunicodechar{⇔}{\ensuremath{\Leftrightarrow}}
\newunicodechar{⟺}{\ensuremath{\Longleftrightarrow}}
\newunicodechar{⇒}{\ensuremath{\Rightarrow}}
\newunicodechar{⟹}{\ensuremath{\Longrightarrow}}
\newunicodechar{∘}{\ensuremath{\circ}}
\newunicodechar{∪}{\ensuremath{\cup}}
\newunicodechar{∩}{\ensuremath{\cap}}
\newunicodechar{≡}{\ensuremath{\equiv}}
\newunicodechar{⊂}{\ensuremath{\subset}}
\newunicodechar{⊥}{\ensuremath{\perp}}
\newunicodechar{∃}{\ensuremath{\exists}}
\newunicodechar{∀}{\ensuremath{\forall}}
\newunicodechar{∛}{\ensuremath{\sqrt[3]{\,}}}
\newunicodechar{∜}{\ensuremath{\sqrt[4]{\,}}}
\newunicodechar{⃗}{\ensuremath{{}^{\to}}}
\newunicodechar{ⁿ}{\ifmmode^{n}\else\textsuperscript{\ensuremath{n}}\fi}
\newunicodechar{ˣ}{\ifmmode^{x}\else\textsuperscript{\ensuremath{x}}\fi}
\newunicodechar{ᵇ}{\ifmmode^{b}\else\textsuperscript{\ensuremath{b}}\fi}
\newunicodechar{ʳ}{\ifmmode^{r}\else\textsuperscript{\ensuremath{r}}\fi}
\newunicodechar{ᵘ}{\ifmmode^{u}\else\textsuperscript{\ensuremath{u}}\fi}
\newunicodechar{ʸ}{\ifmmode^{y}\else\textsuperscript{\ensuremath{y}}\fi}
\newunicodechar{ᵃ}{\ifmmode^{a}\else\textsuperscript{\ensuremath{a}}\fi}
\newunicodechar{ᵉ}{\ifmmode^{e}\else\textsuperscript{\ensuremath{e}}\fi}
\newunicodechar{ᵏ}{\ifmmode^{k}\else\textsuperscript{\ensuremath{k}}\fi}
\newunicodechar{ᵖ}{\ifmmode^{p}\else\textsuperscript{\ensuremath{p}}\fi}
\newunicodechar{ᴺ}{\ifmmode^{N}\else\textsuperscript{\ensuremath{N}}\fi}
\newunicodechar{ᶜ}{\ifmmode^{c}\else\textsuperscript{\ensuremath{c}}\fi}
\newunicodechar{ʰ}{\ifmmode^{h}\else\textsuperscript{\ensuremath{h}}\fi}
\newunicodechar{ᵠ}{\ifmmode^{q}\else\textsuperscript{\ensuremath{q}}\fi}
\newunicodechar{ₙ}{\ifmmode_{n}\else\textsubscript{\ensuremath{n}}\fi}
\newunicodechar{ₐ}{\ifmmode_{a}\else\textsubscript{\ensuremath{a}}\fi}
\newunicodechar{ᵢ}{\ifmmode_{i}\else\textsubscript{\ensuremath{i}}\fi}
\newunicodechar{ᵣ}{\ifmmode_{r}\else\textsubscript{\ensuremath{r}}\fi}
\newunicodechar{ₖ}{\ifmmode_{k}\else\textsubscript{\ensuremath{k}}\fi}
\newunicodechar{ᵦ}{\ifmmode_{\beta}\else\textsubscript{\ensuremath{\beta}}\fi}
\newunicodechar{ₓ}{\ifmmode_{x}\else\textsubscript{\ensuremath{x}}\fi}
\newfontfamily\popdisplay{ArchivoBlack-Regular.ttf}[Path=../../../fonts/]
\newfontfamily\poplogo{SpaceGrotesk-Bold.ttf}[Path=../../../fonts/]
\newfontfamily\popsemi{PlusJakartaSans-SemiBold.ttf}[Path=../../../fonts/]
\newfontfamily\popxbold{PlusJakartaSans-ExtraBold.ttf}[Path=../../../fonts/]
\newfontfamily\popxboldls{PlusJakartaSans-ExtraBold.ttf}[Path=../../../fonts/,LetterSpace=3.0]

\setlist[enumerate]{leftmargin=1.7em,itemsep=5pt,topsep=4pt}
\setlist[itemize]{leftmargin=1.7em,itemsep=4pt,topsep=4pt,label={\color{poporange}\textbullet}}
\setlength{\parindent}{0pt}
\setlength{\parskip}{5pt}
\renewcommand{\arraystretch}{1.25}

% pgfplots : style Pop par défaut
\pgfplotsset{
  every axis/.append style={axis line style={popink,line width=1pt},tick style={popink},
    grid style={popline,line width=0.5pt},label style={font=\small},tick label style={font=\footnotesize}},
  popcurve/.style={poporange,line width=1.8pt,no markers,smooth},
}
% Même style disponible dans un \draw TikZ simple (hors pgfplots)
\tikzset{popcurve/.style={poporange,line width=1.8pt}}
\tikzset{popgrid/.style={popline,very thin}}
\colorlet{popcurve}{poporange} % le nom du style est parfois utilisé comme couleur

% ============================================================
% Pill / squiggle
% ============================================================
\newcommand{\poppill}[3]{%
  \tikz[baseline=-0.55ex]\node[draw=popink,line width=1.2pt,rounded corners=2.6pt,%
    fill=#2,inner xsep=6.5pt,inner ysep=3pt]{%
    \popxboldls\fontsize{7}{7}\selectfont\color{#3}\MakeUppercase{#1}};%
}
\newcommand{\popsquiggle}[1]{%
  \tikz[baseline=-0.4ex]{
    \draw[#1,line width=2.6pt,line cap=round]
      plot [smooth] coordinates {(0,0.09) (0.34,0.01) (0.68,0.21) (1.02,0.01) (1.36,0.10)};
  }%
}
% Mot-clé surligné (marqueur orange sous le texte)
\newcommand{\cle}[1]{{\popxbold\color{popdark}#1}}

% ============================================================
% En-tête / pied
% ============================================================
\pagestyle{fancy}
\fancyhf{}
\renewcommand{\headrulewidth}{0pt}
\renewcommand{\footrulewidth}{0pt}
\lhead{\raisebox{-0.15ex}{\poplogo\fontsize{13}{13}\selectfont\color{popink}OnBuch\color{poporange}+}}
\rhead{\poppill{\DOCMATIERE\ \textbullet\ \DOCNIVEAU\ \textbullet\ Leçon \DOCLECON}{white}{popink}}
\lfoot{\popxbold\fontsize{7.6}{7.6}\selectfont\color{popmuted}\MakeUppercase{Cours signé OnBuch+}\ \textbullet\ {\popsemi\DOCTITRE}}
\rfoot{\tikz[baseline=-0.6ex]\node[rounded corners=8.5pt,fill=popink,minimum height=17pt,minimum width=17pt,inner xsep=5pt,inner sep=0]{\popxbold\fontsize{8}{8}\selectfont\color{popcream}\thepage\,/\,\pageref*{LastPage}};}

% ============================================================
% Titres
% ============================================================
\newcommand{\chaptitre}[2]{%
  \begin{tcolorbox}[enhanced,colback=poporange,colframe=popink,boxrule=2.6pt,arc=11pt,
      drop shadow=popink,left=15pt,right=15pt,top=12pt,bottom=11pt,before skip=3pt,after skip=18pt]
    \poppill{#2}{popink}{popcream}\par\vspace{9pt}
    {\raggedright\hyphenpenalty=10000\exhyphenpenalty=10000\popdisplay\fontsize{22}{24}\selectfont\color{popcream}\MakeUppercase{#1}\par}
  \end{tcolorbox}
}
% Section numérotée : pastille orange avec le numéro + titre Archivo
\newcounter{popsec}
\newcommand{\coursec}[1]{%
  \stepcounter{popsec}\setcounter{popsub}{0}%
  \par\vspace{16pt}\noindent
  \begin{minipage}{\linewidth}
  \tikz[baseline=-0.7ex]\node[draw=popink,line width=1.6pt,fill=poporange,rounded corners=4pt,minimum size=22pt,inner sep=0]{\popdisplay\fontsize{12}{12}\selectfont\color{popcream}\thepopsec};%
  \hspace{8pt}{\popdisplay\fontsize{15}{17}\selectfont\color{popink}\MakeUppercase{#1}}\par
  \vspace{4pt}\noindent{\color{poporange}\rule{36pt}{3.6pt}}
  \end{minipage}\par\nopagebreak\vspace{8pt}
}
\newcounter{popsub}
\newcommand{\courssub}[1]{%
  \stepcounter{popsub}%
  \par\vspace{9pt}\noindent{\popxbold\fontsize{11.5}{13}\selectfont\color{popink}{\color{poporange}\thepopsec.\thepopsub}\hspace{6pt}#1}\par\nopagebreak\vspace{4pt}
}

% ============================================================
% Boîtes pédagogiques — même recette : fond blanc, bordure épaisse colorée,
% ombre dure encre, badge plein. Titre optionnel [#1].
% ============================================================
\tcbset{popbase/.style={breakable,enhanced,colback=white,boxrule=2.3pt,arc=9pt,
  shadow={3pt}{-3pt}{0pt}{popink},left=14pt,right=14pt,top=11pt,bottom=11pt,before skip=11pt,after skip=15pt,
  fontupper=\popsemi\fontsize{10.6}{14.5}\selectfont\color{popink},
  fontlower=\popsemi\fontsize{10.6}{14.5}\selectfont\color{popink}}}
% Coche dessinée (le glyphe ✓ n'existe pas dans les polices du document)
\newcommand{\popcheck}[1]{\tikz[baseline=-0.5ex]\draw[#1,line width=1.6pt,line cap=round,line join=round](-0.13,0.0)--(-0.04,-0.09)--(0.13,0.1);}
\providecommand{\coche}{\popcheck{popgreen}}
\providecommand{\resultat}[1]{\par\smallskip\noindent\fcolorbox{popgreen}{popgreenL}{\parbox{\dimexpr\linewidth-2\fboxsep-2\fboxrule\relax}{\textbf{Résultat.} #1}}\par\smallskip}
\newcommand{\popboxhead}[3]{\poppill{#1}{#2}{white}\ifx\relax#3\relax\else\hspace{6pt}{\popxbold\fontsize{10.6}{12}\selectfont\color{popink}#3}\fi\par\vspace{7pt}}

\newtcolorbox{aretenir}[1][]{popbase,colframe=popgreen,before upper={\popboxhead{À retenir}{popgreen}{#1}}}
\newtcolorbox{exemplebox}[1][]{popbase,colframe=poporange,before upper={\popboxhead{Exemple}{poporange}{#1}}}
\newtcolorbox{definition}[1][]{popbase,colframe=popblue,before upper={\popboxhead{Définition}{popblue}{#1}}}
\newtcolorbox{propriete}[1][]{popbase,colframe=poppurple,before upper={\popboxhead{Propriété}{poppurple}{#1}}}
\newtcolorbox{methode}[1][]{popbase,colframe=popgold,before upper={\popboxhead{Méthode}{popgold}{#1}}}
\newtcolorbox{attention}[1][]{popbase,colframe=poppink,before upper={\popboxhead{Attention}{poppink}{#1}}}
\newtcolorbox{experience}[1][]{popbase,colframe=popblue,colback=popblueL,before upper={\popboxhead{Expérience}{popblue}{#1}}}
% « Le savais-tu ? » : carte violette pleine (contexte, culture, Cameroun)
\newtcolorbox{savaistu}[1][]{popbase,colback=poppurple,colframe=popink,
  fontupper=\popsemi\fontsize{10.4}{14.2}\selectfont\color{white},
  before upper={\poppill{Le savais-tu\,?}{white}{poppurple}\ifx\relax#1\relax\else\hspace{6pt}{\popxbold\color{white}#1}\fi\par\vspace{7pt}}}
% Exercice résolu : énoncé en haut, \tcblower puis la solution sur fond crème
\newcounter{popexo}\newcounter{popexr}
\newtcolorbox{exoresolu}[1][]{popbase,colframe=poporange,colbacklower=popcream,
  segmentation style={popink,line width=1.2pt,dashed},
  before upper={\stepcounter{popexr}\popboxhead{Exercice résolu \thepopexr}{poporange}{#1}},
  before lower={\poppill{Solution}{popink}{popcream}\par\vspace{6pt}}}
% Exercice (non résolu en place) — la correction est dans la partie Corrigés
\newtcolorbox{exercice}[1][]{popbase,colframe=popink,
  before upper={\stepcounter{popexo}\popboxhead{Exercice \thepopexo}{popink}{#1}}}
% Corrigé
\newtcolorbox{corrige}[1][]{popbase,colframe=popgreen,colback=popgreenL,
  before upper={\popboxhead{Corrigé}{popgreen}{#1}}}
% Objectifs / prérequis / bilan
\newtcolorbox{objectifs}{popbase,colframe=popink,colback=poporangeL,
  before upper={\popboxhead{Objectifs de la leçon}{popink}{}}}
\newtcolorbox{prerequis}{popbase,colframe=popink,colback=popblueL,
  before upper={\popboxhead{Prérequis}{popblue}{}}}
\newtcolorbox{bilan}{popbase,colframe=popink,colback=white,boxrule=2.8pt,
  before upper={\popboxhead{Fiche bilan}{poporange}{L'essentiel à connaître pour l'examen}}}
% Figure encadrée avec légende
\newtcolorbox{popfigure}[1][]{enhanced,breakable=false,colback=white,colframe=popline,boxrule=1.4pt,arc=8pt,
  left=10pt,right=10pt,top=10pt,bottom=8pt,before skip=10pt,after skip=12pt,halign=center,
  lower separated=false,halign lower=center,
  fontlower=\popsemi\fontsize{9}{11.5}\selectfont\color{popmuted},
  #1}
\newcommand{\legende}[1]{\par\vspace{6pt}{\popsemi\fontsize{9}{11.5}\selectfont\color{popmuted}#1}}

% ============================================================
% Signature OnBuch+
% ============================================================
\newcommand{\poplogoinline}[1]{{\poplogo\fontsize{#1}{#1}\selectfont\color{popink}OnBuch\color{poporange}+}}
\newcommand{\signatureonbuch}{%
  \begin{tcolorbox}[enhanced,colback=popink,colframe=popink,boxrule=2.2pt,arc=10pt,
      drop shadow=poporange,left=14pt,right=14pt,top=10pt,bottom=10pt,before skip=0pt,after skip=0pt]
    \tikz[baseline=-0.7ex]\node[circle,fill=popgreen,draw=popcream,line width=1.3pt,minimum size=22pt,inner sep=0]{\popcheck{white}};%
    \hspace{9pt}\begin{minipage}[c]{0.62\linewidth}
      {\popxbold\fontsize{12}{13}\selectfont\color{popcream}Signé\ \textbullet\ {\poplogo OnBuch\color{poporange}+}}\par\vspace{2pt}
      {\raggedright\popsemi\fontsize{8}{10.5}\selectfont\color{popline}\MakeUppercase{Équipe pédagogique OnBuch+}\\\MakeUppercase{Conforme au programme officiel MINESEC}\par}
    \end{minipage}\hfill
    {\poplogo\fontsize{22}{22}\selectfont\color{popcream}OnBuch\color{poporange}+}
  \end{tcolorbox}%
}

% ============================================================
% Page de garde
% #1 titre · #2 matière · #3 niveau · #4 module · #5 n° leçon · #6 durée ·
% #7 description / objectifs (texte ou itemize)
% ============================================================
\newcommand{\pagegarde}[7]{%
  \thispagestyle{empty}
  \newgeometry{a4paper,margin=1.7cm,top=1.6cm,bottom=1.4cm}%
  \begin{tikzpicture}[remember picture,overlay]
    \fill[poporange!18] ([xshift=-3.2cm,yshift=-3.6cm]current page.north east) circle (4.6cm);
    \fill[poppurple!14] ([xshift=2.4cm,yshift=7.6cm]current page.south west) circle (3.8cm);
  \end{tikzpicture}%
  {\poplogo\fontsize{30}{30}\selectfont\color{popink}OnBuch\color{poporange}+}\hfill
  \raisebox{0.6ex}{\poppill{Cours signé \textbullet\ Premium}{popink}{popcream}}\par
  \vspace{1pt}\popsquiggle{poppurple}\par
  \vspace{8pt}
  \begin{tcolorbox}[enhanced,colback=poporange,colframe=popink,boxrule=2.8pt,arc=13pt,
      drop shadow=popink,left=18pt,right=18pt,top=12pt,bottom=13pt,after skip=10pt]
    \poppill{#2 \textbullet\ #3}{popink}{popcream}\hspace{5pt}\poppill{#4}{white}{popink}\par\vspace{8pt}
    {\popdisplay\fontsize{14}{14}\selectfont\color{popink}LEÇON #5}\par\vspace{4pt}
    {\raggedright\hyphenpenalty=10000\exhyphenpenalty=10000\popdisplay\fontsize{25}{27}\selectfont\color{popcream}\MakeUppercase{#1}\par}
    \vspace{8pt}
    {\popxbold\fontsize{9.5}{9.5}\selectfont\color{popink}\MakeUppercase{Durée conseillée : #6}}
  \end{tcolorbox}
  \begin{tcolorbox}[enhanced,colback=white,colframe=popink,boxrule=2.2pt,arc=10pt,
      drop shadow=popink,left=16pt,right=16pt,top=10pt,bottom=10pt,after skip=8pt]
    \poppill{Dans cette leçon}{poppurple}{white}\par\vspace{5pt}
    \popsemi\fontsize{9.3}{12.6}\selectfont\color{popink}#7
  \end{tcolorbox}
  \noindent\begin{minipage}[t]{0.487\linewidth}
    \begin{tcolorbox}[enhanced,colback=popgreen,colframe=popink,boxrule=2.2pt,arc=10pt,
        drop shadow=popink,left=11pt,right=11pt,top=10pt,bottom=10pt,height=3.1cm,valign=center]
      \tcbox[colback=white,colframe=popink,boxrule=1.3pt,arc=5pt,boxsep=0pt,nobeforeafter,
        left=3pt,right=3pt,top=3pt,bottom=3pt,valign=center]{\qrcode[height=1.9cm]{https://play.google.com/store/apps/details?id=com.luvvix.onbuch&hl=fr}}%
      \hspace{8pt}\begin{minipage}[c]{0.5\linewidth}
        {\popdisplay\fontsize{11}{12.5}\selectfont\color{white}\MakeUppercase{Télécharge OnBuch}}\par\vspace{4pt}
        {\raggedright\popsemi\fontsize{8.4}{11}\selectfont\color{white}Gratuit \textbullet\ Google Play\\Tuteur IA, annales, concours, résultats\par}
      \end{minipage}
    \end{tcolorbox}
  \end{minipage}\hfill
  \begin{minipage}[t]{0.487\linewidth}
    \begin{tcolorbox}[enhanced,colback=poppurple,colframe=popink,boxrule=2.2pt,arc=10pt,
        drop shadow=popink,left=11pt,right=11pt,top=10pt,bottom=10pt,height=3.1cm,valign=center]
      \tikz[baseline=-0.7ex]\node[circle,fill=white,draw=popink,line width=1.3pt,minimum size=1.6cm,inner sep=0]{\popdisplay\fontsize{24}{24}\selectfont\color{poppurple}+};%
      \hspace{8pt}\begin{minipage}[c]{0.6\linewidth}
        {\popdisplay\fontsize{11}{12.5}\selectfont\color{white}\MakeUppercase{Passe à OnBuch+}}\par\vspace{4pt}
        {\raggedright\popsemi\fontsize{8.4}{11}\selectfont\color{white}Tous les cours signés, Tuteur IA illimité, fiches PDF — onglet OnBuch+ de l'app\par}
      \end{minipage}
    \end{tcolorbox}
  \end{minipage}
  \vspace{8pt}
  \popcontenu
  \vfill
  \signatureonbuch
  \restoregeometry
  \newpage
}

% Rangée « Ce cours contient » (page de garde)
\newcommand{\poptile}[3]{% #1 couleur #2 gros texte #3 légende
  \begin{tcolorbox}[enhanced,colback=#1,colframe=popink,boxrule=2pt,arc=9pt,shadow={2.5pt}{-2.5pt}{0pt}{popink},
      left=6pt,right=6pt,top=9pt,bottom=9pt,halign=center,valign=center,height=1.75cm,width=\linewidth,nobeforeafter]
    {\popdisplay\fontsize{12.5}{13}\selectfont\color{white}\MakeUppercase{#2}}\par\vspace{3pt}
    {\centering\popsemi\fontsize{7.8}{9.5}\selectfont\color{white}#3\par}
  \end{tcolorbox}}
\newcommand{\popcontenu}{%
  \par\noindent{\popxboldls\fontsize{7.5}{7.5}\selectfont\color{popmuted}CE COURS SIGNÉ CONTIENT}\par\vspace{7pt}
  \noindent\begin{minipage}[t]{0.235\linewidth}\poptile{popblue}{Cours}{complet, illustré, pas à pas}\end{minipage}\hfill
  \begin{minipage}[t]{0.235\linewidth}\poptile{poporange}{Méthodes}{et exercices résolus}\end{minipage}\hfill
  \begin{minipage}[t]{0.235\linewidth}\poptile{poppink}{Exercices}{type examen + corrigés}\end{minipage}\hfill
  \begin{minipage}[t]{0.235\linewidth}\poptile{popgreen}{Fiche bilan}{l'essentiel à retenir}\end{minipage}\par}

% Sommaire
\newtcolorbox{sommaire}{enhanced,colback=white,colframe=popink,boxrule=2.2pt,arc=10pt,
  drop shadow=popink,left=18pt,right=18pt,top=13pt,bottom=13pt,before skip=6pt,after skip=20pt,
  before upper={\poppill{Sommaire}{poppurple}{white}\par\vspace{9pt}\popsemi\fontsize{10.6}{14.5}\selectfont\color{popink}}}

% Signature de fin de cours — poussée tout en bas de la dernière page
\newcommand{\findecours}{%
  \par\vfill\nopagebreak
  \begin{center}\popsquiggle{poporange}\end{center}\vspace{4pt}
  \signatureonbuch
}
\AtBeginDocument{\def\llbracket{\mathopen{[\mkern-3.5mu[}}\def\rrbracket{\mathclose{]\mkern-3.5mu]}}} % intervalles d'entiers (glyphe ⟦ absent de la police)
% Glyphes absents de Fira Math : redéfinis à partir de glyphes disponibles
\AtBeginDocument{%
  \def\setminus{\mathbin{\backslash}}\let\smallsetminus\setminus
  \def\vdots{\mathord{\vbox{\baselineskip3.2pt\lineskiplimit0pt\kern4pt\hbox{.}\hbox{.}\hbox{.}}}}%
  \def\ddots{\mathinner{\mkern1mu\raise6.5pt\hbox{.}\mkern2mu\raise3.6pt\hbox{.}\mkern2mu\raise0.7pt\hbox{.}\mkern1mu}}%
  \def\triangle{\mathord{\scalebox{0.9}{$\symup{\Delta}$}}}%
  \def\nabla{\mathord{\raisebox{\depth}{\scalebox{1}[-1]{$\symup{\Delta}$}}}}%
  \def\checkmark{\popcheck{popdark}}%
  \def\star{\mathbin{\ast}}\def\bigstar{\mathord{\ast}}%
  \def\diamond{\mathbin{\scalebox{0.6}{$\Diamond$}}}%
  \def\bigcup{\mathop{\mathchoice{\raisebox{-0.3em}{\scalebox{1.7}{$\cup$}}}{\scalebox{1.25}{$\cup$}}{\cup}{\cup}}\displaylimits}%
  \def\bigcap{\mathop{\mathchoice{\raisebox{-0.3em}{\scalebox{1.7}{$\cap$}}}{\scalebox{1.25}{$\cap$}}{\cap}{\cap}}\displaylimits}%
  \def\lesseqqgtr{\mathrel{\lessgtr}}\def\bigoplus{\mathop{\oplus}\displaylimits}%
}
\providecommand{\ssi}{si et seulement si} % abréviation fréquente
\AtBeginDocument{\providecommand{\wideparen}{\overparen}} % arc de cercle

%% FIGFIX-BEGIN v5 — correctifs globaux des figures (docs/onbuch-generation/tools/figfix)
\usepackage{adjustbox}\usepackage{etoolbox}\tcbuselibrary{raster}
% toute figure TikZ/pgfplots plus large que la ligne est réduite à la largeur disponible
\BeforeBeginEnvironment{tikzpicture}{\begin{adjustbox}{max width=\linewidth}}
\AfterEndEnvironment{tikzpicture}{\end{adjustbox}}
% légendes pgfplots lisibles par-dessus les courbes
\pgfplotsset{every axis/.append style={legend style={fill=white,fill opacity=0.92,text opacity=1,draw=black!15,inner sep=3pt}}}
% étiquettes d'arêtes (syntaxe edge["..."]) avec fond blanc
\tikzset{every edge quotes/.append style={fill=white,inner sep=1.2pt}}
%% FIGFIX-END
\begin{document}

\pagegarde{La lunette astronomique}{Physique}{1ère TI}{Module 3 : Optique}{N10}{5 h}{Cette leçon présente la lunette astronomique : sa constitution (objectif et oculaire), son fonctionnement en système afocal et la marche des rayons lumineux. Elle définit et exploite le grossissement et la puissance, et guide l'élève dans la mise au point et le choix pertinent d'un instrument d'observation.
\vspace{4pt}
\begin{multicols}{2}\small
\begin{itemize}
\item À la fin de cette leçon, je sais décrire la constitution d'une lunette astronomique (objectif, oculaire) et en tracer le schéma de principe
\item À la fin de cette leçon, je sais expliquer le fonctionnement d'une lunette astronomique réglée pour une vision à l'infini (système afocal)
\item À la fin de cette leçon, je sais tracer la marche d'un faisceau lumineux issu d'un objet à l'infini à travers la lunette
\item À la fin de cette leçon, je sais définir le grossissement et la puissance d'une lunette astronomique
\item À la fin de cette leçon, je sais démontrer et utiliser la relation G = f'₁/f'₂ pour calculer le grossissement
\item À la fin de cette leçon, je sais effectuer la mise au point d'une lunette astronomique pour un œil normal
\end{itemize}\end{multicols}}

\begin{sommaire}
\begin{multicols}{2}
\begin{enumerate}
\item Constitution et description de la lunette astronomique
\item Fonctionnement : la lunette afocale
\item Tracé de la marche d'un faisceau lumineux
\item Grossissement et puissance
\item Mise au point et utilisation pratique
\item Bilan et choix raisonné d'un instrument d'optique
\item Activité d'intégration
\item Méthodes et exercices résolus
\item Exercices
\item Corrigés des exercices
\item Fiche bilan
\end{enumerate}
\end{multicols}
\end{sommaire}

\begin{objectifs}
\begin{itemize}
\item À la fin de cette leçon, je sais décrire la constitution d'une lunette astronomique (objectif, oculaire) et en tracer le schéma de principe
\item À la fin de cette leçon, je sais expliquer le fonctionnement d'une lunette astronomique réglée pour une vision à l'infini (système afocal)
\item À la fin de cette leçon, je sais tracer la marche d'un faisceau lumineux issu d'un objet à l'infini à travers la lunette
\item À la fin de cette leçon, je sais définir le grossissement et la puissance d'une lunette astronomique
\item À la fin de cette leçon, je sais démontrer et utiliser la relation G = f'₁/f'₂ pour calculer le grossissement
\item À la fin de cette leçon, je sais effectuer la mise au point d'une lunette astronomique pour un œil normal
\item À la fin de cette leçon, je sais choisir de manière pertinente un instrument d'optique (lunette, jumelles, microscope) et justifier mon choix
\end{itemize}
\end{objectifs}

\begin{prerequis}
\begin{itemize}
\item Lentilles minces convergentes : foyers, distance focale, vergence (leçons précédentes du Module 3)
\item Formules de conjugaison et de grandissement des lentilles minces
\item Construction géométrique des images données par une lentille convergente
\item L'œil normal : accommodation, vision d'un objet à l'infini sans fatigue
\item Notion d'angle sous lequel on voit un objet (diamètre apparent)
\end{itemize}
\end{prerequis}

\newpage

\coursec{Constitution et description de la lunette astronomique}

Lors des veillées nocturnes dans les villages du Nord-Cameroun, les anciens montrent aux enfants les cratères de la Lune et les étoiles filantes, mais l'œil nu ne permet de distinguer que peu de détails. Un élève de Waza qui observe la Lune avec une paire de jumelles découvre soudain ses montagnes et ses mers : comment un simple assemblage de deux lentilles peut-il ainsi \og rapprocher \fg{} les astres ? Pourquoi les astronomes parlent-ils de \cle{grossissement} et de \cle{puissance} ? Cette leçon explique le fonctionnement de la lunette astronomique, instrument qui a révolutionné notre vision du ciel depuis Galilée.

\courssub{Le problème posé : observer un objet très éloigné}

Regardons la Lune un soir de pleine lune. Son diamètre réel est d'environ $3\,474\ \text{km}$ et elle se situe à environ $3{,}84\times 10^{5}\ \text{km}$ de la Terre. Sous quel angle la voyons-nous ?

\begin{definition}[Diamètre apparent]
Le \cle{diamètre apparent} d'un objet est l'angle $\theta$ sous lequel l'objet est vu depuis l'œil de l'observateur. Pour un objet de taille $AB$ situé à la distance $d$ de l'œil, avec $\theta$ petit et exprimé en radians :
\[
\theta \simeq \tan\theta = \frac{AB}{d}.
\]
\end{definition}

\begin{exemplebox}[Diamètre apparent de la Lune]
Pour la Lune : $\theta = \dfrac{3\,474}{3{,}84\times 10^{5}} \simeq 9{,}0\times 10^{-3}\ \text{rad}$, soit environ $0{,}5^{\circ}$. C'est l'angle sous lequel on voit une pièce de 10 francs tenue à bout de bras ! L'œil humain ne distingue deux points que s'ils sont vus sous un angle supérieur à environ $3\times 10^{-4}\ \text{rad}$ (un minute d'arc) : les détails lunaires plus petits que $100\ \text{km}$ nous échappent totalement.
\end{exemplebox}

Deux difficultés se posent donc pour observer un astre :
\begin{itemize}
\item son \cle{diamètre apparent} est extrêmement faible : l'image se forme sur une toute petite zone de la rétine ;
\item la quantité de lumière reçue par la pupille (quelques millimètres de diamètre) est très faible : l'astre paraît peu lumineux.
\end{itemize}

L'idée géniale est alors la suivante : utiliser une première lentille de \emph{grand diamètre} pour \textbf{collecter beaucoup de lumière} et former une image réelle de l'astre, puis examiner cette image à la \textbf{loupe} avec une seconde lentille. C'est exactement le principe de la lunette astronomique.

\courssub{Constitution de la lunette astronomique}

\begin{definition}[Lunette astronomique]
La \cle{lunette astronomique} est un instrument d'optique formé de \textbf{deux lentilles convergentes coaxiales} (même axe optique principal) :
\begin{itemize}
\item l'\cle{objectif} $L_1$, placé du côté de l'objet observé ;
\item l'\cle{oculaire} $L_2$, placé du côté de l'œil de l'observateur.
\end{itemize}
Elle est destinée à l'observation d'objets très éloignés (considérés \og à l'infini \fg{}), qu'elle permet de voir sous un angle plus grand qu'à l'œil nu.
\end{definition}

\textbf{L'objectif $L_1$ : la grande lentille qui collecte la lumière.}

L'objectif est une lentille convergente de \textbf{grande distance focale} $f'_1$ (de l'ordre de $50\ \text{cm}$ à plusieurs mètres) et de \textbf{grand diamètre} (plusieurs centimètres). Son rôle est double :
\begin{itemize}
\item \textbf{collecter} un maximum de lumière venant de l'astre : plus son diamètre est grand, plus l'image est lumineuse ;
\item \textbf{former une image réelle} $A_1B_1$ de l'astre. Comme l'astre est à l'infini, cette image se forme dans le \emph{plan focal image} de l'objectif ; elle est réelle, renversée et très petite.
\end{itemize}

\textbf{L'oculaire $L_2$ : la petite lentille qui joue le rôle de loupe.}

L'oculaire est une lentille convergente de \textbf{courte distance focale} $f'_2$ (de l'ordre du centimètre) et de petit diamètre. Il est placé de manière à examiner l'image intermédiaire $A_1B_1$ donnée par l'objectif : il fonctionne exactement comme une \emph{loupe} et en donne une \textbf{image définitive virtuelle et agrandie}, que l'œil observe sans effort.

\begin{attention}[Ne pas inverser objectif et oculaire]
Erreur très fréquente au Probatoire : confondre les deux lentilles. Retiens bien :
\begin{itemize}
\item l'\textbf{objectif} est toujours la lentille de \textbf{plus grande distance focale} ($f'_1$ grand) et de plus grand diamètre ; il est placé \textbf{du côté de l'objet} ;
\item l'\textbf{oculaire} est la lentille de \textbf{plus courte distance focale} ($f'_2$ petit) ; il est placé \textbf{du côté de l'œil}.
\end{itemize}
Une lunette montée à l'envers (œil derrière l'objectif) ne grossit pas : elle \emph{réduit} !
\end{attention}

\courssub{Schéma de principe et positions relatives des lentilles}

Les deux lentilles sont montées aux extrémités d'un tube, de sorte que leurs axes optiques principaux coïncident : on dit qu'elles sont \cle{coaxiales}. La distance $O_1O_2$ entre les centres optiques est réglable (tube coulissant) pour permettre la \emph{mise au point}, que nous étudierons en section 5.

Dans le réglage le plus important, dit réglage \cle{afocal} (étudié en détail en section 2), le \textbf{foyer image $F'_1$ de l'objectif coïncide avec le foyer objet $F_2$ de l'oculaire}. La longueur de la lunette est alors :
\[
O_1O_2 = f'_1 + f'_2.
\]

\begin{popfigure}
\begin{center}
\begin{tikzpicture}[scale=1.0, every node/.style={font=\small}]
% Axe optique
\draw[popmuted, dashed, thick] (-0.8,0) -- (11.5,0) node[below left=2pt and 0pt] {};
\node[popmuted, below] at (10.8,0) {axe optique};
% Objectif L1 (grande lentille)
\draw[popblue, very thick] (1.5,-2.2) -- (1.5,2.2);
\draw[popblue, very thick, -{Stealth}] (1.5,-2.2) -- (1.35,-1.95);
\draw[popblue, very thick, -{Stealth}] (1.5,2.2) -- (1.35,1.95);
\node[popblue, above] at (1.5,2.25) {$L_1$ objectif};
\node[popblue, below] at (1.5,-2.25) {$f'_1$ grand};
% Oculaire L2 (petite lentille)
\draw[poporange, very thick] (9.5,-1.0) -- (9.5,1.0);
\draw[poporange, very thick, -{Stealth}] (9.5,-1.0) -- (9.35,-0.8);
\draw[poporange, very thick, -{Stealth}] (9.5,1.0) -- (9.35,0.8);
\node[poporange, above] at (9.5,1.05) {$L_2$ oculaire};
\node[poporange, below] at (9.5,-1.05) {$f'_2$ petit};
% Centres optiques
\fill[popdark] (1.5,0) circle (1.5pt) node[below=2pt] {$O_1$};
\fill[popdark] (9.5,0) circle (1.5pt) node[below=2pt] {$O_2$};
% Foyer F'1 = F2 confondus
\fill[popgreen] (7.5,0) circle (2pt);
\node[popgreen, below=2pt] at (7.5,0) {$F'_1 \equiv F_2$};
% Foyer F1
\fill[popmuted] (-4.5,0) circle (0pt);
\fill[popmuted] (7.5,0) circle (0pt);
% Distance focale f'1
\draw[popblue, <->] (1.5,1.6) -- (7.5,1.6) node[midway, above] {$f'_1$};
% Distance focale f'2
\draw[poporange, <->] (7.5,0.7) -- (9.5,0.7) node[midway, above] {$f'_2$};
% Astre à l'infini
\node[popdark, left] at (-0.8,1.5) {astre à l'infini};
\draw[popgold, thick, ->] (-0.6,1.4) -- (0.9,0.6);
\draw[popgold, thick, ->] (-0.6,1.1) -- (0.9,0.35);
% Oeil
\node[popdark, right] at (11.4,0.3) {\oe il};
\end{tikzpicture}
\end{center}
\legende{Schéma de principe de la lunette astronomique réglée en afocal : l'objectif $L_1$ (grande distance focale $f'_1$, grand diamètre) et l'oculaire $L_2$ (courte distance focale $f'_2$) sont coaxiaux ; le foyer image $F'_1$ de l'objectif coïncide avec le foyer objet $F_2$ de l'oculaire.}
\end{popfigure}

Remarque sur le schéma : l'objectif est dessiné \emph{plus grand} que l'oculaire, ce qui traduit physiquement son grand diamètre (sa capacité à collecter la lumière). La distance $f'_1$ est très supérieure à $f'_2$ : c'est la signature géométrique de toute lunette astronomique.

\courssub{Rôle de chaque élément : la chaîne des images}

Suivons le destin de la lumière venue de l'astre, étape par étape. C'est la clé de tout le fonctionnement de l'instrument.

\textbf{Étape 1 : l'objectif forme une image réelle réduite.}

L'astre $AB$ est à l'infini ; on le voit sous le diamètre apparent $\theta$. Tous les rayons issus du point $B$ (à l'infini) arrivent sur l'objectif en \emph{faisceau parallèle} incliné de l'angle $\theta$ par rapport à l'axe. L'objectif les fait converger en un point $B_1$ de son plan focal image : l'image $A_1B_1$ est \textbf{réelle}, \textbf{renversée} et \textbf{réduite}. Sa taille se calcule immédiatement :
\[
A_1B_1 = f'_1 \times \tan\theta \simeq f'_1\,\theta \qquad (\theta \text{ en radians}).
\]
Vérification dimensionnelle : $f'_1$ est une longueur, $\theta$ est sans dimension, donc $f'_1\theta$ est bien une longueur.

\begin{exemplebox}[Taille de l'image de la Lune donnée par l'objectif]
Avec $f'_1 = 1{,}00\ \text{m}$ et $\theta = 9{,}0\times 10^{-3}\ \text{rad}$ :
\[
A_1B_1 = 1{,}00 \times 9{,}0\times 10^{-3} = 9{,}0\times 10^{-3}\ \text{m} = 9{,}0\ \text{mm}.
\]
Toute la Lune tient dans une image réelle de $9\ \text{mm}$, située dans le plan focal de l'objectif. C'est cette image que l'oculaire va examiner.
\end{exemplebox}

\textbf{Étape 2 : l'oculaire joue le rôle de loupe.}

L'image intermédiaire $A_1B_1$ sert d'\emph{objet} pour l'oculaire. Celui-ci est placé de telle sorte que $A_1B_1$ soit dans son plan focal objet (réglage afocal) ou légèrement à l'intérieur : il en donne alors une image définitive $A'B'$ \textbf{virtuelle}, \textbf{agrandie}, rejetée à l'infini (ou au punctum remotum de l'œil), que l'observateur regarde \emph{sans accommoder}, donc sans fatigue.

\begin{propriete}[Nature de l'image définitive]
L'image définitive $A'B'$ donnée par une lunette astronomique est :
\begin{itemize}
\item \textbf{virtuelle} (elle ne peut pas être recueillie sur un écran) ;
\item \textbf{renversée} par rapport à l'objet (deux renversements successifs ? Non : l'objectif renverse, la loupe ne renverse pas, donc au total l'image est renversée) ;
\item vue sous un angle $\theta'$ \textbf{plus grand} que le diamètre apparent $\theta$ de l'objet : c'est tout l'intérêt de l'instrument.
\end{itemize}
Le renversement de l'image est sans gravité en astronomie : qu'il y ait des cratères \og en haut \fg{} ou \og en bas \fg{} de la Lune importe peu ! En revanche, pour l'observation terrestre (jumelles, longues-vues), on redresse l'image à l'aide de prismes ou d'une lentille supplémentaire.
\end{propriete}

\begin{popfigure}
\begin{center}
\begin{tikzpicture}[scale=1.0, every node/.style={font=\small}]
% Grande lentille (objectif)
\draw[popblue, very thick] (0,-2.6) -- (0,2.6);
\draw[popblue, very thick, -{Stealth}] (0,-2.6) -- (-0.18,-2.3);
\draw[popblue, very thick, -{Stealth}] (0,2.6) -- (-0.18,2.3);
\node[popblue, below] at (0,-2.7) {objectif : grand diamètre, $f'$ grande};
% Petite lentille (oculaire)
\draw[poporange, very thick] (4.5,-1.1) -- (4.5,1.1);
\draw[poporange, very thick, -{Stealth}] (4.5,-1.1) -- (4.32,-0.85);
\draw[poporange, very thick, -{Stealth}] (4.5,1.1) -- (4.32,0.85);
\node[poporange, below] at (4.5,-1.2) {oculaire : petit diamètre, $f'$ courte};
% flèche lumière
\draw[popgold, very thick, ->] (-2.2,0) -- (-0.6,0) node[midway, above, popdark] {lumière de l'astre};
% flèche vers oeil
\draw[popgold, very thick, ->] (5.1,0) -- (6.6,0) node[right, popdark] {vers l'\oe il};
\end{tikzpicture}
\end{center}
\legende{Comparaison des deux lentilles : l'objectif, de grand diamètre, collecte la lumière de l'astre ; l'oculaire, de petite taille et de courte focale, sert de loupe pour examiner l'image intermédiaire.}
\end{popfigure}

\courssub{Exemple chiffré : identifier les éléments d'une lunette scolaire}

\begin{exoresolu}[La lunette du club d'astronomie de Garoua]
Le club de sciences d'un lycée de Garoua possède une lunette dont les deux lentilles ont pour vergences $C_a = +1{,}0\ \delta$ et $C_b = +50\ \delta$.
\begin{enumerate}
\item Calculer les distances focales des deux lentilles.
\item Identifier l'objectif et l'oculaire. Justifier.
\item Quelle est la longueur minimale du tube en réglage afocal ?
\end{enumerate}
\tcblower
\textbf{Solution.}
\begin{enumerate}
\item La distance focale image est l'inverse de la vergence : $f' = \dfrac{1}{C}$.
\begin{align*}
f'_a &= \frac{1}{C_a} = \frac{1}{1{,}0} = 1{,}0\ \text{m} = 100\ \text{cm},\\
f'_b &= \frac{1}{C_b} = \frac{1}{50} = 0{,}020\ \text{m} = 2{,}0\ \text{cm}.
\end{align*}
\item L'\textbf{objectif} est la lentille de \emph{plus grande distance focale} : c'est la lentille $a$ ($f'_1 = 100\ \text{cm}$), placée du côté de l'astre. L'\textbf{oculaire} est la lentille $b$ de courte focale ($f'_2 = 2{,}0\ \text{cm}$), placée du côté de l'œil.
\item En réglage afocal, $F'_1 \equiv F_2$, donc :
\[
O_1O_2 = f'_1 + f'_2 = 100 + 2{,}0 = 102\ \text{cm}.
\]
Le tube mesure environ $1{,}02\ \text{m}$ : c'est bien l'allure d'une lunette astronomique scolaire, longue et fine.
\end{enumerate}
\end{exoresolu}

\begin{attention}[Pièges de rédaction au Probatoire]
\begin{itemize}
\item Ne jamais écrire que l'objectif \og grossit \fg{} : l'objectif forme une image réelle \emph{réduite} de l'astre ; c'est l'ensemble objectif + oculaire qui produit un grossissement angulaire.
\item Ne pas confondre distance focale et diamètre : l'objectif possède les deux \og records \fg{} (grande focale \emph{et} grand diamètre), mais ce sont deux grandeurs différentes.
\item L'image définitive est \textbf{virtuelle} : elle ne peut pas être projetée sur un écran. Seule l'image intermédiaire $A_1B_1$ est réelle.
\item Dans les calculs, convertir toutes les distances focales dans la même unité (le mètre en SI) avant toute application numérique.
\end{itemize}
\end{attention}

\begin{savaistu}[1609 : Galilée retourne le ciel]
En 1609, le savant italien \textbf{Galilée} entend parler d'un instrument hollandais qui \og rapproche les objets éloignés \fg{}. Il en construit une version améliorée et, le premier, la pointe vers le ciel : il découvre les \textbf{cratères et montagnes de la Lune}, les phases de Vénus, les taches solaires et quatre \textbf{satellites de Jupiter}. Ces observations ruinent l'idée d'un ciel parfait et immuable et apportent un soutien décisif au système héliocentrique de Copernic. Quatre siècles plus tard, les plus grands télescopes du monde — et même les modestes lunettes des clubs de sciences de Yaoundé ou de Garoua — fonctionnent toujours sur le même principe : un objectif qui collecte la lumière, un oculaire qui examine l'image.
\end{savaistu}

\begin{aretenir}[L'essentiel de cette section]
\begin{itemize}
\item La lunette astronomique comporte \textbf{deux lentilles convergentes coaxiales} : l'\textbf{objectif} (grande focale $f'_1$, grand diamètre, côté objet) et l'\textbf{oculaire} (courte focale $f'_2$, côté œil).
\item L'objectif forme de l'astre une image \textbf{réelle, renversée et réduite} $A_1B_1$ dans son plan focal image, de taille $A_1B_1 = f'_1\,\theta$.
\item L'oculaire fonctionne en \textbf{loupe} : il donne de $A_1B_1$ une image définitive \textbf{virtuelle, agrandie et renversée}, observée sans fatigue.
\item En réglage afocal : $F'_1 \equiv F_2$ et $O_1O_2 = f'_1 + f'_2$.
\end{itemize}
\end{aretenir}

\coursec{Fonctionnement : la lunette afocale}

Dans la section précédente, nous avons décrit la constitution d'une lunette astronomique : un \cle{objectif} de grande distance focale $f'_1$, un \cle{oculaire} de courte distance focale $f'_2$, montés sur le même axe optique à l'intérieur d'un tube. Reste à comprendre \emph{comment} ces deux lentilles coopèrent pour permettre à un œil humain d'observer confortablement un astre situé, pour ainsi dire, infiniment loin. C'est l'objet de cette section : nous allons suivre la lumière pas à pas, de l'étoile jusqu'à la rétine, et découvrir la condition géométrique fondamentale qui fait d'une lunette un \cle{système afocal}.

\courssub{Un objet à l'infini : le faisceau parallèle incident}

\textbf{Une évidence... qu'il faut justifier.}\par

Regarde la Lune depuis la cour du lycée : elle se trouve à environ $\SI{384000}{km}$. Une étoile comme Sirius est à $8,6$ années-lumière, soit environ $8\times 10^{13}\ \si{km}$. Même pour la planète Mars au plus près de la Terre, la distance dépasse $\SI{5e7}{km}$. Or le diamètre de la monture d'un objectif de lunette scolaire est de quelques centimètres, disons $D = \SI{5}{cm}$.

Considérons un point $B$ d'un astre, hors de l'axe optique. Ce point émet de la lumière dans toutes les directions ; seule une infime partie de ces rayons pénètre dans la lunette, ceux qui traversent la monture de diamètre $D$. L'angle maximal entre deux de ces rayons est de l'ordre de
\[
\frac{D}{d} = \frac{\num{5e-2}}{\num{3,84e8}} \approx \num{1,3e-10}\ \si{rad}
\]
pour la Lune : c'est absolument négligeable. Les rayons issus de $B$ qui entrent dans la lunette sont donc, à toute fin pratique, \emph{parallèles entre eux}.

\begin{aretenir}[Objet à l'infini]
Un astre étant situé à une distance immense devant les dimensions de l'instrument, tous les rayons issus d'un même point $B$ de l'astre qui pénètrent dans la lunette forment un \cle{faisceau parallèle}. Comme $B$ n'est pas sur l'axe optique, ce faisceau est \textbf{incliné d'un angle $\theta$} par rapport à l'axe optique. Cet angle $\theta$ est précisément le \emph{diamètre apparent} sous lequel on voit l'astre (ou la portion d'astre) à l'œil nu.
\end{aretenir}

C'est ce point qui donne tout son intérêt à la lunette : l'angle $\theta$ est minuscule (le diamètre apparent de la Lune vaut environ $0,5^\circ$, celui de Jupiter au mieux $50''$, soit $0,014^\circ$). L'œil nu ne distingue pas de détails sous un si petit angle. Toute la stratégie de la lunette consiste à \emph{amplifier cet angle}, comme nous le verrons dans la section 4.

\courssub{Première étape : l'objectif forme une image réelle dans son plan focal image}

Que devient ce faisceau parallèle incliné en traversant l'objectif, lentille convergente de distance focale $f'_1$ ? Appliquons une propriété fondamentale des lentilles minces : un faisceau de rayons parallèles entre eux converge, après traversée d'une lentille convergente, en un point de son \cle{plan focal image}. Si le faisceau est parallèle à l'axe, ce point est le foyer image $F'_1$ lui-même ; s'il est incliné d'un angle $\theta$, le point de convergence $B_1$ est situé dans le plan focal image, \emph{à côté} de $F'_1$.

Autrement dit, l'objectif joue le rôle d'une chambre noire à lentille : il forme de l'astre une \cle{image intermédiaire} $A_1B_1$ qui est :
\begin{itemize}
\item \textbf{réelle} : les rayons convergent effectivement en ses points, on pourrait la recueillir sur un écran translucide placé dans le plan focal image ;
\item située \textbf{dans le plan focal image de l'objectif}, c'est-à-dire à la distance $f'_1$ derrière $O_1$ ;
\item \textbf{renversée} par rapport à l'objet : si $B$ est « au-dessus » de l'axe (faisceau venant d'en haut), $B_1$ est en dessous, et réciproquement ;
\item de taille $A_1B_1 = f'_1 \tan\theta \approx f'_1\,\theta$ (pour les petits angles, en radians).
\end{itemize}

Vérifions l'ordre de grandeur de cette dernière relation, qui sera précieuse pour le grossissement : pour la Lune ($\theta \approx \num{4,4e-3}\ \si{rad}$ en prenant le rayon apparent) et un objectif de $f'_1 = \SI{80}{cm}$, l'image intermédiaire mesure $A_1B_1 \approx 0,80 \times \num{4,4e-3} \approx \SI{3,5}{mm}$. Toute la richesse de l'observation est concentrée dans ces quelques millimètres : c'est cette petite image que l'oculaire va « examiner » comme une loupe.

\begin{attention}[L'image intermédiaire n'est pas l'image définitive]
Erreur fréquente au Probatoire : affirmer que « l'image se forme sur la rétine au foyer de l'objectif ». Non ! L'image $A_1B_1$ est une étape \emph{intermédiaire}, invisible pour l'observateur tant que l'oculaire n'est pas en place. C'est l'oculaire qui la transforme, et c'est \textbf{l'œil} qui forme l'image définitive sur la rétine à partir du faisceau sortant de l'oculaire. La lunette seule ne produit jamais d'image sur la rétine.
\end{attention}

\courssub{Deuxième étape : l'oculaire et la vision sans fatigue}

\textbf{Le problème de l'observateur.}\par

L'image intermédiaire $A_1B_1$ flotte dans l'air, à l'intérieur du tube. L'observateur la regarde \emph{à travers} l'oculaire, qui joue exactement le rôle d'une \textbf{loupe} : une lentille convergente de courte focale $f'_2$ qui donne d'un petit objet une image agrandie et confortable à observer.

Rappelons alors une contrainte physiologique essentielle. Un \cle{œil normal} (emmetrope) voit nettement \emph{sans accommoder} — c'est-à-dire sans fatigue, cristallin au repos — uniquement les objets situés \textbf{à l'infini}. Observer longtemps un objet proche impose au cristallin un effort permanent de convergence : c'est la fatigue visuelle que tu ressens après une heure de lecture de trop près. Pour une observation astronomique prolongée et confortable, il faut donc que l'image définitive $A'B'$ donnée par l'oculaire soit \textbf{rejetée à l'infini}.

\begin{propriete}[Condition de vision sans fatigue]
Pour qu'un œil normal observe \emph{sans accommoder} (vision sans fatigue), l'image définitive $A'B'$ donnée par la lunette doit être située \textbf{à l'infini} : les rayons qui sortent de l'oculaire doivent former un faisceau parallèle.
\end{propriete}

Or, d'après les propriétés des lentilles minces, une lentille convergente donne d'un objet une image à l'infini si et seulement si cet objet est placé \textbf{dans son plan focal objet}. Appliqué à l'oculaire : l'image intermédiaire $A_1B_1$ doit se trouver dans le plan focal objet de l'oculaire.

Mais nous avons vu que $A_1B_1$ se forme dans le plan focal \emph{image} de l'objectif. Les deux contraintes se combinent alors magnifiquement :

\begin{definition}[Lunette afocale]
Une lunette astronomique est dite \cle{afocale} lorsque le \textbf{foyer image de l'objectif} et le \textbf{foyer objet de l'oculaire} sont confondus :
\[
F'_1 \equiv F_2
\]
L'image intermédiaire $A_1B_1$, formée dans le plan focal image de l'objectif, se trouve alors automatiquement dans le plan focal objet de l'oculaire : l'oculaire la renvoie à l'infini, et l'œil normal l'observe sans accommoder.
\end{definition}

\courssub{Le système afocal : un faisceau parallèle en devient un autre}

Résumons la chaîne complète des transformations, qui justifie le nom « afocal » :
\begin{enumerate}
\item L'astre, à l'infini, envoie un \textbf{faisceau parallèle} incliné de $\theta$ sur l'axe ;
\item l'objectif le fait converger au point $B_1$ de son plan focal image ;
\item de $B_1$, les rayons repartent en divergeant et traversent l'oculaire ; comme $B_1$ est dans le plan focal \emph{objet} de l'oculaire, ils ressortent en un \textbf{faisceau parallèle}, incliné d'un nouvel angle $\theta'$ sur l'axe ;
\item l'œil, placé derrière l'oculaire, reçoit ce faisceau parallèle et le focalise sur la rétine sans effort.
\end{enumerate}

\begin{propriete}[Propriété du système afocal — admise, justifiée par la construction des rayons]
Un système afocal transforme tout \textbf{faisceau parallèle incident en un faisceau parallèle émergent} : l'image d'un objet à l'infini est elle-même à l'infini. Un tel système n'a \textbf{ni foyer, ni distance focale} propre (d'où le nom « a-focal », littéralement « sans foyer ») : on ne peut pas lui appliquer les formules de conjugaison des lentilles globalement ; on raisonne lentille par lentille.
\end{propriete}

La construction des rayons ci-dessous rend cette propriété évidente : le parallélisme est « conservé » par la double traversée, seule l'\emph{inclinaison} change ($\theta' \neq \theta$) — et c'est précisément ce changement d'inclinaison qui produira le grossissement étudié en section 4.

\begin{popfigure}
\begin{center}
\begin{tikzpicture}[scale=1.0,>=Stealth]
% Axe optique
\draw[popmuted,->] (-1.2,0) -- (9.8,0) node[below left,font=\small] {axe optique};
% Objectif
\draw[popblue,very thick] (2,-1.9) -- (2,1.9);
\draw[popblue,->] (2,1.9) -- (1.82,1.62);
\draw[popblue,->] (2,-1.9) -- (2.18,-1.62);
\node[popblue,above,font=\small] at (2,2.0) {Objectif $L_1$};
\node[below,font=\small] at (2,-0.12) {$O_1$};
% Oculaire
\draw[poporange,very thick] (6.4,-1.3) -- (6.4,1.3);
\draw[poporange,->] (6.4,1.3) -- (6.22,1.02);
\draw[poporange,->] (6.4,-1.3) -- (6.58,-1.02);
\node[poporange,above,font=\small] at (6.4,1.4) {Oculaire $L_2$};
\node[below,font=\small] at (6.4,-0.12) {$O_2$};
% Foyers
\fill[popink] (5.2,0) circle (1.6pt);
\node[below,font=\small] at (5.2,-0.05) {$F'_1 \equiv F_2$};
\fill[popink] (7.6,0) circle (1.6pt);
\node[below,font=\small] at (7.6,-0.05) {$F'_2$};
\fill[popink] (3.2,0) circle (1.6pt);
\node[below,font=\small] at (3.2,-0.05) {$F_1$};
% Distances focales (echelle : f'1 = 3.2, f'2 = 1.2)
\draw[<->,popmuted] (2,-2.2) -- (5.2,-2.2) node[midway,below,font=\small] {$f'_1$};
\draw[<->,popmuted] (5.2,-2.2) -- (6.4,-2.2) node[midway,below,font=\small] {$f'_2$};
% Faisceau parallele incident incline (angle theta)
% Theta = -0.1875 rad. Ray heights at x=2: 1.0, 0.2, -0.6. At x=-1: 1.56, 0.76, -0.04.
\draw[popgreen,thick,->] (-1.0,1.56) -- (2,1.0);
\draw[popgreen,thick,->] (-1.0,0.76) -- (2,0.2);
\draw[popgreen,thick,->] (-1.0,-0.04) -- (2,-0.6);
\node[popgreen,above left,font=\small] at (-0.6,1.5) {faisceau parallèle issu de $B$ (à l'infini)};
% Angle theta
\draw[popgreen] (-0.2,1.28) arc[start angle=180,end angle=193,radius=1.0];
\node[popgreen,font=\small] at (-0.35,1.05) {$\theta$};
\draw[popgreen,dashed] (-1.0,1.30) -- (1.4,1.30);
% Convergence vers B1 dans le plan focal image (x = 5.2, y = -0.6)
\draw[popgreen,thick] (2,1.0) -- (5.2,-0.6);
\draw[popgreen,thick] (2,0.2) -- (5.2,-0.6);
\draw[popgreen,thick] (2,-0.6) -- (5.2,-0.6);
% Plan focal image
\draw[popmuted,dashed] (5.2,-1.9) -- (5.2,1.9);
% Image intermediaire A1B1
\fill[popink] (5.2,-0.6) circle (1.6pt);
\node[popink,right,font=\small] at (5.25,-0.67) {$B_1$};
\node[popink,above right,font=\small] at (5.2,0.05) {$A_1$};
% Faisceau emergent parallele (pente m = 0.5)
% Ray 1: through center O2 (6.4, 0). Enters at (5.2,-0.6), exits (6.4,0) slope 0.5. At x=9.6, y=1.6.
\draw[popgreen,thick] (5.2,-0.6) -- (6.4,0) -- (9.6,1.6);
% Ray 2: incident parallel to axis on ocular. Hits ocular at y=-0.6. Refracts through F'2 (7.6,0). Slope 0.5. At x=9.6, y=1.0.
\draw[popgreen,thick] (5.2,-0.6) -- (6.4,-0.6) -- (7.6,0) -- (9.6,1.0);
% Ray 3: arbitrary, hits ocular at y=0.6. Exits slope 0.5. At x=9.6, y=2.2.
\draw[popgreen,thick] (5.2,-0.6) -- (6.4,0.6) -- (9.6,2.2);
% Angle theta'
\draw[popgreen,dashed] (6.4,1.0) -- (9.6,1.0); % reference horizontal through ray 2 at ocular exit? No, ray 2 exits at -0.6.
% Better: draw angle at eye level.
\draw[popgreen,dashed] (9.6,1.0) -- (9.6,1.6); % vertical ref? No.
% Angle theta' measured from axis.
\draw[popgreen] (8.6,1.0) arc[start angle=0,end angle=14,radius=2.0]; % approx atan(0.5)=26.5 deg. radius 2.0 -> arc length.
\node[popgreen,font=\small] at (8.9,1.3) {$\theta'$};
\node[popgreen,font=\small] at (9.0,2.3) {faisceau émergent parallèle};
% Oeil
\draw[popdark] (9.7,1.3) ellipse (0.28 and 0.42);
\node[popdark,below,font=\small] at (9.7,0.85) {œil};
\end{tikzpicture}
\end{center}
\legende{Lunette afocale : le faisceau parallèle incliné de $\theta$ converge en $B_1$ dans le plan focal image de l'objectif ; comme $F'_1 \equiv F_2$, l'oculaire renvoie un faisceau parallèle incliné de $\theta'$, reçu par l'œil sans accommodation.}
\end{popfigure}

\begin{attention}[Deux pièges classiques]
\begin{itemize}
\item \textbf{Ne pas confondre les foyers} : c'est le foyer \emph{image} de l'objectif ($F'_1$) qui coïncide avec le foyer \emph{objet} de l'oculaire ($F_2$). Écrire « $F_1 = F'_2$ » est une erreur grave qui fausse tout le tracé.
\item \textbf{Le système afocal n'a pas de distance focale} : vouloir appliquer la relation de conjugaison $\frac{1}{\overline{O A'}} - \frac{1}{\overline{OA}} = \frac{1}{f'}$ à la lunette entière est dépourvu de sens. On applique les formules \emph{séparément} à l'objectif puis à l'oculaire.
\end{itemize}
\end{attention}

\courssub{Longueur de la lunette afocale}

La condition $F'_1 \equiv F_2$ fixe entièrement la géométrie de l'instrument. La distance entre les deux lentilles, c'est-à-dire la \cle{longueur de la lunette} (ou \emph{tirage}), vaut :
\[
L = \overline{O_1O_2} = \overline{O_1F'_1} + \overline{F_2O_2} = f'_1 + f'_2
\]

\begin{aretenir}[Longueur d'une lunette afocale]
Pour une lunette réglée afocale (objet à l'infini, œil normal sans accommodation) :
\[
\boxed{L = f'_1 + f'_2}
\]
La longueur de l'instrument est imposée par les distances focales : un objectif de grande focale donne mécaniquement une lunette longue.
\end{aretenir}

\begin{exemplebox}[Longueur d'une lunette scolaire]
Une lunette astronomique de laboratoire porte un objectif de distance focale $f'_1 = \SI{80}{cm}$ et un oculaire de distance focale $f'_2 = \SI{4}{cm}$. Réglée pour la vision à l'infini, sa longueur vaut :
\[
L = f'_1 + f'_2 = 80 + 4 = \SI{84}{cm}
\]
Le tube doit donc mesurer environ $\SI{84}{cm}$ : c'est pourquoi les lunettes astronomiques sont des instruments allongés, contrairement aux petites longues-vues à prismes redresseurs.
\end{exemplebox}

\begin{exoresolu}[Réglage d'une lunette de TP]
Énoncé. Au laboratoire du lycée, on dispose d'un objectif de vergence $C_1 = \num{1,25}\ \delta$ et d'un oculaire de vergence $C_2 = \num{20}\ \delta$. On veut observer la Lune avec un œil normal sans accommoder. (1) Calculer les distances focales des deux lentilles. (2) À quelle distance doit-on placer l'oculaire de l'objectif ? (3) Où se forme l'image intermédiaire de la Lune ?
\tcblower
Solution détaillée.
(1) La distance focale image est l'inverse de la vergence :
\[
f'_1 = \frac{1}{C_1} = \frac{1}{\num{1,25}} = \SI{0,80}{m} = \SI{80}{cm} \qquad ; \qquad f'_2 = \frac{1}{C_2} = \frac{1}{20} = \SI{0,05}{m} = \SI{5}{cm}
\]
(2) La Lune est à l'infini et l'œil ne doit pas accommoder : la lunette doit être afocale, donc
\[
L = f'_1 + f'_2 = 80 + 5 = \SI{85}{cm}
\]
(3) La Lune étant à l'infini, son image intermédiaire $A_1B_1$ se forme dans le plan focal image de l'objectif, à $\SI{80}{cm}$ derrière $O_1$ — soit exactement $\SI{5}{cm}$ devant l'oculaire, dans son plan focal objet. Tout est cohérent : l'oculaire rejette alors l'image définitive à l'infini.
\end{exoresolu}

\courssub{La mise au point : régler le tirage}

Dans la pratique, les lentilles ont des focales fixes ; c'est la \emph{distance} objectif–oculaire que l'on ajuste. L'oculaire est monté dans une petite lunette coulissante (le \cle{porte-oculaire}) que l'on fait glisser : c'est la \cle{mise au point}, et la distance réglée s'appelle le \cle{tirage}.

\begin{methode}[Mise au point d'une lunette pour un œil normal]
\begin{enumerate}
\item Pointer la lunette vers l'objet à observer et placer l'œil derrière l'oculaire.
\item Régler d'abord l'oculaire sur la position théorique $L = f'_1 + f'_2$ (repères gravés sur le tube).
\item Ajuster finement le tirage par petites translations de l'oculaire jusqu'à obtenir une image \emph{nette}, perçue sans effort : $F'_1$ et $F_2$ coïncident alors, l'image définitive est à l'infini.
\item L'œil observe alors \textbf{sans accommoder} : on peut observer longtemps sans fatigue.
\end{enumerate}
\end{methode}

\textbf{Et si l'objet n'est pas à l'infini ?} C'est le cas lorsqu'on vise le clocher d'une église à quelques centaines de mètres, ou un objet du laboratoire. L'objectif donne alors de l'objet une image intermédiaire située \emph{un peu au-delà} de son plan focal image (d'après la relation de conjugaison, un objet qui se rapproche voit son image s'éloigner de la lentille). Pour ramener cette image dans le plan focal objet de l'oculaire, il faut donc \textbf{éloigner l'oculaire de l'objectif} : on \emph{allonge légèrement le tirage}. Plus l'objet est proche, plus l'allongement est important — c'est exactement ce que fait le photographe quand sa bague de mise au point « sort » l'objectif pour photographier de près.

\begin{popfigure}
\begin{center}
\begin{tikzpicture}[scale=1.0,>=Stealth]
% --- Cas 1 : objet a l'infini ---
\draw[popmuted,->] (0,0) -- (8.6,0);
\draw[popblue,very thick] (1.5,-1.0) -- (1.5,1.0);
\node[popblue,below,font=\small] at (1.5,-1.05) {$L_1$};
\draw[poporange,very thick] (5.7,-0.8) -- (5.7,0.8);
\node[poporange,below,font=\small] at (5.7,-0.85) {$L_2$};
\fill[popink] (4.5,0) circle (1.5pt);
\node[above,font=\small] at (4.5,0.05) {$F'_1 \equiv F_2$};
\draw[<->,popdark,thick] (1.5,1.25) -- (5.7,1.25) node[midway,above,font=\small] {$L = f'_1 + f'_2$};
\node[font=\small\bfseries,popdark] at (4.2,-1.6) {Objet à l'infini (astre) : tirage minimal};
% --- Cas 2 : objet rapproche ---
\begin{scope}[yshift=-3.4cm]
\draw[popmuted,->] (0,0) -- (8.6,0);
\draw[popblue,very thick] (1.5,-1.0) -- (1.5,1.0);
\node[popblue,below,font=\small] at (1.5,-1.05) {$L_1$};
\draw[poporange,very thick] (6.3,-0.8) -- (6.3,0.8);
\node[poporange,below,font=\small] at (6.3,-0.85) {$L_2$};
\draw[poporange,dashed,thick] (5.7,-0.8) -- (5.7,0.8);
\draw[->,poporange,thick] (5.7,1.0) -- (6.3,1.0) node[midway,above,font=\small] {$\Delta$};
\fill[popink] (4.5,0) circle (1.5pt);
\node[above,font=\small] at (4.5,0.05) {$F'_1$};
\fill[popink] (5.1,0) circle (1.5pt);
\node[above,font=\small] at (5.15,-0.02) {$F_2$};
\fill[popgreen] (5.1,0.45) circle (1.5pt);
\node[popgreen,right,font=\small] at (5.15,0.45) {$B_1$ (image reculée)};
\draw[<->,popdark,thick] (1.5,1.45) -- (6.3,1.45) node[midway,above,font=\small] {$L' = f'_1 + f'_2 + \Delta$};
\node[font=\small\bfseries,popdark] at (4.2,-1.6) {Objet rapproché : on allonge le tirage de $\Delta$};
\end{scope}
\end{tikzpicture}
\end{center}
\legende{Mise au point : pour un objet à l'infini, $F'_1 \equiv F_2$ et le tirage vaut $f'_1 + f'_2$. Pour un objet rapproché, l'image intermédiaire recule derrière $F'_1$ : on translate l'oculaire vers l'arrière pour que $F_2$ coïncide avec la nouvelle position de l'image.}
\end{popfigure}

\courssub{Le cercle oculaire : où placer l'œil ?}

Dernière question pratique : une fois la mise au point faite, où l'œil doit-il se placer exactement derrière l'oculaire ? Toute la lumière qui entre dans la lunette passe par la monture circulaire de l'objectif, de diamètre $D$. Cette monture joue le rôle de \emph{diaphragme} : c'est elle qui limite les faisceaux. Son image à travers l'oculaire est un petit cercle lumineux situé un peu derrière $F'_2$, appelé \cle{cercle oculaire} (ou \emph{pupille de sortie}).

\begin{aretenir}[Cercle oculaire]
Le \cle{cercle oculaire} est l'image de la monture de l'objectif à travers l'oculaire. C'est le cercle de plus petit diamètre que traversent tous les faisceaux émergents. C'est là que l'on doit placer la pupille de l'œil : on y reçoit \textbf{le maximum de lumière} et le \textbf{champ le plus large}. Son diamètre, de l'ordre de quelques millimètres, est idéalement adapté à la pupille de l'œil dans l'obscurité (environ $\SI{6}{mm}$).
\end{aretenir}

Tu l'as sans doute remarqué en regardant dans des jumelles : quand l'œil est trop loin ou mal centré, l'image se rétrécit, s'assombrit ou présente un vignettage (bords noirs). Placé au niveau du cercle oculaire, l'œil capte en revanche la totalité du faisceau : l'image est lumineuse et le champ maximal.

\begin{savaistu}[Des lunettes géantes pour scruter le ciel]
La condition $L = f'_1 + f'_2$ explique les dimensions démesurées des lunettes historiques : la grande lunette de l'observatoire de Yerkes (États-Unis, 1897) possède un objectif de $f'_1 \approx \SI{19}{m}$ — le tube mesure donc près de $\SI{19}{m}$ de long ! Aujourd'hui encore, les astronomes amateurs camerounais qui observent le ciel limpide des savanes de l'Adamaoua utilisent des lunettes dont la longueur trahit directement la focale de leur objectif. Et lors des éclipses visibles au Cameroun, c'est la mise au point soigneuse du tirage qui fait la différence entre une tache floue et un croissant de Soleil parfaitement net (projeté sur écran, jamais observé directement !).
\end{savaistu}

\begin{exercice}[Pour s'entraîner]
Une lunette afocale est constituée d'un objectif de distance focale $f'_1 = \SI{60}{cm}$ et d'un oculaire de vergence $C_2 = \num{25}\ \delta$.
\begin{enumerate}
\item Calculer la longueur $L$ de la lunette réglée pour un œil normal observant une étoile.
\item On observe maintenant un oiseau perché à $\SI{30}{m}$ de l'objectif. En utilisant la relation de conjugaison sur l'objectif, déterminer la position de l'image intermédiaire et en déduire de combien il faut allonger le tirage.
\item Où l'observateur doit-il placer son œil, et pourquoi ?
\end{enumerate}
\end{exercice}

\begin{corrige}[Exercice 1]
1. $f'_2 = 1/C_2 = 1/25 = \SI{0,04}{m} = \SI{4}{cm}$. Réglage afocal : $L = f'_1 + f'_2 = 60 + 4 = \SI{64}{cm}$.

2. Relation de conjugaison pour l'objectif (distances algébriques : $\overline{O_1A} < 0$ pour un objet réel) :
\[
\frac{1}{\overline{O_1A_1}} - \frac{1}{\overline{O_1A}} = \frac{1}{f'_1}
\]
Avec $\overline{O_1A} = -\SI{30}{m}$ et $f'_1 = \SI{0,60}{m}$ :
\[
\frac{1}{\overline{O_1A_1}} = \frac{1}{f'_1} + \frac{1}{\overline{O_1A}} = \frac{1}{0,60} + \frac{1}{-30} = \num{1,6667} - \num{0,0333} = \num{1,6333}\ \si{m^{-1}}
\]
d'où $\overline{O_1A_1} \approx \SI{0,612}{m} = \SI{61,2}{cm}$. L'image intermédiaire se forme à $\SI{61,2}{cm}$ de l'objectif, soit $\SI{1,2}{cm}$ derrière $F'_1$. Il faut donc allonger le tirage de $\Delta \approx \SI{1,2}{cm}$ : $L' \approx \SI{65,2}{cm}$.

3. L'œil doit être placé au niveau du \cle{cercle oculaire}, image de la monture de l'objectif à travers l'oculaire, juste derrière le plan focal image de l'oculaire : c'est la position où la pupille reçoit tout le faisceau émergent, garantissant luminosité et champ maximaux.
\end{corrige}

En définitive, la lunette afocale réalise une prouesse élégante : deux lentilles, une seule condition géométrique ($F'_1 \equiv F_2$), et la lumière d'un astre lointain entre et sort de l'instrument en faisceaux parallèles — mais avec une inclinaison amplifiée. C'est cette amplification angulaire, source du grossissement, que nous quantifierons dans la section 4.

\coursec{Tracé de la marche d'un faisceau lumineux}

Dans la section précédente, nous avons établi le principe de la lunette \cle{afocale} : l'objectif $L_1$ forme d'un astre très lointain une image intermédiaire $A_1B_1$ dans son plan focal image, et l'oculaire $L_2$, dont le foyer objet $F_2$ coïncide avec $F'_1$, renvoie cette image à l'infini où l'\oe il l'observe sans accommoder. Cette description reste toutefois abstraite tant qu'on ne l'a pas \emph{dessinée}. Or, en optique géométrique, savoir tracer la marche des rayons, c'est comprendre l'instrument : le tracé justifie à lui seul le grossissement, le sens de l'image et la manière de faire la mise au point. C'est l'objet de cette section, qui correspond à l'une des compétences exigibles au Probatoire : \emph{tracer la marche d'un faisceau lumineux à travers une lunette astronomique}.

\courssub{Le problème posé : un faisceau incliné issu d'un point à l'infini}

Considérons un astre observé à travers la lunette, par exemple la Lune vue depuis Yaoundé. L'astre est à distance finie mais immense devant les dimensions de l'instrument : chacun de ses points est pratiquement \cle{à l'infini}. Le point $A$ de l'astre situé sur l'axe optique envoie sur l'objectif un faisceau de rayons \emph{parallèles à l'axe}. Un autre point $B$ de l'astre, situé hors de l'axe (le bord du disque lunaire, par exemple), envoie lui aussi un faisceau de rayons parallèles entre eux — car $B$ est à l'infini — mais ce faisceau est \emph{incliné} d'un angle $\theta$ par rapport à l'axe optique. Cet angle $\theta$ est la \cle{position angulaire} du point $B$ (ou la moitié du diamètre apparent si $B$ est au bord du disque).

Toute la question est donc la suivante : \emph{que devient un faisceau parallèle incliné de $\theta$ lorsqu'il traverse successivement l'objectif $L_1$ puis l'oculaire $L_2$ d'une lunette afocale ?} Nous allons répondre étape par étape, en n'utilisant que les deux règles fondamentales des lentilles minces convergentes :

\begin{itemize}
\item un faisceau incident \emph{parallèle} converge, après traversée d'une lentille convergente, en un point de son \emph{plan focal image} ;
\item un rayon passant par le \emph{centre optique} d'une lentille mince n'est \emph{pas dévié}.
\end{itemize}

\courssub{Étape 1 : le faisceau incident parallèle incliné de $\theta$}

Le point $B$ de l'astre étant à l'infini, tous les rayons qu'il envoie sur l'objectif sont parallèles entre eux et font l'angle $\theta$ avec l'axe optique. On en trace en général \emph{trois} : un rayon qui passera par le centre optique $O_1$ (le \cle{rayon utile}), et deux rayons qui frappent la lentille $L_1$ de part et d'autre de $O_1$.

\begin{attention}[Un faisceau incliné ne converge pas en $F'_1$]
L'erreur la plus fréquente — et la plus lourde de conséquences au Probatoire — consiste à faire converger \emph{tout} faisceau parallèle au foyer image $F'_1$ situé sur l'axe. C'est faux : seul un faisceau parallèle \textbf{à l'axe optique} converge en $F'_1$. Un faisceau parallèle \emph{incliné} de $\theta$ converge en un point $B_1$ du \emph{plan focal image}, mais \textbf{hors de l'axe}. Retiens : le plan focal image est un \emph{plan}, pas un point ; $F'_1$ n'en est que le point situé sur l'axe.
\end{attention}

\courssub{Étape 2 : convergence en $B_1$, dans le plan focal image de l'objectif}

Après traversée de l'objectif $L_1$, les trois rayons convergent en un même point $B_1$ du plan focal image de $L_1$. Mais \emph{où} exactement ? C'est le rayon utile qui le dit : le rayon passant par le centre optique $O_1$ n'est pas dévié ; il poursuit donc en ligne droite avec son inclinaison $\theta$ et vient percer le plan focal image en $B_1$. Les deux autres rayons sont alors déviés par la lentille de sorte à passer eux aussi par $B_1$.

La position de $B_1$ se calcule immédiatement. Dans le triangle rectangle formé par $O_1$, $F'_1$ et $B_1$ :
\[
\tan\theta = \frac{F'_1B_1}{O_1F'_1} = \frac{A_1B_1}{f'_1}
\qquad\text{soit, pour les petits angles ($\theta$ en radians),}\qquad
A_1B_1 \approx f'_1\,\theta.
\]
\begin{propriete}[Image d'un point à l'infini hors axe]
Tout point objet $B$ situé à l'infini, vu sous l'angle $\theta$ par rapport à l'axe, a pour image par l'objectif le point $B_1$ du plan focal image de $L_1$, situé à la distance $A_1B_1 = f'_1\tan\theta$ de l'axe, \emph{du côté opposé} à $B$ par rapport à l'axe. L'image intermédiaire est donc déjà \textbf{renversée} par l'objectif seul.
\end{propriete}

\courssub{Étape 3 : l'oculaire renvoie un faisceau parallèle incliné de $\theta'$}

La lunette est réglée de façon \emph{afocale} : le foyer objet $F_2$ de l'oculaire coïncide avec le foyer image $F'_1$ de l'objectif. Le point $B_1$, qui se trouve dans le plan focal image de $L_1$, se trouve donc aussi dans le \emph{plan focal objet} de $L_2$. Or $B_1$ joue maintenant le rôle d'\textbf{objet} pour l'oculaire.

\begin{propriete}[Objet dans le plan focal objet]
Tout point objet situé dans le plan focal objet d'une lentille convergente a son image \emph{à l'infini} : les rayons issus de ce point émergent de la lentille en un faisceau \textbf{parallèle}.
\end{propriete}

Chacun des trois rayons arrivant en $B_1$ se comporte, après $B_1$, comme un rayon issu de $B_1$ : à la sortie de l'oculaire, ils ressortent donc \emph{parallèles entre eux}. Reste à déterminer leur direction commune $\theta'$. On utilise à nouveau le rayon utile, cette fois pour l'oculaire : le rayon (réel ou virtuel) passant par le centre optique $O_2$ n'est pas dévié. La droite $(B_1O_2)$ fixe donc la direction d'émergence, et tous les rayons émergents sont parallèles à $(B_1O_2)$. Dans le triangle $O_2F_2B_1$ :
\[
\tan\theta' = \frac{F_2B_1}{O_2F_2} = \frac{A_1B_1}{f'_2}
\qquad\text{soit, pour les petits angles,}\qquad
\theta' \approx \frac{A_1B_1}{f'_2}.
\]

\courssub{Le tracé complet, en quatre temps}

\begin{methode}[Tracé de la marche d'un faisceau à travers une lunette afocale]
\begin{enumerate}
\item \textbf{Faisceau incident} : tracer l'axe optique, les deux lentilles $L_1$ et $L_2$ distantes de $f'_1+f'_2$, placer $F'_1$ et $F_2$ confondus. Tracer trois rayons parallèles entre eux, inclinés de $\theta$, arrivant sur $L_1$ (l'un d'eux passe par $O_1$).
\item \textbf{Convergence en $B_1$} : prolonger le rayon passant par $O_1$ sans le dévier jusqu'au plan focal image de $L_1$ ; l'intersection est $B_1$. Faire converger les deux autres rayons en $B_1$ (changement de pente au niveau de $L_1$).
\item \textbf{Point focal objet de $L_2$} : vérifier que $B_1$ est bien dans le plan focal objet de l'oculaire (condition afocale $F'_1 \equiv F_2$).
\item \textbf{Faisceau émergent} : tracer la droite $(B_1O_2)$ qui n'est pas déviée ; tous les rayons émergent de $L_2$ parallèles à cette droite, avec l'inclinaison $\theta'$. L'\oe il placé derrière l'oculaire reçoit un faisceau parallèle : il observe sans accommoder.
\end{enumerate}
\end{methode}

\begin{popfigure}
\begin{center}
\begin{tikzpicture}[scale=1.0,>=Stealth]
% Axe optique
\draw[popmuted,->] (-1.5,0) -- (10.5,0) node[below left] {\footnotesize axe optique};
% Lentille L1 (objectif) en x=0
\draw[popblue,thick,<->] (0,-2.2) -- (0,2.2);
\node[popblue,above] at (0,2.2) {\footnotesize $L_1$ objectif};
\fill (0,0) circle (1.5pt) node[below right,popdark] {\footnotesize $O_1$};
% Lentille L2 (oculaire) en x=6.5
\draw[popgreen!70!black,thick,<->] (6.5,-1.6) -- (6.5,1.6);
\node[popgreen!70!black,above] at (6.5,1.6) {\footnotesize $L_2$ oculaire};
\fill (6.5,0) circle (1.5pt) node[below right,popdark] {\footnotesize $O_2$};
% F'1 = F2 en x=5
\fill[poppink] (5,0) circle (1.5pt);
\node[poppink,below] at (5,-0.05) {\footnotesize $F'_1 \equiv F_2$};
% Plan focal commun
\draw[poppink,dashed] (5,-1.6) -- (5,1.6);
% B1
\fill[poporange] (5,-1.0) circle (2pt) node[below right,poporange] {\footnotesize $B_1$};
% ETAPE 1 : faisceau incident parallele (popblue), pente -0.2
\draw[popblue,thick,->] (-1.4,0.28) -- (0,0);
\draw[popblue,thick,->] (-1.4,0.78) -- (0,0.5);
\draw[popblue,thick,->] (-1.4,1.78) -- (0,1.5);
\node[popblue,above] at (-1.2,1.7) {\footnotesize étape 1 : faisceau parallèle incliné de $\theta$};
% ETAPE 2 : convergence en B1 (poporange)
\draw[poporange,thick,->] (0,0) -- (5,-1.0);
\draw[poporange,thick,->] (0,0.5) -- (5,-1.0);
\draw[poporange,thick,->] (0,1.5) -- (5,-1.0);
\node[poporange] at (2.6,-1.35) {\footnotesize étape 2 : convergence en $B_1$};
% ETAPE 3 : faisceau emergent parallele (popgreen), direction (B1 O2) pente +2/3
% Rayon chef : B1(5,-1) -> O2(6.5,0) -> suite
\draw[popgreen!70!black,thick,->] (5,-1.0) -- (6.5,0) -- (9.8,2.2);
% Rayon haut : passe par L2 en y=0.5
\draw[popgreen!70!black,thick,->] (5,-1.0) -- (6.5,0.5) -- (9.8,2.7);
% Rayon bas : passe par L2 en y=-0.5
\draw[popgreen!70!black,thick,->] (5,-1.0) -- (6.5,-0.5) -- (9.8,1.7);
\node[popgreen!70!black,above right] at (8.2,2.35) {\footnotesize étape 3 : faisceau parallèle incliné de $\theta'$};
\end{tikzpicture}
\end{center}
\legende{Construction pas à pas de la marche d'un faisceau à travers une lunette afocale : faisceau incident parallèle (bleu), convergence en $B_1$ dans le plan focal image de $L_1$ (orange), émergence en faisceau parallèle (vert). Les angles sont volontairement exagérés pour la lisibilité.}
\end{popfigure}

\courssub{Angles $\theta$ et $\theta'$ : le sens de l'image}

Regarde attentivement la construction : le faisceau incident arrive \emph{en descendant} vers l'axe (le point $B$ est vu au-dessus de l'axe, $\theta < 0$ si l'on adopte la convention trigonométrique), tandis que le faisceau émergent repart \emph{en montant} ($\theta' > 0$). Les angles $\theta$ et $\theta'$ sont de \emph{signes opposés}. Or l'\oe il interprète la direction des rayons émergents comme la direction de l'objet : il croit voir le point $B$ dans la direction $\theta'$, c'est-à-dire \emph{de l'autre côté de l'axe}.

\begin{aretenir}[Image renversée]
Dans une lunette astronomique à deux lentilles convergentes, $\theta$ et $\theta'$ sont de signes opposés : l'image définitive est \textbf{renversée} par rapport à l'objet. C'est sans importance en astronomie (il n'y a pas de « haut » ni de « bas » sur la Lune), mais gênant pour observer un paysage terrestre.
\end{aretenir}

\begin{popfigure}
\begin{center}
\begin{tikzpicture}[scale=1.0,>=Stealth]
\draw[popmuted,->] (-1.5,0) -- (10.5,0);
\draw[popblue,thick,<->] (0,-2.0) -- (0,2.0);
\node[popblue,above] at (0,2.0) {\footnotesize $L_1$};
\draw[popgreen!70!black,thick,<->] (6.5,-1.5) -- (6.5,1.5);
\node[popgreen!70!black,above] at (6.5,1.5) {\footnotesize $L_2$};
\fill (5,0) circle (1.5pt) node[below,poppink] {\footnotesize $F'_1\equiv F_2$};
\fill[poporange] (5,-1.0) circle (2pt) node[below right,poporange] {\footnotesize $B_1$};
% Rayon central incident
\draw[popblue,thick,->] (-1.4,0.28) -- (0,0);
\draw[poporange,thick,->] (0,0) -- (5,-1.0);
\draw[popgreen!70!black,thick,->] (5,-1.0) -- (6.5,0) -- (9.8,2.2);
% Arc angle theta en O1
\draw[popblue,->] (1.5,0) arc (0:-11.3:1.5);
\node[popblue] at (1.9,-0.4) {$\theta$};
% Arc angle theta' en O2
\draw[popgreen!70!black,->] (6.5,0.3) arc (0:33.7:1.5);
\node[popgreen!70!black] at (7.5,0.6) {$\theta'$};
\node[popdark,align=center] at (8.6,-0.5) {\footnotesize $\theta$ et $\theta'$ de signes opposés \\ \footnotesize $\Rightarrow$ image renversée};
\end{tikzpicture}
\end{center}
\legende{Schéma annoté des angles : le faisceau incident fait l'angle $\theta$ avec l'axe (mesuré au niveau de $O_1$), le faisceau émergent fait l'angle $\theta'$ (mesuré au niveau de $O_2$). Le rayon non dévié par $O_1$ fixe $B_1$ ; la droite $(B_1O_2)$ fixe $\theta'$.}
\end{popfigure}

\courssub{Application : construction à l'échelle}

\begin{exemplebox}[Construction à l'échelle $1/10$]
On donne une lunette afocale avec $f'_1 = \SI{50}{cm}$ et $f'_2 = \SI{5}{cm}$, éclairée par un faisceau issu d'un point $B$ à l'infini vu sous $\theta = 1°$. Réalisons la construction sur papier millimétré à l'échelle $1/10$ (\SI{1}{cm} sur le papier pour \SI{10}{cm} réels).

\textbf{Données reportées sur le papier :} $f'_1 \to \SI{5,0}{cm}$, $f'_2 \to \SI{0,5}{cm}$, distance $O_1O_2 = f'_1+f'_2 \to \SI{5,5}{cm}$.

\textbf{Position de $B_1$ :}
\[
A_1B_1 = f'_1\tan\theta = 50\times\tan 1° = 50\times\num{0,01746} \approx \SI{0,87}{cm},
\]
soit seulement $\SI{0,87}{mm}$ sur le papier : $B_1$ est très proche de l'axe, ce qui est réaliste (les diamètres apparents des astres sont minuscules ; pour le dessin, on exagérera $\theta$).

\textbf{Angle émergent :}
\[
\tan\theta' = \frac{A_1B_1}{f'_2} = \frac{\num{0,87}}{5} = \num{0,174}
\quad\Rightarrow\quad \theta' \approx \num{9,9}° \approx 10°.
\]
L'angle a été multiplié par environ $10$ : c'est le \cle{grossissement} $G = \theta'/\theta \approx f'_1/f'_2 = 10$, que nous étudierons en détail dans la section suivante. Le tracé le \emph{démontre} déjà géométriquement.
\end{exemplebox}

\begin{exoresolu}[Vérification par le calcul d'un tracé]
Énoncé. Une lunette afocale a pour objectif $f'_1 = \SI{80}{cm}$ et pour oculaire $f'_2 = \SI{2}{cm}$. Un rayon issu d'un point $B$ à l'infini passe par $O_1$ en faisant $\theta = \num{0,50}°$ avec l'axe. \textbf{1.} À quelle distance de l'axe ce rayon coupe-t-il le plan focal image de $L_1$ ? \textbf{2.} Sous quel angle $\theta'$ émerge-t-il de l'oculaire ? \textbf{3.} L'image est-elle droite ou renversée ?
\tcblower
Solution. \textbf{1.} Le rayon passant par $O_1$ n'est pas dévié ; il perce le plan focal image en $B_1$ avec
\[
A_1B_1 = f'_1\tan\theta = 80\times\tan(\num{0,50}°) = 80\times\num{0,008727} \approx \SI{0,70}{cm}.
\]
\textbf{2.} $B_1$ est dans le plan focal objet de $L_2$ (montage afocal) ; la direction d'émergence est donnée par $(B_1O_2)$ :
\[
\tan\theta' = \frac{A_1B_1}{f'_2} = \frac{\num{0,70}}{2} = \num{0,35}
\quad\Rightarrow\quad
\theta' \approx \num{19,3}°.
\]
Vérification : $G = f'_1/f'_2 = 40$ et $40\times\num{0,50}° = 20° \approx \num{19,3}°$ (l'écart vient de l'approximation des petits angles, moins bonne à $20°$). \textbf{3.} $\theta$ et $\theta'$ sont de signes opposés : l'image est \textbf{renversée}.
\end{exoresolu}

\begin{attention}[Pièges de rédaction au Probatoire]
\begin{itemize}
\item Les rayons changent de pente \emph{au niveau de la lentille}, jamais avant ni après : ne fais pas « casser » le rayon en dehors du plan de $L_1$ ou $L_2$.
\item N'oublie pas les \emph{flèches} indiquant le sens de propagation de la lumière sur chaque rayon.
\item Le faisceau émergent doit être \emph{rigoureusement parallèle} : si tes trois rayons émergents ne sont pas parallèles entre eux, le tracé est faux.
\item Précise sur le schéma que $F'_1$ et $F_2$ sont confondus : c'est la signature du réglage afocal.
\end{itemize}
\end{attention}

\begin{savaistu}[Et pour voir les images à l'endroit ?]
Le renversement de l'image est sans conséquence pour observer les astres, mais il est rédhibitoire pour observer un paysage, un match de football au stade de Japoma ou une régate sur le Wouri. Les \emph{lunettes terrestres} ajoutent donc une troisième lentille, dite \emph{lentille redresseuse}, qui renverse l'image une seconde fois — au prix d'un allongement de l'instrument. Les \emph{jumelles} résolvent le problème plus élégamment : chaque tube contient deux prismes (montage dit « de Porro ») qui redressent l'image par réflexions totales tout en repliant le trajet de la lumière, ce qui raccourcit l'instrument. C'est pourquoi tes jumelles sont compactes alors qu'une lunette terrestre ancienne était un long tube !
\end{savaistu}

En résumé, le tracé de la marche d'un faisceau à travers la lunette afocale repose sur deux outils seulement — le plan focal et le rayon non dévié par le centre optique — mais il livre tous les secrets de l'instrument : l'image intermédiaire $B_1$ en $f'_1\tan\theta$ de l'axe, l'émergence parallèle sous $\theta'$ tel que $\tan\theta' = A_1B_1/f'_2$, et le renversement de l'image. Il ne reste plus qu'à exploiter quantitativement le rapport $\theta'/\theta$ : ce sera le \emph{grossissement}, objet de la section suivante.

\coursec{Grossissement et puissance}

Dans la section précédente, nous avons tracé la marche d'un faisceau lumineux à travers une lunette afocale : un faisceau parallèle issu d'un point de l'astre, arrivant sous l'incidence $\theta$, ressort de l'oculaire sous la forme d'un faisceau parallèle incliné d'un angle $\theta'$, visiblement plus grand que $\theta$. L'image définitive est à l'infini : l'œil de l'observateur la regarde sans accommoder, donc sans fatigue. Mais une question demeure, et c'est la question que tout utilisateur d'instrument pose en premier : \emph{de combien la lunette grossit-elle ?} C'est précisément l'objet de cette section : définir rigoureusement le \cle{grossissement}, le calculer à partir des caractéristiques de l'instrument, puis introduire une grandeur complémentaire, la \cle{puissance}, très appréciée des examinateurs au Probatoire.

\courssub{Diamètre apparent : l'angle qui mesure la taille d'un astre}

\begin{experience}[Observer la Lune à l'œil nu]
Par une belle nuit claire à Yaoundé ou à Maroua, lève les yeux vers la pleine Lune. Elle te semble « grosse », mais cette impression est trompeuse : tends le bras, ton petit doigt suffit à la masquer entièrement ! La Lune a pourtant un diamètre réel d'environ $\SI{3474}{km}$. Ce qui compte pour l'œil, ce n'est pas la taille réelle de l'objet, mais l'\textbf{angle} sous lequel on le voit : c'est le \cle{diamètre apparent}. Pour la Lune, située à environ $\SI{384000}{km}$, cet angle vaut :
\[
\theta \approx \dfrac{\SI{3474}{km}}{\SI{384000}{km}} \approx \num{9,0e-3}~\text{rad} \approx 0,5^{\circ}.
\]
C'est un angle très petit : la rétine ne peut pas distinguer les détails lunaires (cratères, mers) car ils sont vus sous des angles encore plus faibles, inférieurs au pouvoir séparateur de l'œil (environ $1'$, soit un soixantième de degré).
\end{experience}

\begin{definition}[Diamètres apparents de l'objet et de l'image]
On note :
\begin{itemize}
\item $\theta$ le \cle{diamètre apparent} de l'astre observé \textbf{à l'œil nu} : c'est l'angle sous lequel on voit l'objet (situé à l'infini) sans instrument ;
\item $\theta'$ l'angle sous lequel on voit l'\textbf{image} de ce même astre \textbf{à travers la lunette}.
\end{itemize}
Ces deux angles sont des grandeurs \emph{algébriques} en optique géométrique, mais pour le grossissement on ne retient que leurs valeurs absolues. Ils s'expriment en radians (rad) dans les calculs.
\end{definition}

\begin{attention}[Ne pas confondre taille réelle et diamètre apparent]
Deux objets de tailles très différentes peuvent avoir le même diamètre apparent : la Lune ($\SI{3474}{km}$ à $\SI{384000}{km}$) et le Soleil ($\SI{1,4e6}{km}$ à $\SI{1,5e8}{km}$) sont vus sous quasiment le même angle, $0,5^{\circ}$ — c'est ce qui rend les éclipses totales possibles ! En astronomie, c'est toujours l'\textbf{angle} qui compte, jamais la taille absolue.
\end{attention}

\courssub{Grossissement d'une lunette astronomique}

\begin{definition}[Grossissement]
Le \cle{grossissement} d'une lunette astronomique est le rapport, en valeur absolue, de l'angle sous lequel on voit l'image à travers l'instrument à l'angle sous lequel on voit l'objet à l'œil nu :
\[
\boxed{G = \dfrac{|\theta'|}{|\theta|}}
\]
Le grossissement est un nombre \textbf{sans unité} (rapport de deux angles). On dit par exemple « cette lunette grossit 60 fois », ce qu'on note parfois $G = 60\times$ sur les catalogues.
\end{definition}

Autrement dit : si la Lune est vue à l'œil nu sous $0,5^{\circ}$ et à travers la lunette sous $30^{\circ}$, la lunette grossit $60$ fois. Le grossissement mesure le \emph{gain angulaire} apporté par l'instrument.

\courssub{Démonstration : expression du grossissement d'une lunette afocale}

Cette démonstration est un grand classique du Probatoire : sache la refaire proprement, avec le schéma et l'approximation des petits angles clairement annoncée.

\begin{experience}[Démonstration guidée : $G = f'_1/f'_2$]
Considérons une lunette \textbf{afocale} : le foyer image $F'_1$ de l'objectif coïncide avec le foyer objet $F_2$ de l'oculaire. Un point $B$ de l'astre, situé à l'infini, envoie sur l'objectif un faisceau parallèle faisant l'angle $\theta$ avec l'axe optique.

\textbf{Étape 1 — Image intermédiaire dans le plan focal de l'objectif.} L'image $B_1$ de $B$ donnée par l'objectif se forme dans son plan focal image, en $F'_1$. L'image intermédiaire $A_1B_1$ a pour taille, dans le triangle $O_1F'_1B_1$ rectangle en $F'_1$ :
\[
\tan\theta = \dfrac{A_1B_1}{O_1F'_1} = \dfrac{A_1B_1}{f'_1}.
\]
Comme l'astre est vu sous un très petit angle, on utilise l'\textbf{approximation des petits angles} ($\theta$ en radians) : $\tan\theta \approx \theta$. D'où :
\[
\theta \approx \dfrac{A_1B_1}{f'_1} \qquad (1)
\]

\textbf{Étape 2 — L'oculaire renvoie l'image à l'infini.} L'image intermédiaire $A_1B_1$ joue le rôle d'objet pour l'oculaire ; comme elle est dans son plan focal objet (lunette afocale), l'image définitive est rejetée à l'infini et vue sous l'angle $\theta'$. Dans le triangle $O_2F_2B_1$ rectangle en $F_2$ :
\[
\tan\theta' = \dfrac{A_1B_1}{O_2F_2} = \dfrac{A_1B_1}{f'_2}.
\]
L'angle $\theta'$ reste petit (quelques degrés), donc $\tan\theta' \approx \theta'$ :
\[
\theta' \approx \dfrac{A_1B_1}{f'_2} \qquad (2)
\]

\textbf{Étape 3 — Rapport des deux angles.} En divisant membre à membre (2) par (1), la taille $A_1B_1$ de l'image intermédiaire s'élimine :
\[
G = \dfrac{\theta'}{\theta} = \dfrac{A_1B_1/f'_2}{A_1B_1/f'_1} = \dfrac{f'_1}{f'_2}.
\]
\end{experience}

\begin{propriete}[Grossissement d'une lunette afocale — démonstration exigible]
Pour une lunette astronomique réglée de façon \textbf{afocale}, le grossissement est le rapport de la distance focale de l'objectif à celle de l'oculaire :
\[
\boxed{G = \dfrac{f'_1}{f'_2}}
\]
où $f'_1$ est la distance focale de l'\textbf{objectif} (côté astre) et $f'_2$ celle de l'\textbf{oculaire} (côté œil), exprimées dans la \emph{même unité}. $G$ est sans unité. La démonstration (triangles rectangles + approximation $\tan\alpha \approx \alpha$) est exigible au Probatoire.
\end{propriete}

\begin{popfigure}
\begin{center}
\begin{tikzpicture}[scale=1.0, >=Stealth]
% Axe optique
\draw[->, popmuted] (-1.5,0) -- (9.8,0) node[below left] {\footnotesize axe optique};
% Objectif L1
\draw[popblue, very thick] (0,-1.9) -- (0,1.9);
\draw[popblue, ->] (0,1.9) -- (0.25,1.75);
\draw[popblue, ->] (0,-1.9) -- (0.25,-1.75);
\node[popblue, above] at (0,2.0) {\footnotesize Objectif $L_1$};
\node[below, popdark] at (0,-0.15) {\footnotesize $O_1$};
\fill[popdark] (0,0) circle (1.3pt);
% F'1 = F2
\fill[poporange] (5,0) circle (1.6pt);
\node[below, poporange] at (5,-0.15) {\footnotesize $F'_1 = F_2$};
% Image intermediaire B1
\fill[popink] (5,1.5) circle (1.6pt);
\node[right, popink] at (5.1,1.5) {\footnotesize $B_1$};
\draw[popink, thick, ->] (5,0) -- (5,1.5);
\node[left, popink] at (4.9,0.75) {\footnotesize $A_1B_1$};
% Oculaire L2
\draw[poppurple, very thick] (7,-1.4) -- (7,1.4);
\draw[poppurple, ->] (7,1.4) -- (7.25,1.26);
\draw[poppurple, ->] (7,-1.4) -- (7.25,-1.26);
\node[poppurple, above] at (7,1.5) {\footnotesize Oculaire $L_2$};
\node[below, popdark] at (7,-0.15) {\footnotesize $O_2$};
\fill[popdark] (7,0) circle (1.3pt);
% F'2
\fill[poporange] (9,0) circle (1.6pt);
\node[below, poporange] at (9,-0.15) {\footnotesize $F'_2$};

% Faisceau entrant (angle theta) : 3 rayons paralleles
% Rayon principal (par centre O1)
\draw[popgreen, thick, ->] (-1.2,-0.36) -- (0,0) -- (5,1.5);
% Rayon marginal haut
\draw[popgreen, thick, ->] (-1.2,0.64) -- (0,1.0) -- (5,1.5);
% Rayon marginal bas
\draw[popgreen, thick, ->] (-1.2,-1.36) -- (0,-1.0) -- (5,1.5);

% Faisceau emergent (angle theta') : 3 rayons paralleles
% Rayon principal (par centre O2)
\draw[popgreen, thick, ->] (5,1.5) -- (7,0) -- (9.5,-1.125);
% Rayon par F'2 (entrant parallele a l'axe depuis B1)
\draw[popgreen, thick] (5,1.5) -- (7,1.5);
\draw[popgreen, thick, ->] (7,1.5) -- (9,0) -- (9.5,-0.375);
% Rayon symetrique bas
\draw[popgreen, thick] (5,1.5) -- (7,-0.5);
\draw[popgreen, thick, ->] (7,-0.5) -- (9,0) -- (9.5,-1.875);

% Angles theta (entree)
\draw[poporange, thick] (1.2,0) arc (0:16.7:1.2);
\node[poporange] at (1.7,0.25) {\footnotesize $\theta$};
% Angle theta' (sortie)
\draw[poporange, thick] (8.2,0) arc (180:196.7:1.2);
\node[poporange] at (7.8,-0.55) {\footnotesize $\theta'$};

% Distances focales
\draw[<->, popblue] (0,-2.3) -- (5,-2.3) node[midway, below] {\footnotesize $f'_1$};
\draw[<->, poppurple] (5,-2.3) -- (7,-2.3) node[midway, below] {\footnotesize $f'_2$};
\end{tikzpicture}
\end{center}
\legende{Géométrie de la démonstration : dans le triangle $O_1F'_1B_1$, $\theta \approx A_1B_1/f'_1$ ; dans le triangle $O_2F_2B_1$, $\theta' \approx A_1B_1/f'_2$. Le faisceau émerge de l'oculaire sous l'angle $\theta' > \theta$.}
\end{popfigure}

\begin{attention}[L'erreur classique : inverser le rapport]
Beaucoup d'élèves écrivent $G = f'_2/f'_1$ : c'est \textbf{faux} et cela se voit physiquement ! Retiens la logique : pour grossir beaucoup, il faut un objectif de \textbf{grande} focale (il forme une image intermédiaire grande, puisque $A_1B_1 = f'_1\,\theta$) et un oculaire de \textbf{courte} focale (c'est une loupe puissante qui examine cette image). Le grand nombre est donc au numérateur :
\[
G = \dfrac{f'_1 \ (\text{grande})}{f'_2 \ (\text{petite})}.
\]
Vérifie toujours que ton résultat est \emph{supérieur à 1} pour une lunette astronomique : un grossissement inférieur à 1 trahit immédiatement l'inversion.
\end{attention}

\begin{exemplebox}[Conséquence pratique : comment concevoir une lunette puissante]
La relation $G = f'_1/f'_2$ dicte la conception des instruments :
\begin{itemize}
\item une \textbf{longue-vue} de randonnée : $f'_1 = \SI{400}{mm}$, $f'_2 = \SI{20}{mm}$, soit $G = 20$ ;
\item une \textbf{lunette astronomique} d'amateur : $f'_1 = \SI{900}{mm}$, $f'_2 = \SI{15}{mm}$, soit $G = 60$ ;
\item en changeant simplement d'oculaire ($f'_2 = \SI{6}{mm}$), la même lunette atteint $G = 150$ : c'est pourquoi les astronomes possèdent une \emph{collection d'oculaires} !
\end{itemize}
Le prix à payer : une grande focale d'objectif signifie un instrument long et encombrant, et un fort grossissement exige une excellente stabilité (le moindre tremblement est grossi lui aussi).
\end{exemplebox}

\begin{popfigure}
\begin{center}
\begin{tikzpicture}
\begin{axis}[
width=0.85\linewidth, height=6.2cm,
xlabel={$f'_1$ (mm)}, ylabel={$G$},
xmin=0, xmax=1000, ymin=0, ymax=110,
xtick={0,200,400,600,800,1000},
ytick={0,20,40,60,80,100},
grid=both, grid style={popline!40},
legend pos=north west,
legend style={font=\footnotesize, draw=popline},
axis lines=left, thick
]
\addplot[popcurve, popblue, domain=0:1000, samples=50]{x/10};
\addlegendentry{$f'_2 = \SI{10}{mm}$}
\addplot[popcurve, poporange, domain=0:1000, samples=50]{x/15};
\addlegendentry{$f'_2 = \SI{15}{mm}$}
\addplot[popcurve, popgreen, domain=0:1000, samples=50]{x/25};
\addlegendentry{$f'_2 = \SI{25}{mm}$}
\addplot[mark=*, popink, mark size=1.8pt] coordinates {(900,60)};
\node[popink, anchor=west, font=\footnotesize] at (axis cs:910,60) {$(900\,;\,60)$};
\end{axis}
\end{tikzpicture}
\end{center}
\legende{Grossissement $G = f'_1/f'_2$ en fonction de la focale de l'objectif, pour trois oculaires. Les droites passent par l'origine : $G$ est proportionnel à $f'_1$, et la pente $1/f'_2$ est d'autant plus forte que l'oculaire est de courte focale. Le point marqué correspond à l'exemple $f'_1 = \SI{900}{mm}$, $f'_2 = \SI{15}{mm}$, $G = 60$.}
\end{popfigure}

\courssub{Puissance d'une lunette}

Au-delà du grossissement, les opticiens caractérisent un instrument par sa \cle{puissance} : de quel angle l'instrument « amplifie-t-il » un objet intermédiaire de taille donnée ?

\begin{definition}[Puissance d'une lunette]
La \cle{puissance} d'une lunette est le quotient de l'angle $\theta'$ sous lequel on voit l'image à travers l'instrument par la taille $A_1B_1$ de l'image intermédiaire (ou de l'objet pour l'oculaire) :
\[
P = \dfrac{\theta'}{A_1B_1}
\]
Pour une lunette \textbf{afocale}, la relation (2) de la démonstration donne immédiatement :
\[
\boxed{P = \dfrac{1}{f'_2}}
\]
Si $f'_2$ est exprimée en \textbf{mètres}, $P$ s'exprime en \textbf{dioptries} (symbole $\delta$) : $1~\delta = 1~\text{m}^{-1}$. La puissance de la lunette afocale est donc simplement la \emph{vergence de l'oculaire}.
\end{definition}

\begin{propriete}[Relation entre puissance et grossissement]
En combinant $G = f'_1/f'_2$ et $P = 1/f'_2$, on obtient :
\[
\boxed{G = P \times f'_1}
\]
avec $P$ en dioptries ($\text{m}^{-1}$) et $f'_1$ en \textbf{mètres}. Cette relation est très pratique pour les exercices à deux étapes : on calcule d'abord $P$, puis $G$.
\end{propriete}

\begin{attention}[Unités : le piège du millimètre]
La puissance $P = 1/f'_2$ n'est en dioptries que si $f'_2$ est en \textbf{mètres}. Avec $f'_2 = \SI{15}{mm}$, il faut convertir : $f'_2 = \SI{0,015}{m}$, d'où $P = 1/0,015 \approx 66,7~\delta$. Écrire $P = 1/15$ serait une erreur de \textbf{trois ordres de grandeur} (facteur 1000) ! De même, dans $G = P f'_1$, convertis $f'_1$ en mètres : $G = 66,7 \times 0,900 = 60$. Analyse dimensionnelle : $[P]\times[f'_1] = \text{m}^{-1}\times\text{m} = 1$, sans dimension — cohérent avec un grossissement sans unité.
\end{attention}

\begin{exoresolu}[Lunette d'un club d'astronomie de Douala]
Énoncé. Le club d'astronomie d'un lycée de Douala possède une lunette dont l'objectif a pour distance focale $f'_1 = \SI{900}{mm}$ et l'oculaire $f'_2 = \SI{15}{mm}$, réglée de façon afocale.
\begin{enumerate}
\item Calculer le grossissement $G$ de la lunette.
\item Calculer sa puissance $P$ en dioptries et retrouver $G$ par la relation $G = Pf'_1$.
\item La pleine Lune est vue à l'œil nu sous un diamètre apparent $\theta \approx 0,5^{\circ}$. Sous quel angle $\theta'$ est-elle vue à travers la lunette ? Conclure.
\end{enumerate}
\tcblower
Solution détaillée.
\begin{enumerate}
\item Les deux focales étant dans la même unité, on applique directement :
\[
G = \dfrac{f'_1}{f'_2} = \dfrac{900}{15} = 60.
\]
La lunette grossit 60 fois (résultat sans unité, supérieur à 1 : cohérent).
\item Conversion en mètres : $f'_2 = \SI{0,015}{m}$. D'où :
\[
P = \dfrac{1}{f'_2} = \dfrac{1}{\num{0,015}} \approx 66,7~\delta.
\]
Vérification du grossissement : $f'_1 = \SI{0,900}{m}$, donc
\[
G = P \times f'_1 = 66,7 \times 0,900 = 60.
\]
On retrouve bien la même valeur : les deux relations sont cohérentes.
\item Par définition du grossissement, $\theta' = G\,\theta = 60 \times 0,5^{\circ} = 30^{\circ}$.
\end{enumerate}
\textbf{Conclusion.} La Lune, qui tenait à l'œil nu derrière un petit doigt tendu, est vue à travers la lunette sous un angle de $30^{\circ}$ : elle remplit une large partie du champ de l'instrument, et les cratères, vus eux aussi 60 fois plus « grands » angulairement, deviennent discernables puisqu'ils dépassent le pouvoir séparateur de l'œil ($1'$).
\end{exoresolu}

\courssub{Limite pratique du grossissement}

Peut-on augmenter $G$ indéfiniment en choisissant un oculaire de focale de plus en plus courte ? Non, et c'est une remarque qualitative à savoir formuler au Probatoire.

\begin{aretenir}[Grossissement utile et diamètre de l'objectif]
Le grossissement \emph{utile} d'une lunette est limité par le \textbf{diamètre $D$ de l'objectif}, pour deux raisons :
\begin{itemize}
\item \textbf{luminosité} : l'objectif collecte la lumière de l'astre sur sa surface ; au-delà d'un certain grossissement, l'image devient trop sombre pour être observable ;
\item \textbf{diffraction} : la lumière, en traversant l'ouverture circulaire de l'objectif, subit une diffraction d'autant plus marquée que $D$ est petit ; les détails trop fins sont brouillés, et grossir davantage ne révèle rien de nouveau (image « floue » agrandie).
\end{itemize}
Règle pratique des astronomes amateurs : le grossissement maximal utile est de l'ordre de $1,5$ à $2$ fois le diamètre de l'objectif exprimé en millimètres (soit $G_{\max} \approx 150$ pour $D = \SI{80}{mm}$). Un fort grossissement exige donc un \emph{objectif de grand diamètre}, pas seulement de grande focale.
\end{aretenir}

\begin{savaistu}[Des lunettes aux grands télescopes africains]
Le principe « grande focale pour fort grossissement » atteint vite ses limites avec les lentilles : une lentille de plusieurs mètres de focale et de grand diamètre est lourde, coûteuse et se déforme sous son propre poids. C'est pourquoi tous les grands instruments modernes sont des \textbf{télescopes} à miroirs : un miroir concave joue le rôle de l'objectif, avec des focales de plusieurs dizaines de mètres. L'Afrique n'est pas en reste : l'Observatoire de l'\textbf{Oukaïmeden}, dans le Haut Atlas marocain, accueille des télescopes professionnels qui ont participé à la découverte de comètes et d'exoplanètes, et l'Afrique du Sud héberge le télescope SALT, l'un des plus grands de l'hémisphère sud. Le jour où un observatoire verra le jour sur les hauts plateaux de l'Adamaoua, réputés pour leurs nuits limpides, ce sera encore la même physique — celle de cette leçon — qui guidera son objectif vers les étoiles.
\end{savaistu}

\begin{exercice}[Pour t'entraîner]
Une lunette afocale a une puissance $P = 40~\delta$ et un grossissement $G = 32$.
\begin{enumerate}
\item Calculer la distance focale $f'_2$ de l'oculaire (en mm).
\item Calculer la distance focale $f'_1$ de l'objectif (en mm).
\item En déduire la longueur totale de la lunette (distance objectif--oculaire).
\item Jupiter est vu à l'œil nu sous un angle maximal de $50''$ (secondes d'arc). Sous quel angle est-il vu à travers cette lunette ? Exprime le résultat en minutes d'arc ($1' = 60''$).
\end{enumerate}
\end{exercice}

\begin{corrige}[Exercice]
\begin{enumerate}
\item $P = 1/f'_2$ donc $f'_2 = 1/P = 1/40 = \SI{0,025}{m} = \SI{25}{mm}$.
\item $G = Pf'_1$ donc $f'_1 = G/P = 32/40 = \SI{0,80}{m} = \SI{800}{mm}$. (Vérification : $G = f'_1/f'_2 = 800/25 = 32$.)
\item Pour une lunette afocale, la longueur est $L = f'_1 + f'_2 = 800 + 25 = \SI{825}{mm}$.
\item $\theta' = G\,\theta = 32 \times 50'' = 1600'' = 1600/60 \approx 26,7'$, soit environ $0,44^{\circ}$ : les bandes nuageuses et les satellites galiléens de Jupiter deviennent accessibles.
\end{enumerate}
\end{corrige}

\begin{aretenir}[L'essentiel de la section]
\begin{itemize}
\item Diamètre apparent $\theta$ : angle sous lequel l'astre est vu à l'œil nu ; $\theta'$ : angle sous lequel l'image est vue à travers la lunette.
\item Grossissement : $G = \dfrac{|\theta'|}{|\theta|}$, sans unité ; pour une lunette afocale : $G = \dfrac{f'_1}{f'_2}$ (démonstration exigible : triangles rectangles et $\tan\alpha \approx \alpha$).
\item Fort grossissement $\Rightarrow$ objectif de \textbf{grande} focale, oculaire de \textbf{courte} focale.
\item Puissance : $P = \dfrac{\theta'}{A_1B_1} = \dfrac{1}{f'_2}$, en dioptries ($\delta$) si $f'_2$ est en mètres ; relation $G = P\,f'_1$.
\item Le grossissement utile est limité par le diamètre de l'objectif (luminosité et diffraction).
\end{itemize}
\end{aretenir}

\coursec{Mise au point et utilisation pratique}

Nous savons désormais qu'une lunette astronomique afocale grossit les astres d'un facteur $G = f'_1/f'_2$ et que sa puissance s'exprime en dioptries. Mais entre la théorie et l'observation réelle, il reste un geste essentiel : la \cle{mise au point}. Une lunette mal réglée donne une image floue, fatigue l'œil et décourage l'observateur. Cette section t'apprend à régler correctement une lunette, à adapter ce réglage à la distance de l'objet et à ta propre vue, puis à choisir l'instrument d'optique le plus pertinent selon la situation.

\courssub{Le réglage initial : pointer et placer l'œil}

Avant toute mise au point, deux opérations préliminaires sont indispensables.

\textbf{Pointer la lunette vers l'objet.} Le champ d'une lunette astronomique est étroit : quelques degrés tout au plus. Il est donc illusoire de chercher un astre « au jugé ». On commence par viser l'objet à l'œil nu le long du tube, ou mieux à l'aide du \cle{chercheur} (petite lunette auxiliaire de faible grossissement et de grand champ, solidaire du tube principal). Une fois l'objet centré dans le chercheur, il l'est aussi dans la lunette.

\textbf{Placer l'œil au niveau du cercle oculaire.} Le faisceau qui sort de l'oculaire se rétrécit en un disque lumineux appelé \cle{cercle oculaire}, situé un peu en arrière de $F'_2$. C'est là que tout le flux lumineux est concentré : l'œil doit y placer sa pupille pour recevoir la totalité de la lumière et voir l'image la plus lumineuse et la plus large possible. Si l'œil est trop près ou trop loin, on observe des ombres qui envahissent le champ (le fameux « effet de trou de serrure »).

\begin{aretenir}[Cercle oculaire]
Le \cle{cercle oculaire} est l'image de la monture de l'objectif donnée par l'oculaire. C'est le plus petit cercle que traverse la lumière après l'oculaire. Pour une observation confortable, on place la pupille de l'œil exactement à cet endroit.
\end{aretenir}

\courssub{Mise au point pour un œil normal : le réglage afocal}

Rappelons l'exigence physiologique : un œil \cle{normal} (emmétrope) voit sans effort un objet placé à l'infini, c'est-à-dire \emph{sans accommoder}. Pour que l'observation soit confortable et non fatigante, l'image définitive donnée par la lunette doit donc être rejetée à l'infini.

Pour un objet à l'infini (étoile, Lune, bâtiment lointain), l'image intermédiaire $A_1B_1$ se forme dans le plan focal image de l'objectif : $A_1$ est confondu avec $F'_1$. Pour que l'oculaire renvoie cette image à l'infini, il faut que $A_1$ soit aussi confondu avec $F_2$, le foyer objet de l'oculaire. Le réglage correct est donc :

\begin{propriete}[Condition de mise au point afocale]
Pour un objet à l'infini et un œil normal, la mise au point est réalisée lorsque les foyers sont confondus : $F'_1 \equiv F_2$. La lunette est alors \cle{afocale} et le \cle{tirage} (distance objectif--oculaire) vaut :
\[ \overline{O_1O_2} = f'_1 + f'_2. \]
\end{propriete}

Concrètement, l'oculaire est monté sur une coulisse (ou commandé par une molette de mise au point) qui permet de faire varier le tirage de quelques centimètres autour de cette valeur de référence.

\begin{methode}[La mise au point en trois étapes]
\begin{enumerate}
\item \textbf{Viser l'objet} : pointer la lunette vers l'objet à l'aide du chercheur, puis placer l'œil derrière l'oculaire, au niveau du cercle oculaire.
\item \textbf{Tirer l'oculaire au maximum} (tirage le plus grand possible), puis le \textbf{rentrer progressivement} en tournant doucement la molette. Partir du tirage maximal garantit que l'image définitive arrive d'abord « de l'infini », ce qui interdit à l'œil d'accommoder.
\item \textbf{S'arrêter dès la première netteté}, l'œil parfaitement relâché, comme si l'on regardait l'horizon. Ne jamais « forcer » pour préciser davantage : la première position nette est la bonne.
\end{enumerate}
\end{methode}

\begin{attention}[Erreur fréquente : accommoder pendant la mise au point]
Si tu rentres l'oculaire trop loin, l'image définitive devient virtuelle à distance finie : ton œil peut encore la voir nette \emph{en accommodant}. La mise au point semble réussie, mais elle n'est valable que pour ton œil, à cet instant : un camarade emmétrope trouvera l'image floue, et ta vue fatiguera rapidement (picotements, maux de tête). En partant du tirage maximal et en s'arrêtant à la première netteté, on garantit une image à l'infini, donc une observation sans accommodation, confortable et transmissible.
\end{attention}

\courssub{Mise au point pour un objet rapproché : augmenter le tirage}

Que se passe-t-il si l'objet n'est plus à l'infini, par exemple un clocher à $\SI{200}{m}$ ? Appliquons la relation de conjugaison de Descartes à l'objectif, avec $\overline{O_1A} = -\SI{200}{m}$ et $f'_1 = \SI{0,50}{m}$ :
\[ \frac{1}{\overline{O_1A_1}} - \frac{1}{\overline{O_1A}} = \frac{1}{f'_1} \quad\Longrightarrow\quad \frac{1}{\overline{O_1A_1}} = \frac{1}{f'_1} + \frac{1}{\overline{O_1A}} = \frac{1}{0{,}50} + \frac{1}{-200} \;(\text{en m}^{-1}). \]
D'où $\overline{O_1A_1} = \dfrac{1}{2 - 0{,}005} \approx \SI{0,5013}{m} = \SI{50,13}{cm}$.

L'image intermédiaire se forme donc \emph{au-delà} de $F'_1$, à $1{,}3$ mm derrière le plan focal. Pour que l'oculaire continue de la rejeter à l'infini, il faut que $F_2$ vienne coïncider avec ce nouveau $A_1$ : on \textbf{éloigne l'oculaire de l'objectif}, c'est-à-dire qu'on \cle{augmente le tirage}.

\begin{propriete}[Mise au point sur un objet rapproché]
Plus l'objet observé est proche, plus l'image intermédiaire $A_1$ recule derrière $F'_1$, et plus il faut \textbf{augmenter le tirage} (sortir l'oculaire) pour conserver une image définitive à l'infini. C'est pourquoi les oculaires sont montés sur coulisse.
\end{propriete}

\begin{exemplebox}[Allongement du tirage pour un objet à 200 m]
Objectif de distance focale $f'_1 = \SI{50}{cm}$, objet à $D = \SI{200}{m}$.
\begin{itemize}
\item Objet à l'infini : tirage afocal $\overline{O_1O_2} = f'_1 + f'_2$.
\item Objet à $\SI{200}{m}$ : $\overline{O_1A_1} \approx \SI{50,13}{cm}$, soit un recul de l'image intermédiaire de $\SI{0,13}{cm} = \SI{1,3}{mm}$.
\end{itemize}
Il faut donc \textbf{allonger le tirage de seulement $\SI{1,3}{mm}$}. Remarque la sensibilité du réglage : une variation millimétrique suffit à passer de l'infini à $200$ m. C'est pourquoi la molette de mise au point doit être tournée avec douceur.
\end{exemplebox}

\courssub{Mise au point pour un œil myope}

Un œil \cle{myope} ne voit net sans accommodation que jusqu'à son \cle{punctum remotum} (PR), point le plus éloigné visible nettement, situé à distance finie (par exemple $\SI{50}{cm}$ devant l'œil). Pour qu'un observateur myope observe \emph{sans ses verres correcteurs} et sans accommoder, la lunette doit fournir une image définitive \emph{virtuelle} située exactement au PR de l'œil, et non plus à l'infini.

Pour cela, l'image intermédiaire $A_1$ doit se trouver \emph{légèrement à l'intérieur} du foyer objet de l'oculaire (entre $F_2$ et $O_2$) : l'oculaire fonctionne alors comme une loupe et donne une image virtuelle à distance finie. On y parvient en \textbf{réduisant légèrement le tirage} (on rentre l'oculaire) par rapport au réglage afocal.

\begin{aretenir}[Réglage selon l'œil]
\begin{itemize}
\item Œil normal (emmétrope) : tirage afocal, image définitive à l'infini.
\item Œil myope : tirage légèrement \emph{réduit}, image définitive virtuelle au punctum remotum.
\item Œil hypermétrope : tirage légèrement \emph{augmenté}, l'image intermédiaire étant placée légèrement au-delà de $F_2$ pour que l'œil compense en accommodant, comme il le fait déjà au quotidien.
\end{itemize}
Au Probatoire, on te demandera souvent de justifier le sens de déplacement de l'oculaire : raisonne toujours sur la position de $A_1$ par rapport à $F_2$.
\end{aretenir}

\begin{popfigure}
\begin{center}
\begin{tikzpicture}[scale=1,>=Stealth]
% Cas 1 : objet à l'infini
\draw[thick] (0,0) -- (0,1.4); \node[below,font=\small] at (0,0) {$O_1$};
\draw[thick,popblue] (2.6,0) -- (2.6,1.0); \node[below,font=\small,text=popblue] at (2.6,0) {$O_2$};
\draw[<->,popdark] (0,-0.7) -- (2.6,-0.7) node[midway,below,font=\small] {$f'_1+f'_2$};
\fill[poporange] (2.0,0) circle (1.2pt); \node[above,font=\small,text=poporange] at (2.0,0.05) {$F'_1{\equiv}F_2$};
\node[font=\small\bfseries,text=popdark] at (1.3,1.7) {Objet à l'infini};
% Cas 2 : objet rapproché
\begin{scope}[xshift=5.2cm]
\draw[thick] (0,0) -- (0,1.4); \node[below,font=\small] at (0,0) {$O_1$};
\draw[thick,popblue] (2.9,0) -- (2.9,1.0); \node[below,font=\small,text=popblue] at (2.9,0) {$O_2$};
\draw[<->,popdark] (0,-0.7) -- (2.9,-0.7) node[midway,below,font=\small] {tirage augmenté};
\fill[poporange] (2.0,0) circle (1.2pt); \node[above,font=\small,text=poporange] at (2.0,0.05) {$F'_1$};
\fill[popgreen] (2.9,0) circle (1.2pt); \node[above,font=\small,text=popgreen] at (2.9,1.05) {$F_2$};
\node[font=\small\bfseries,text=popdark] at (1.45,1.7) {Objet rapproché};
\end{scope}
% Cas 3 : œil myope
\begin{scope}[xshift=10.9cm]
\draw[thick] (0,0) -- (0,1.4); \node[below,font=\small] at (0,0) {$O_1$};
\draw[thick,popblue] (2.3,0) -- (2.3,1.0); \node[below,font=\small,text=popblue] at (2.3,0) {$O_2$};
\draw[<->,popdark] (0,-0.7) -- (2.3,-0.7) node[midway,below,font=\small] {tirage réduit};
\fill[poporange] (2.0,0) circle (1.2pt); \node[above,font=\small,text=poporange] at (2.0,0.05) {$F'_1$};
\fill[popgreen] (2.3,0) circle (1.2pt); \node[above,font=\small,text=popgreen] at (2.3,1.05) {$F_2$};
\node[font=\small\bfseries,text=popdark] at (1.15,1.7) {Œil myope};
\end{scope}
\end{tikzpicture}
\end{center}
\legende{Les trois réglages du tirage : position afocale pour un objet à l'infini (foyers confondus), tirage augmenté pour un objet rapproché ($A_1$ recule au-delà de $F'_1$), tirage réduit pour un œil myope ($A_1$ entre $F_2$ et $O_2$, image définitive virtuelle au punctum remotum).}
\end{popfigure}

\courssub{Vérification expérimentale du grossissement}

Comment vérifier concrètement le grossissement annoncé d'une lunette ? L'astuce classique consiste à observer \textbf{avec les deux yeux ouverts simultanément} : un œil regarde à travers la lunette, l'autre regarde directement à l'œil nu. Le cerveau superpose alors les deux images.

\begin{experience}[Mesure du grossissement sur un mur crénelé]
\textbf{Matériel} : une lunette astronomique sur pied, un bâtiment éloigné présentant un motif régulier (mur crénelé, rangées de fenêtres, briques).
\begin{enumerate}
\item Régler la lunette au point sur le bâtiment (méthode des trois étapes).
\item Ouvrir les deux yeux : l'œil derrière l'oculaire voit le motif grossi, l'autre œil voit le motif à taille réelle. Par entraînement, le cerveau superpose les deux images.
\item Compter combien de motifs vus à l'œil nu « tiennent » dans un seul motif vu dans la lunette. Ce nombre est le grossissement $G$.
\end{enumerate}
\textbf{Exemple} : si un créneau vu dans la lunette couvre exactement 8 créneaux vus à l'œil nu, alors $G = 8$. Pour une lunette avec $f'_1 = \SI{80}{cm}$ et $f'_2 = \SI{10}{cm}$, on retrouve bien $G = f'_1/f'_2 = 8$.
\end{experience}

\begin{exoresolu}[Réglage et grossissement d'une lunette de TP]
Une lunette de laboratoire est constituée d'un objectif de distance focale $f'_1 = \SI{60}{cm}$ et d'un oculaire de distance focale $f'_2 = \SI{5,0}{cm}$. Un élève observe d'abord un paysage très lointain, puis un panneau situé à $\SI{30}{m}$.
\begin{enumerate}
\item Quel est le tirage en réglage afocal ? Quel est le grossissement ?
\item De combien et dans quel sens faut-il modifier le tirage pour observer le panneau à $\SI{30}{m}$ ?
\end{enumerate}
\tcblower
\textbf{1.} En réglage afocal : $\overline{O_1O_2} = f'_1 + f'_2 = 60 + 5{,}0 = \SI{65}{cm}$. Le grossissement vaut :
\[ G = \frac{f'_1}{f'_2} = \frac{60}{5{,}0} = 12. \]
\textbf{2.} Pour le panneau à $D = \SI{30}{m}$, la relation de conjugaison appliquée à l'objectif donne :
\[ \overline{O_1A_1} = \frac{D\, f'_1}{D - f'_1} = \frac{30 \times 0{,}60}{30 - 0{,}60} = \frac{18}{29{,}4} \approx \SI{0,6122}{m} = \SI{61,22}{cm}. \]
L'image intermédiaire recule de $61{,}22 - 60 = \SI{1,22}{cm}$ derrière $F'_1$. Il faut donc \textbf{sortir l'oculaire de $\SI{1,2}{cm}$ environ} : le nouveau tirage vaut $\approx \SI{66,2}{cm}$. Le sens du déplacement (augmentation du tirage) est conforme à la règle : objet plus proche $\Rightarrow$ tirage plus grand.
\end{exoresolu}

\courssub{Choisir l'instrument adapté à la situation d'observation}

La lunette astronomique n'est pas l'instrument universel. Le choix d'un instrument d'optique se justifie par trois questions : \emph{à quelle distance est l'objet ? quelle taille angulaire ou quel détail veut-on voir ? accepte-t-on une image renversée ?}

\begin{itemize}
\item \textbf{Lunette astronomique} : fort grossissement ($G$ de 10 à plusieurs centaines), mais image \emph{renversée} et encombrement important. Réservée aux \textbf{objets célestes} (Lune, planètes, étoiles), pour lesquels le renversement de l'image n'a aucune importance.
\item \textbf{Jumelles} : grossissement modéré (typiquement $7\times$ à $12\times$), image \emph{droite} grâce à des prismes redresseurs, grand champ, légères et maniables. Idéales pour l'\textbf{observation terrestre} : paysages du mont Cameroun, faune du parc de Waza, match au stade.
\item \textbf{Loupe} : objet \emph{proche} placé au foyer d'une lentille convergente unique ; grossissement commercial modeste (2 à 10). Pour lire une carte, examiner un timbre, un insecte.
\item \textbf{Microscope} : objet \emph{extrêmement proche et minuscule} ; association objectif + oculaire à très courtes distances focales, grossissement de 40 à plus de 1000. Pour observer des cellules, des micro-organismes.
\end{itemize}

\begin{savaistu}[Les jumelles : deux lunettes astronomiques déguisées]
Une paire de jumelles, ce sont littéralement \textbf{deux lunettes astronomiques accolées}, une pour chaque œil ! Mais chaque tube contient en plus un jeu de \textbf{prismes} (prismes de Porro ou prismes en toit) qui remplissent deux fonctions : \emph{redresser l'image} (la rendre droite, indispensable pour l'observation terrestre) et \emph{replier le trajet de la lumière} en zigzag, ce qui raccourcit fortement l'encombrement : une lunette de tirage $50$ cm tient ainsi dans un tube de $15$ cm. Les inscriptions « $8\times 40$ » sur des jumelles signifient : grossissement $8$, diamètre d'objectif $40$ mm.
\end{savaistu}

\begin{popfigure}
\begin{center}
\begin{tikzpicture}[font=\small,>=Stealth,
sit/.style={draw=popdark,fill=popcream,rounded corners=2pt,align=center,minimum height=0.9cm,text width=3.1cm},
ins/.style={draw=popblue,fill=popblueL,rounded corners=2pt,align=center,minimum height=0.9cm,text width=3.1cm}]
\node[sit, align=center] (astre) at (0,3){Astre très lointain\\(Lune, planètes)};
\node[sit, align=center] (pays) at (4.6,3){Paysage, animal,\\match (loin, terrestre)};
\node[sit, align=center] (proche) at (0,0){Objet proche\\(texte, insecte)};
\node[sit, align=center] (micro) at (4.6,0){Objet microscopique\\(cellule, microbe)};
\node[ins, align=center] (l1) at (0,1.5){\textbf{Lunette astronomique}\\fort $G$, image renversée};
\node[ins, align=center] (l2) at (4.6,1.5){\textbf{Jumelles}\\$G$ modéré, image droite};
\node[ins, align=center] (l3) at (0,-1.5){\textbf{Loupe}\\objet au foyer, $G$ faible};
\node[ins, align=center] (l4) at (4.6,-1.5){\textbf{Microscope}\\très fort grossissement};
\draw[->,popdark] (astre) -- (l1);
\draw[->,popdark] (pays) -- (l2);
\draw[->,popdark] (proche) -- (l3);
\draw[->,popdark] (micro) -- (l4);
\end{tikzpicture}
\end{center}
\legende{Organigramme de choix d'un instrument d'optique selon la situation d'observation : la distance de l'objet et la nécessité d'une image droite guident le choix.}
\end{popfigure}

\begin{attention}[Pièges de rédaction au Probatoire]
\begin{itemize}
\item Ne jamais écrire que « la mise au point consiste à déplacer l'objectif » : c'est l'\textbf{oculaire} (ou le tirage) que l'on règle.
\item Pour un objet qui se rapproche, le tirage \textbf{augmente} ; beaucoup de candidats écrivent le contraire. Retiens : image intermédiaire qui recule $\Rightarrow$ oculaire qui recule.
\item Une lunette astronomique donne une image \textbf{renversée} : citer ce défaut est souvent demandé pour justifier l'usage des jumelles en observation terrestre.
\item Le grossissement $G = f'_1/f'_2$ n'a de sens que pour la lunette \textbf{afocale} (objet à l'infini, image à l'infini).
\end{itemize}
\end{attention}

\begin{exercice}[Pour t'entraîner]
Une lunette possède un objectif de $f'_1 = \SI{75}{cm}$ et un oculaire de $f'_2 = \SI{15}{mm}$.
\begin{enumerate}
\item Calculer le tirage afocal et le grossissement.
\item Un élève myope dont le punctum remotum est à $\SI{40}{cm}$ observe sans ses lunettes. Dans quel sens doit-il déplacer l'oculaire par rapport au réglage afocal ? Justifier.
\item Cet élève veut observer les éléphants du parc de Waza à $\SI{500}{m}$. Quel instrument conseilles-tu à la place de la lunette astronomique ? Donner deux arguments.
\end{enumerate}
\end{exercice}

\begin{corrige}[Exercice : lunette et choix d'instrument]
\textbf{1.} Tirage afocal : $\overline{O_1O_2} = f'_1 + f'_2 = 75 + 1{,}5 = \SI{76,5}{cm}$. Grossissement : $G = f'_1/f'_2 = 75/1{,}5 = 50$.
\textbf{2.} L'image définitive doit être virtuelle, située au PR ($\SI{40}{cm}$ devant l'œil), donc à distance finie et non à l'infini. Il faut que l'image intermédiaire $A_1$ se trouve entre $F_2$ et $O_2$ : l'élève doit \textbf{rentrer l'oculaire} (réduire légèrement le tirage).
\textbf{3.} Des \textbf{jumelles} : (i) elles donnent une image \emph{droite}, indispensable pour suivre des animaux ; (ii) leur champ est plus large et leur maniabilité bien meilleure qu'une lunette de grossissement 50, dont le champ très étroit rendrait le repérage des éléphants quasi impossible.
\end{corrige}

En résumé, la mise au point d'une lunette astronomique est un geste précis mais simple : viser, partir du tirage maximal, rentrer l'oculaire jusqu'à la première netteté, l'œil relâché. Ce réglage s'adapte à la distance de l'objet (tirage augmenté pour un objet proche) et à la vue de l'observateur (tirage réduit pour un myope). Enfin, le physicien averti choisit son instrument en fonction de la situation : lunette pour le ciel, jumelles pour la Terre, loupe et microscope pour le monde du proche et de l'infiniment petit.

\coursec{Bilan et choix raisonné d'un instrument d'optique}

Après avoir maîtrisé la mise au point de la lunette astronomique, il est temps de prendre du recul. Un physicien ne se contente pas de savoir \emph{comment} régler un instrument ; il doit comprendre \emph{pourquoi} il fonctionne ainsi, \emph{quand} l'utiliser et \emph{comment} le choisir parmi d'autres solutions optiques. Cette section synthétise la chaîne optique complète, établit le bilan des formules clés, t'offre une méthode rigoureuse pour justifier un choix instrumental et ouvre la porte vers les grands télescopes professionnels.

\courssub{Synthèse de la chaîne de fonctionnement : de l'objet à la rétine}

Rappelons le cheminement de la lumière dans une lunette astronomique réglée en \emph{vision normale} (image finale à l'infini), situation de référence pour l'observation astronomique confortable.

\begin{popfigure}
\begin{tikzpicture}[scale=0.95, >=Stealth]
% Styles
\tikzset{
    lens/.style={draw=popblue, very thick, fill=popblue!10},
    ray/.style={draw=poporange, thick, ->},
    axis/.style={draw=popdark, thin},
    img/.style={draw=poppurple, very thick, fill=poppurple!20},
    eye/.style={draw=poppink, thick, fill=poppink!20},
    label/.style={font=\small, text=popdark},
    focus/.style={fill=popgreen, circle, inner sep=1.5pt}
}

% Paramètres géométriques
\def\fOne{4.0}   % f'_1
\def\fTwo{3.0}   % f'_2
\def\L{\fOne+\fTwo} % 7.0
\def\h{1.2}      % Hauteur rayon incident (proportionnel à alpha)

% Axes optiques
\draw[axis] (-1,0) -- (11.5,0) node[right, label] {Axe optique};

% Objectif L1 (x=0)
\draw[lens] (0,-\fOne*0.4) -- (0,\fOne*0.4) node[midway, above=3mm, label] {Objectif $L_1$};
\draw[dashed, thin, popblue] (0,-\fOne*0.5) -- (0,\fOne*0.5);
\fill[focus] (0,0) node[below=4pt, label] {$O_1$};
\fill[focus] (-\fOne,0) node[below=4pt, label] {$F_1$};
\fill[focus] (\fOne,0) node[below=4pt, label] {$F'_1$};

% Oculaire L2 (x=L)
\draw[lens] (\L,-\fTwo*0.4) -- (\L,\fTwo*0.4) node[midway, above=3mm, label] {Oculaire $L_2$};
\draw[dashed, thin, popblue] (\L,-\fTwo*0.5) -- (\L,\fTwo*0.5);
\fill[focus] (\L,0) node[below=4pt, label] {$O_2$};
\fill[focus] (\L-\fTwo,0) node[below=4pt, label] {$F_2$};
\fill[focus] (\L+\fTwo,0) node[below=4pt, label] {$F'_2$};

% Distances
\draw[<->, thin, popmuted] (0, -2.0) -- (\L, -2.0) node[midway, below, label] {$L = f'_1 + f'_2$};
\draw[<->, thin, popmuted] (0, -2.0) -- (\fOne, -2.0) node[midway, below, label] {$f'_1$};
\draw[<->, thin, popmuted] (\L, -2.0) -- (\L+\fTwo, -2.0) node[midway, below, label] {$f'_2$};

% Faisceaux incidents (objet à l'infini, angle alpha)
% Haut (angle +alpha)
\draw[ray] (-1, \h) -- (0, \h);
\draw[ray] (0, \h) -- (\fOne, 0); % Vers F'_1
\draw[ray] (\fOne, 0) -- (\L, \h - \L*\h/\fOne); % Vers L2 (hauteur = h*(1 - L/f1) = h*(1 - (f1+f2)/f1) = -h*f2/f1)
\draw[ray] ({\L}, {\h - \L*\h/\fOne}) -- (11, {\h - \L*\h/\fOne - (11-\L)*\h/\fTwo}); % Sortant parallèle (pente -h/f2)

% Centre (angle 0)
\draw[ray] (-1, 0) -- (0, 0);
\draw[ray] (0, 0) -- (\fOne, 0);
\draw[ray] (\fOne, 0) -- (\L, 0);
\draw[ray] (\L, 0) -- (11, 0);

% Bas (angle -alpha)
\draw[ray] (-1, -\h) -- (0, -\h);
\draw[ray] (0, -\h) -- (\fOne, 0); % Vers F'_1
\draw[ray] (\fOne, 0) -- (\L, -\h + \L*\h/\fOne); % Vers L2
\draw[ray] ({\L}, {-\h + \L*\h/\fOne}) -- (11, {-\h + \L*\h/\fOne + (11-\L)*\h/\fTwo}); % Sortant parallèle

% Image intermédiaire I1 (réelle, renversée en F'_1)
\draw[img] (\fOne, -0.6) -- (\fOne, 0.6) node[midway, right=3mm, label] {Image réelle $A_1B_1$ (renversée)};
\node[label] at (\fOne, -0.9) {$F'_1 \equiv F_2$};

% Œil
\draw[eye] (11.5, -1) arc (180:360:1 and 0.5) -- (11.5, 1) arc (0:180:1 and 0.5) -- cycle;
\node[label] at (11.5, 0) {Œil};
\node[label, align=center] at (11.5, 1.5) {Image virtuelle\\ à l'infini\\ (droite par rapport à $I_1$)};

% Angles alpha et alpha'
% Alpha (entrée) : rayon haut passant par O1
\draw[popdark, ->] (0.3, 0.1) arc (0:16.7:0.35); % approx atan(h/f1) = atan(1.2/4)=16.7 deg
\node[label] at (-0.2, 0.5) {$\alpha$};
% Alpha' (sortie) : rayon haut sortant (pente négative)
\draw[popdark, ->] (\L+0.3, -0.1) arc (0:-22:0.35); % approx atan(h/f2) = atan(1.2/3)=21.8 deg
\node[label] at (\L+0.6, -0.5) {$\alpha'$};

% Légende intégrée
\node[draw=popline, fill=popcream, rounded corners, align=left, font=\small, anchor=north west] at (-0.8, 2.8) {
\textbf{Chaîne optique (vision normale) :}\\
1. Objet $\infty$ $\xrightarrow{L_1}$ Image réelle $I_1$ en $F'_1$ (renversée)\\
2. $I_1$ au foyer $F_2$ de $L_2$ $\xrightarrow{L_2}$ Faisceaux parallèles\\
3. Œil reçoit faisceaux $\parallel$ $\to$ Vision nette sans accommodation
};
\end{tikzpicture}
\legende{Chaîne de formation des images dans une lunette astronomique afocale (vision normale). L'objectif $L_1$ forme une image réelle renversée dans son plan focal image $F'_1$, qui coïncide avec le foyer objet $F_2$ de l'oculaire $L_2$. L'oculaire renvoie un faisceau émergent parallèle (image virtuelle à l'infini) vers l'œil.}
\end{popfigure}

\begin{aretenir}[Synthèse du fonctionnement]
La lunette astronomique est un système \textbf{afocal} : elle transforme un faisceau incident parallèle (objet à l'infini) en un faisceau émergent parallèle (image finale à l'infini).
\begin{enumerate}
    \item \textbf{Objectif $L_1$ (convergent, grande focale $f'_1$) :} Capte la lumière de l'objet lointain et forme une \textbf{image réelle, renversée et réduite} dans son plan focal image $F'_1$. C'est l'\textbf{image intermédiaire}.
    \item \textbf{Oculaire $L_2$ (convergent, petite focale $f'_2$) :} Placé de sorte que son foyer objet $F_2$ coïncide avec $F'_1$. L'image intermédiaire sert d'objet à l'oculaire. Comme cet objet est au foyer de $L_2$, l'oculaire émet un \textbf{faisceau émergent parallèle} : l'\textbf{image finale est virtuelle, à l'infini, droite par rapport à l'image intermédiaire} (donc renversée par rapport à l'objet initial).
    \item \textbf{Œil observateur :} Reçoit des faisceaux parallèles. Le cristallin converge ces rayons sur la rétine sans effort d'accommodation (vision de loin reposée).
\end{enumerate}
La distance entre les deux lentilles (longueur de la lunette) est $L = f'_1 + f'_2$.
\end{aretenir}

\courssub{Bilan des formules fondamentales}

Toutes les relations utiles découlent de la géométrie des faisceaux parallèles au voisinage de l'axe optique (approximation de Gauss, petits angles).

\begin{propriete}[Formules de la lunette astronomique afocale]
Pour une lunette astronomique constituée d'un objectif de focale image $f'_1$ et d'un oculaire de focale image $f'_2$, réglée en vision normale (afocale) :
\begin{itemize}
    \item \textbf{Longueur mécanique (distance lentilles) :} $\boxed{L = f'_1 + f'_2}$
    \item \textbf{Grossissement (puissance angulaire) :} $\boxed{G = \dfrac{\alpha'}{\alpha} = -\dfrac{f'_1}{f'_2}}$
    \item \textbf{Puissance commerciale (valeur positive) :} $\boxed{P = |G| = \dfrac{f'_1}{f'_2}}$
    \item \textbf{Puissance vergente de l'oculaire :} $\boxed{V_2 = \dfrac{1}{f'_2} = P \cdot \dfrac{1}{f'_1}}$ (souvent notée $P$ dans le commerce, en dioptries $\delta = \text{m}^{-1}$).
\end{itemize}
\textbf{Analyse dimensionnelle :} $f'_1$ et $f'_2$ sont des longueurs (m). $G$ et $P$ sont des nombres sans dimension. $L$ est une longueur (m). $V_2$ est en $\text{m}^{-1}$ (dioptrie).
\par\medskip
\textbf{Convention de signe :} $G < 0$ car l'image finale est renversée par rapport à l'objet ($\alpha'$ et $\alpha$ ont des signes contraires). En pratique, on cite la \textbf{puissance} $P = |G|$ (ex: « une lunette 50$\times$ »).
\end{propriete}

\begin{attention}[Grossissement vs Grandissement : piège classique !]
\begin{itemize}
    \item \textbf{Grandissement $\gamma$ (ou $g$) :} Rapport de la \textbf{taille de l'image} à la \textbf{taille de l'objet} ($\gamma = \frac{A'B'}{AB}$). Il s'applique pour un objet à \textbf{distance finie} (loupe, microscope, objectif photo).
    \item \textbf{Grossissement $G$ (ou $P$) :} Rapport de l'\textbf{angle sous-tendu par l'image} à l'\textbf{angle sous-tendu par l'objet} ($G = \frac{\alpha'}{\alpha}$). Il est \textbf{seul pertinent pour un objet à l'infini} (étoiles, Lune, cratères), car la « taille » de l'objet n'est pas définie (ou infinie).
\end{itemize}
\emph{Erreur fréquente au Probatoire :} Écrire $G = \frac{A'B'}{AB}$ pour une lunette astronomique. \textbf{Interdit.} On doit écrire $G = \frac{\alpha'}{\alpha} = -\frac{f'_1}{f'_2}$.
\end{attention}

\courssub{Critères de choix et comparaison des instruments d'optique}

Le choix d'un instrument ne se résume pas à son grossissement. Voici les critères physiques et pratiques à évaluer.

\begin{popfigure}
\begin{center}
\begin{tabularx}{\linewidth}{|l|X|X|X|}\hline
\rowcolor{popblue}
\textcolor{white}{\textbf{Instrument}} & \textcolor{white}{\textbf{Constitution typique}} & \textcolor{white}{\textbf{Grossissement / Grandissement}} & \textcolor{white}{\textbf{Usage principal \& Image}} \\ \hline
\textbf{Lunette astronomique} & Objectif $f'_1$ grand (500--1500 mm) + Oculaire $f'_2$ petit (5--25 mm) & $G = f'_1/f'_2$ (20$\times$--300$\times$) & Astronomie (Lune, planètes). Image \textbf{renversée}. Encombrant ($L \approx 1$ m). \\ \hline
\textbf{Jumelles (prismes Porro/Toit)} & 2 lunettes afocales couplées + prismes redressants & $G = f'_1/f'_2$ (7$\times$--12$\times$ typique) & Observation terrestre, ornithologie, marine. Image \textbf{droite}. Compact, stéréoscopique. \\ \hline
\textbf{Loupe (verre grossissant)} & 1 lentille convergente $f'$ court (10--50 mm) & $\gamma \approx \frac{25\text{ cm}}{f'}$ (2,5$\times$--25$\times$) & Lecture, philatélie, inspection près. Image \textbf{droite}, virtuelle, à 25 cm ou $\infty$. \\ \hline
\textbf{Microscope optique} & Objectif $f'_1$ très court (mm) + Oculaire $f'_2$ court (cm) & $\gamma \approx \frac{\Delta}{f'_1} \times \frac{25\text{ cm}}{f'_2}$ (100$\times$--1000$\times$) & Biologie, médecine (cellules, bactéries). Image \textbf{renversée}. Objet \textbf{proche} (mm). \\ \hline
\end{tabularx}
\end{center}
\legende{Tableau comparatif des principaux instruments d'optique à lentilles. Le choix dépend de la distance de l'objet (infini vs proche), de l'orientation de l'image souhaitée (droite vs renversée) et de l'encombrement.}
\end{popfigure}

\begin{methode}[Justifier le choix d'un instrument d'optique (3 étapes)]
Face à une situation concrète (sujet d'examen, projet personnel), structure ta réponse ainsi :
\begin{enumerate}
    \item \textbf{Analyser le besoin (Situation) :} Identifier la \textbf{distance de l'objet} (infini : astres, paysage lointain / finie : insecte, texte, cellule), l'\textbf{orientation de l'image} requise (droite pour observation terrestre, renversée acceptable pour astronomie), et la \textbf{transportabilité} souhaitée.
    \item \textbf{Identifier la grandeur pertinente :}
    \begin{itemize}
        \item Objet à l'infini $\to$ \textbf{Grossissement $G = f'_1/f'_2$} (ou Puissance $P$).
        \item Objet proche $\to$ \textbf{Grandissement $\gamma$} (loupe : $\gamma \approx 25\text{ cm}/f'$ ; microscope : $\gamma \approx \frac{\Delta}{f'_1}\frac{25\text{ cm}}{f'_2}$).
    \end{itemize}
    \item \textbf{Vérifier l'adéquation instrument/besoin :} Calculer la grandeur attendue avec les caractéristiques données (focales, $\Delta$). Comparer aux valeurs typiques du tableau. Mentionner la \textbf{luminosité} (diamètre de l'objectif/entrée de pupille) et l'\textbf{encombrement} ($L = f'_1+f'_2$).
\end{enumerate}
\emph{Exemple de phrase type :} « Pour observer les anneaux de Saturne (objet à l'infini, détails angulaires fins), une lunette astronomique est adaptée car elle fournit un grossissement $G = f'_1/f'_2$ élevé avec une image renversée (sans importance pour les astres). Une loupe serait inopérante (objet non proche) et des jumelles insuffisantes (grossissement limité $\approx 10\times$). »
\end{methode}

\courssub{Exemple chiffré : choix d'oculaires pour l'observation lunaire}

\begin{exemplebox}[Observation des cratères de la Lune]
\textbf{Situation :} Tu disposes d'une lunette astronomique dont l'objectif a une focale $f'_1 = 700\,\text{mm}$. Tu as deux oculaires interchangeables : $f'_{2a} = 25\,\text{mm}$ et $f'_{2b} = 10\,\text{mm}$. Le diamètre de la Lune sous-tend un angle d'environ $\alpha_{\text{Lune}} \approx 0,5^\circ \approx 8,7 \times 10^{-3}\,\text{rad}$. Le pouvoir de résolution de l'œil nu est d'environ $1' = 1/60^\circ \approx 2,9 \times 10^{-4}\,\text{rad}$. Quel oculaire choisir pour distinguer des cratères de taille angulaire $\approx 10''$ (secondes d'arc) sur la Lune ?

\textbf{Solution détaillée :}

1. \textbf{Calcul des grossissements :}
\[
G_a = \frac{f'_1}{f'_{2a}} = \frac{700}{25} = 28 \quad ; \quad G_b = \frac{f'_1}{f'_{2b}} = \frac{700}{10} = 70
\]

2. \textbf{Taille angulaire apparente de la Lune :}
\[
\alpha'_{a} = G_a \times \alpha_{\text{Lune}} = 28 \times 0,5^\circ = 14^\circ \quad ; \quad \alpha'_{b} = 70 \times 0,5^\circ = 35^\circ
\]
Les deux offrent un champ confortable (champ apparent oculaire typique $> 40^\circ$).

3. \textbf{Résolution des détails (cratères) :}
Un cratère de $10''$ (secondes d'arc) = $10/3600^\circ \approx 2,78 \times 10^{-3}{}^\circ \approx 4,85 \times 10^{-5}\,\text{rad}$.
Angle apparent vu dans la lunette :
\[
\alpha'_{\text{cratère}, a} = 28 \times 4,85 \times 10^{-5} \approx 1,36 \times 10^{-3}\,\text{rad} \approx 4,7'
\]
\[
\alpha'_{\text{cratère}, b} = 70 \times 4,85 \times 10^{-5} \approx 3,40 \times 10^{-3}\,\text{rad} \approx 11,6'
\]
L'œil nu résout $\approx 1'$. Les deux grossissements permettent \emph{théoriquement} de résoudre le cratère ($4,7' > 1'$ et $11,6' > 1'$).

4. \textbf{Contraintes pratiques (le choix final) :}
\begin{itemize}
    \item \textbf{Luminosité :} Pupille de sortie $D' = \frac{D_{\text{objectif}}}{G}$. Si $D_{\text{obj}} = 60\,\text{mm}$ : $D'_a = 2,1\,\text{mm}$, $D'_b = 0,86\,\text{mm}$. L'oculaire $10\,\text{mm}$ donne une image plus sombre (pupille de sortie $<$ pupille œil $\approx 2-3\,\text{mm}$ la nuit).
    \item \textbf{Champ de vue :} Plus $G$ est grand, plus le champ réel est petit. $G=70$ montre une toute petite portion de Lune.
    \item \textbf{Turbulence atmosphérique (seeing) :} Au Cameroun, la turbulence limite souvent le grossissement utile à $\approx 1,5 \times D_{\text{obj}}(\text{mm}) \approx 90\times$. $70\times$ est à la limite, $28\times$ est plus stable.
\end{itemize}
\textbf{Conclusion :} L'oculaire \textbf{$25\,\text{mm}$ ($G=28$)} est préférable pour une vue d'ensemble nette, lumineuse et stable de la Lune. L'oculaire $10\,\text{mm}$ ($G=70$) est réservée aux nuits exceptionnellement stables pour un détail précis sur une petite zone.
\end{exemplebox}

\courssub{Ouverture culturelle : des lentilles aux miroirs — les grands télescopes}

\begin{savaistu}[Pourquoi les grands télescopes professionnels sont-ils des miroirs, pas des lentilles ?]
La lunette astronomique (réfracteur) a dominé l'astronomie jusqu'au XIXe siècle (ex: lunette de Meudon 83 cm, Yerkes 102 cm). Mais elle bute sur des limites physiques infranchissables pour les géants actuels (VLT 8,4 m, ELT 39 m, James Webb 6,5 m) :
\begin{enumerate}
    \item \textbf{Aberration chromatique :} Une lentille simple dévie les couleurs différemment. Corriger cela demande des doublets/triplets complexes, coûteux et lourds.
    \item \textbf{Déformation par gravité :} Une lentille ne peut être soutenue que par son bord (son centre doit être libre pour laisser passer la lumière). Un miroir de 10 m pèse des tonnes ; s'il fléchit, l'image floute. Un miroir, lui, peut être soutenu par tout son dos (système actif/opto-mécanique).
    \item \textbf{Masse et volume :} Le verre optique de qualité (ex: FK51, SF6) est dense et cher. Un miroir n'a besoin de qualité optique que sur sa surface réfléchissante (quelques $\mu$m d'aluminium ou argent sur substrat léger type ZERODUR ou nid d'abeille).
    \item \textbf{Longueur focale :} Un réfracteur de 10 m de diamètre aurait une longueur $L \approx f'_1 \sim 15-20\,\text{m}$ (tube énorme, instable au vent). Un télescope Ritchey-Chrétien (deux miroirs) plie le chemin optique : tube court ($\approx 10-12\,\text{m}$ pour 8 m de diamètre), stable, compact.
\end{enumerate}
\textbf{Au Cameroun :} L'Observatoire astronomique de l'Université de Yaoundé I et les clubs d'astronomie (ex: \emph{Association Camerounaise d'Astronomie}) utilisent souvent des télescopes Newton (miroir primaire + miroir plan secondaire) ou Schmidt-Cassegrain (miroirs + lame de Schmidt) pour leur compacité et leur rapport performance/prix. Le principe reste celui de la lunette : \textbf{former une image réelle au foyer d'un objectif (miroir ici), puis la grossir par un oculaire}. La formule $G = f'_{\text{objectif}} / f'_{\text{oculaire}}$ reste valide !
\end{savaistu}

\begin{exoresolu}[Choix d'instrument pour un naturaliste au Parc de la Waza]
\textbf{Énoncé :} Un naturaliste souhaite observer des oiseaux perchés à $\approx 50\,\text{m}$ (détails plumage) et repérer des troupeaux d'éléphants à $\approx 2\,\text{km}$ dans la savane. Il hésite entre : (A) Lunette astronomique $f'_1=1000\,\text{mm}, f'_2=20\,\text{mm}$ ; (B) Jumelles $8\times42$ ; (C) Loupe $f'=50\,\text{mm}$. Justifier le meilleur choix.

\textbf{Solution détaillée :}
\tcblower
\textbf{Analyse des besoins :}
\begin{itemize}
    \item Objets à distance \textbf{finie} (50 m, 2 km) $\to$ notion de \textbf{grandissement $\gamma$} ou grossissement angulaire $G$ pertinent.
    \item Observation \textbf{terrestre} $\to$ image \textbf{droite} impérative.
    \item Mobilité (marche en brousse) $\to$ \textbf{compacité, légèreté, robustesse}.
    \item Vision \textbf{binoculaire} (deux yeux) $\to$ confort, perception relief, champ large.
\end{itemize}

\textbf{Évaluation des options :}
\begin{itemize}
    \item \textbf{(A) Lunette astronomique :} $G = 50\times$. Image \textbf{renversée} (inadaptée terrain). Longueur $L=1,02\,\text{m}$ (encombrant). Monoculaire (fatigue). \textbf{Rejetée.}
    \item \textbf{(C) Loupe :} $\gamma \approx 250/50 = 5\times$ (vision à 25 cm). Objet doit être \emph{très proche} (distance focale 5 cm). Inutilisable pour oiseaux à 50 m. \textbf{Rejetée.}
    \item \textbf{(B) Jumelles $8\times42$ :} Grossissement $G=8\times$ (standard terrain). Diamètre objectif $42\,\text{mm}$ $\to$ luminosité excellente (pupille sortie $5,25\,\text{mm}$). Prismes redressants $\to$ image \textbf{droite}. Compacts, légers, étanches, vision binoculaire. Champ large ($\approx 120\,\text{m}$ à $1000\,\text{m}$).
    \item \textbf{Vérification grossissement :} À $50\,\text{m}$, $G=8$ donne taille apparente $\times 8$ (ex: oiseau 20 cm $\to$ 1,6 m apparent). À $2\,\text{km}$, éléphant 3 m $\to$ 24 m apparent. Suffisant pour identification.
\end{itemize}
\textbf{Choix :} \textbf{Les jumelles $8\times42$ (B)}. C'est l'instrument \emph{par excellence} de l'ornithologue et du naturaliste de terrain.
\end{exoresolu}

\begin{aretenir}[Bilan final : Choisir son instrument]
\begin{center}
\begin{tabular}{|c|c|c|c|}\hline
\textbf{Question clé} & \textbf{Réponse} & \textbf{Instrument adapté} & \textbf{Grandeur clé} \\ \hline
Objet à l'infini (astres) ? & Oui, image renversée OK & Lunette astronomique / Télescope & $G = f'_1/f'_2$ \\ \hline
Objet à l'infini (terrestre) ? & Oui, image \textbf{droite} requise & \textbf{Jumelles} (prismes) & $G = f'_1/f'_2$ \\ \hline
Objet proche (lecture, insecte) ? & Oui, main libre, image droite & \textbf{Loupe} & $\gamma \approx 25\,\text{cm}/f'$ \\ \hline
Objet très proche (cellule, $\mu$m) ? & Oui, grossissement fort & \textbf{Microscope} & $\gamma \approx \frac{\Delta}{f'_1}\frac{25\,\text{cm}}{f'_2}$ \\ \hline
\end{tabular}
\end{center}
\textbf{Règle d'or :} \emph{La distance de l'objet dicte l'instrument. Le grossissement se calcule ensuite. La luminosité (diamètre entrée) et l'ergonomie (image droite, binoculaire, poids) valident le choix final.}
\end{aretenir}

\courssub{Exercices d'entraînement (Probatoire type)}

\begin{exercice}[Exercice 1 : Caractéristiques d'une lunette]
Une lunette astronomique a une longueur $L = 1,20\,\text{m}$ et un grossissement $G = -40$.
\begin{enumerate}
    \item Calculer les focales $f'_1$ (objectif) et $f'_2$ (oculaire).
    \item L'observateur remplace l'oculaire par un modèle de focale $f''_2 = 10\,\text{mm}$. Quelle est la nouvelle longueur $L'$ à régler pour la vision normale ? Quel est le nouveau grossissement $G'$ ?
    \item Pourquoi l'image fournie par cette lunette est-elle renversée ? Répondre par un argument géométrique (marche des rayons).
\end{enumerate}
\end{exercice}

\begin{exercice}[Exercice 2 : Choix justifié — Surveillance volcanique]
Un volcanologue au Mont Cameroun doit surveiller l'activité du cratère (distance $\approx 5\,\text{km}$) de jour comme de nuit. Il dispose de trois instruments :
\begin{itemize}
    \item Instrument 1 : Lunette astronomique $f'_1 = 800\,\text{mm}$, $f'_2 = 16\,\text{mm}$, $D_{\text{obj}} = 50\,\text{mm}$.
    \item Instrument 2 : Jumelles $10 \times 50$.
    \item Instrument 3 : Télescope Newton (miroir) $D = 114\,\text{mm}$, $f' = 900\,\text{mm}$, oculaires $25\,\text{mm}$ et $10\,\text{mm}$.
\end{itemize}
Justifier par un tableau comparatif (grossissement, luminosité, orientation image, encombrement, usage nuit/jour) l'instrument le plus polyvalent pour cette mission.
\end{exercice}

\begin{corrige}[Exercice 1]
\begin{enumerate}
    \item Système : $\begin{cases} f'_1 + f'_2 = 1200\,\text{mm} \\ f'_1/f'_2 = 40 \end{cases} \Rightarrow f'_2 = 1200/41 \approx 29,3\,\text{mm}$, $f'_1 \approx 1170,7\,\text{mm}$.
    \item $L' = f'_1 + f''_2 = 1170,7 + 10 = 1180,7\,\text{mm} \approx 1,18\,\text{m}$. $G' = -1170,7/10 \approx -117$.
    \item L'objectif forme une image réelle renversée en $F'_1$. L'oculaire, utilisé comme loupe sur cette image renversée, produit une image virtuelle à l'infini qui conserve l'orientation de son objet (l'image intermédiaire). Donc image finale renversée par rapport à l'objet initial.
\end{enumerate}
\end{corrige}

\begin{corrige}[Exercice 2]
\begin{center}
\begin{tabularx}{\linewidth}{|l|X|X|X|}\hline
\rowcolor{popblue}
\textcolor{white}{\textbf{Critère}} & \textcolor{white}{\textbf{Inst. 1 (Lunette)}} & \textcolor{white}{\textbf{Inst. 2 (Jumelles)}} & \textcolor{white}{\textbf{Inst. 3 (Newton)}} \\ \hline
Grossissement & $G=50\times$ & $G=10\times$ & $G_{25}=36\times$, $G_{10}=90\times$ \\ \hline
Pupille sortie (luminosité) & $50/50 = 1,0\,\text{mm}$ (faible nuit) & $50/10 = 5,0\,\text{mm}$ (excellente) & $114/36=3,2\,\text{mm}$ ; $114/90=1,3\,\text{mm}$ \\ \hline
Orientation image & Renversée & \textbf{Droite} & Renversée \\ \hline
Encombrement & $\approx 1,2\,\text{m}$ tube & Compact ($\approx 18\,\text{cm}$) & $\approx 90\,\text{cm}$ tube \\ \hline
Vision & Monoculaire & \textbf{Binoculaire} & Monoculaire \\ \hline
Polyvalence Jour/Nuit & Mauvais (sombre, renverse) & \textbf{Excellent} & Bon nuit, mauvais jour (renverse) \\ \hline
\end{tabularx}
\end{center}
\textbf{Choix : Instrument 2 (Jumelles $10\times50$).} Seule option image droite, binoculaire, compacte, lumineuse jour/nuit. Grossissement $10\times$ suffisant pour détails cratère à 5 km (taille angulaire $\sim$ minutes d'arc).
\end{corrige}

\newpage

\coursec{Activité d'intégration}

\courssub{Objectifs de l'activité}

À l'issue de cette activité d'intégration, tu seras capable de :
\begin{itemize}
\item \cle{décrire} une lunette astronomique et justifier le choix de ses deux lentilles (objectif et oculaire) en fonction du besoin d'observation ;
\item \cle{calculer} la longueur d'une lunette en configuration afocale et vérifier la faisabilité technique d'un instrument ;
\item \cle{calculer} le grossissement d'une lunette et l'angle sous lequel un astre est vu à travers l'instrument ;
\item \cle{tracer} à l'échelle la marche d'un faisceau lumineux issu d'un point de l'objet à travers les deux lentilles ;
\item \cle{comparer} de manière argumentée deux instruments d'optique (lunette et jumelles) et choisir le plus pertinent selon la situation d'observation.
\end{itemize}

Cette activité mobilise l'ensemble des savoir-faire de la leçon dans une situation authentique de conception : tu n'es plus seulement un utilisateur de la lunette, tu deviens son \cle{concepteur}.

\courssub{Situation}

Le club scientifique du Lycée de Ngaoundéré prépare la \og Nuit des étoiles \fg{}, une soirée d'observation ouverte à tous les élèves et aux parents d'élèves. Le proviseur, enthousiaste, met à la disposition du club le matériel suivant, récupéré dans la réserve du laboratoire de physique :

\begin{itemize}
\item trois lentilles minces convergentes de distances focales connues : $f_1 = \SI{5}{cm}$, $f_2 = \SI{20}{cm}$ et $f_3 = \SI{100}{cm}$ ;
\item un tube en PVC de diamètre intérieur \SI{6}{cm}, dont la longueur est réglable par coulissement entre \SI{20}{cm} et \SI{150}{cm} ;
\item deux supports de lentilles coulissants et un trépied photo.
\end{itemize}

Le proviseur formule deux souhaits précis :
\begin{enumerate}
\item observer la \cle{Lune}, dont le diamètre apparent vu à l'œil nu est $\theta = \SI{0,5}{\degree}$, avec un grossissement aussi grand que possible, afin de distinguer les détails de son relief ;
\item suivre, depuis la colline qui surplombe la ville, la couverture d'un match de football au stade municipal, situé à environ \SI{2}{km}.
\end{enumerate}

Le budget du club est nul : il faut donc \cle{concevoir} la lunette avec le matériel disponible, en choisissant judicieusement les lentilles, en calculant ses caractéristiques (longueur, grossissement), puis en argumentant le choix de l'instrument le mieux adapté à chacune des deux observations.

Par groupes de quatre, vous êtes les ingénieurs du club. Votre production sera restituée en trois minutes au tableau, schéma à l'appui, devant la classe qui jouera le rôle du proviseur.

\courssub{Consignes}

\begin{enumerate}
\item \textbf{Choix des lentilles.} Rappeler les rôles respectifs de l'objectif et de l'oculaire dans une lunette astronomique. Parmi les trois lentilles disponibles, choisir celle qui servira d'objectif et celle qui servira d'oculaire, en justifiant chaque choix par la condition sur les distances focales. Que devient la troisième lentille ?

\item \textbf{Longueur de la lunette afocale.} Donner l'expression de la distance entre les centres optiques $O_1$ et $O_2$ des deux lentilles en configuration afocale. Calculer la longueur du tube nécessaire et vérifier qu'elle est compatible avec le tube réglable du club ($\SI{20}{cm} \leq L \leq \SI{150}{cm}$).

\item \textbf{Grossissement et angle d'observation.} Calculer le grossissement $G$ de la lunette ainsi conçue. En déduire l'angle $\theta'$ sous lequel la Lune sera vue à travers l'instrument. Comparer avec le diamètre apparent de la Lune à l'œil nu et conclure sur l'intérêt de l'instrument.

\item \textbf{Tracé de la marche d'un faisceau lumineux.} Sur une feuille, à l'échelle $1/10$ (1 cm sur la feuille pour 10 cm réels), placer les deux lentilles et leurs foyers. Tracer la marche à travers la lunette d'un faisceau parallèle issu du bord supérieur du disque lunaire, faisant un angle $\theta$ avec l'axe optique. Repérer l'image intermédiaire et faire apparaître l'angle $\theta'$ du faisceau émergent. Que constate-t-on quant au sens de l'image ?

\item \textbf{Argumentaire de choix d'instrument.} Rédiger un court texte (une dizaine de lignes) destiné au proviseur, expliquant :
\begin{itemize}
\item pourquoi la lunette conçue convient parfaitement à l'observation de la Lune ;
\item pourquoi des jumelles seraient préférables pour suivre le match au stade (penser au sens de l'image, au champ de vision et à la maniabilité).
\end{itemize}

\item \textbf{Restitution.} Préparer une présentation orale de trois minutes : schéma au tableau, résultats chiffrés annoncés avec leurs unités, et conclusion argumentée.
\end{enumerate}

\courssub{Ressources à mobiliser}

\begin{aretenir}[Boîte à outils de la leçon]
\begin{itemize}
\item \cle{Constitution} : une lunette astronomique est formée d'un \textbf{objectif} (lentille convergente de grande distance focale $f'_1$, côté objet) et d'un \textbf{oculaire} (lentille convergente de courte distance focale $f'_2$, côté œil).
\item \cle{Lunette afocale} : le foyer image $F'_1$ de l'objectif coïncide avec le foyer objet $F_2$ de l'oculaire. Un faisceau parallèle incident donne un faisceau parallèle émergent : l'image finale est rejetée à l'infini et l'œil observe sans accommoder. La longueur de la lunette est alors :
\[ L = O_1O_2 = f'_1 + f'_2 \]
\item \cle{Grossissement} : $G = \dfrac{\theta'}{\theta} = \dfrac{f'_1}{f'_2}$, où $\theta$ est l'angle sous lequel l'objet est vu à l'œil nu et $\theta'$ l'angle sous lequel l'image est vue à travers la lunette. $G$ est sans unité.
\item \cle{Sens de l'image} : la lunette astronomique donne d'un objet une image \textbf{renversée} (inversée).
\item \cle{Tracé} : un rayon passant par le centre optique d'une lentille n'est pas dévié ; un rayon parallèle à l'axe émerge en passant par le foyer image ; un faisceau parallèle incliné converge dans le plan focal image.
\item \cle{Analyse dimensionnelle} : $f'_1$ et $f'_2$ sont des longueurs (m), donc $L$ est une longueur (m) et $G$ est adimensionnel. Cohérence vérifiée.
\end{itemize}
\end{aretenir}

\courssub{Démarche et corrigé commenté}

\begin{exoresolu}[Conception de la lunette du club d'astronomie]
\textbf{Énoncé.} Le club dispose de trois lentilles convergentes de distances focales $f_1 = \SI{5}{cm}$, $f_2 = \SI{20}{cm}$, $f_3 = \SI{100}{cm}$ et d'un tube réglable de \SI{20}{cm} à \SI{150}{cm}. On veut observer la Lune ($\theta = \SI{0,5}{\degree}$) avec un grossissement maximal. (1) Choisir objectif et oculaire. (2) Calculer la longueur afocale. (3) Calculer le grossissement et l'angle $\theta'$. (4) Tracer le faisceau à l'échelle $1/10$. (5) Argumenter le choix lunette/jumelles.
\tcblower
\textbf{Solution détaillée.}

\textbf{1. Choix des lentilles.} Dans une lunette astronomique, l'objectif doit avoir une \emph{grande} distance focale (pour former une image intermédiaire grande et lumineuse) et l'oculaire une \emph{courte} distance focale (il joue le rôle de loupe sur l'image intermédiaire). De plus, le grossissement $G = f'_1/f'_2$ est d'autant plus grand que $f'_1$ est grand et $f'_2$ petit. On choisit donc :
\begin{itemize}
\item \textbf{objectif} : la lentille de $f'_1 = f_3 = \SI{100}{cm}$ ;
\item \textbf{oculaire} : la lentille de $f'_2 = f_1 = \SI{5}{cm}$.
\end{itemize}
La lentille de \SI{20}{cm} n'est pas utilisée : en oculaire, elle donnerait un grossissement cinq fois plus faible ($G=5$) ; en objectif, un grossissement quatre fois plus faible ($G=4$).

\textbf{2. Longueur de la lunette afocale.} En configuration afocale, $F'_1$ et $F_2$ sont confondus, donc :
\[ L = f'_1 + f'_2 = \SI{100}{cm} + \SI{5}{cm} = \SI{105}{cm} \]
Cette longueur est bien comprise entre \SI{20}{cm} et \SI{150}{cm} : le tube réglable du club convient, réglé sur \SI{105}{cm}. Le projet est techniquement réaliste.

\textbf{3. Grossissement et angle d'observation.}
\[ G = \frac{f'_1}{f'_2} = \frac{100}{5} = 20 \]
L'angle sous lequel la Lune est vue à travers la lunette est :
\[ \theta' = G \times \theta = 20 \times \SI{0,5}{\degree} = \SI{10}{\degree} \]
La Lune apparaît sous un angle vingt fois plus grand qu'à l'œil nu : son diamètre apparent passe de $\SI{0,5}{\degree}$ à $\SI{10}{\degree}$, ce qui permet de distinguer les cratères et les mers lunaires. L'objectif du proviseur est atteint.

\textbf{4. Tracé de la marche d'un faisceau lumineux (échelle 1/10).}
\begin{popfigure}
\begin{tikzpicture}[scale=1, >=Stealth]
% Échelle : 1 cm TikZ = 1 cm sur la feuille = 10 cm réel
% Objectif O1 à x=0, Oculaire O2 à x=10.5
% Foyers : F'1 (obj) à x=10, F2 (ocul) à x=10 (confondus)
% F'2 (ocul) à x=11
% Rayon incident parallèle à l'axe (haut) -> passe par F'1
% Rayon incident par centre O1 -> non dévié
% Intersection en B1 (image intermédiaire) à x=10, y = -f1*tan(theta) ~ -10*0.087 = -0.87 (exagéré pour visibilité)
% Pour le schéma, on prend theta ~ 10 deg pour voir le tracé.
\def\focObj{10} % cm sur feuille
\def\focOcu{0.5} % cm sur feuille
\def\Ltube{10.5} % cm sur feuille
\def\angleInc{12} % degrés, exagéré pour le dessin
\def\hauteurRayon{2} % hauteur d'incidence sur l'objectif (arbitraire pour le dessin)

% Lentilles
\draw[thick, popblue] (0, -1.5) -- (0, 1.5) node[above left] {$L_1$ (Objectif)};
\draw[thick, poporange] (\Ltube, -1) -- (\Ltube, 1) node[above right] {$L_2$ (Oculaire)};

% Foyers
\draw[dashed, popmuted] (\focObj, -1.5) -- (\focObj, 1.5) node[above] {$F'_1 = F_2$};
\draw[dashed, popmuted] (\Ltube+\focOcu, -1) -- (\Ltube+\focOcu, 1) node[above] {$F'_2$};
\draw[dashed, popmuted] (-\focObj, -1.5) -- (-\focObj, 1.5) node[above] {$F_1$};

% Axe optique
\draw[popline, <->] (-1, 0) -- (\Ltube+\focOcu+1, 0) node[right] {Axe optique};

% Rayons incidents (venant de l'infini, angle theta)
% Rayon 1 : passant par O1 (centre objectif)
\draw[popgreen, thick, ->] (-1, {-\hauteurRayon}) -- (0,0) -- (\Ltube, {-\hauteurRayon * \Ltube / 1}); % Prolongé vers oculaire
% Correction : le rayon passant par O1 sort avec même pente.
% Pente = -hauteurRayon / 1 (depuis x=-1). A x=0, y=0. A x=Ltube, y = -hauteurRayon * Ltube.
% Mais on veut qu'il passe par B1 (image intermédiaire).
% Image intermédiaire B1 est au foyer F'1 (x=10). Hauteur y_B1 = -f'_1 * tan(theta).
% tan(12deg) ~ 0.21. y_B1 = -10 * 0.21 = -2.1.
% Rayon par O1 : y = x * tan(-theta) = -x * 0.21. A x=10, y=-2.1. OK.
% Rayon parallèle à l'axe (hauteur y=2) : converge vers F'1 (10, 0).
% Rayon incident parallèle incliné (theta) : hauteur y=2 à x=0. Converge vers point focal dans plan focal : (10, -2.1).

% Définition des points pour le tracé principal (bord supérieur de la Lune -> angle theta)
% On trace 2 rayons caractéristiques pour le bord supérieur (angle +theta vers le haut? Non, bord supérieur -> rayon descendant vers l'axe).
% Objet en haut -> rayon arrive en bas. Angle theta = 0.5 deg. 
% Pour le dessin, on prend un faisceau parallèle incliné vers le BAS (bord supérieur de la lune).
% Rayon 1 : Par centre O1 (0,0). Pente -tan(theta).
% Rayon 2 : Parallèle à l'axe optique incident à hauteur h (ex: 1.5cm). Passe par F'1 (10, 0).
% Rayon 3 : Parallèle à l'axe optique incident à hauteur 0 (axe). Passe par F'1 (10, 0).
% Intersection en B1 (10, -h_B1).

\pgfmathsetmacro{\tanTheta}{tan(\angleInc)}
\pgfmathsetmacro{\yB}{-\focObj * \tanTheta} % Hauteur image intermédiaire (négative = bas)

% Rayon incident 1 : Par centre O1
\draw[popgreen, thick, ->] (-1, {\tanTheta}) -- (0,0) -- (\Ltube, {\tanTheta * \Ltube}) -- (\Ltube+2, {\tanTheta * (\Ltube+2)});
% Rayon incident 2 : Parallèle à l'axe, hauteur y=1.5 (passe par F'1)
\draw[popgreen, thick, ->] (-1, 1.5) -- (0, 1.5) -- (\focObj, 0) -- (\Ltube, {\yB * (\Ltube - \focObj) / (-\focObj)}); % Ligne de (10,0) à B1(10, yB) ? Non.
% Rayon parallèle à l'axe (hauteur 1.5) -> réfracté vers F'1 (10, 0).
% Ensuite de F'1 vers oculaire. Comme F'1 = F2, le rayon sort parallèle à l'axe ? Non.
% Si le rayon passe par F2 (foyer objet oculaire), il sort PARALLÈLE à l'axe.
% Mais ici le rayon passe par F'1 (qui est F2). Donc il sort PARALLÈLE à l'axe.
% Mais l'image est renversée. Le faisceau émergent doit être incliné de -theta' (vers le haut si incident vers le bas).
% Attendez : Bord supérieur de la Lune -> rayons arrivent en descendant (angle -theta par rapport à l'axe).
% Image intermédiaire en BAS (y négatif).
% Oculaire : objet (image intermédiaire) au foyer F2. 
% Rayon partant de B1 (bas) et passant par centre O2 -> sort non dévié (vers le bas).
% Rayon partant de B1 (bas) et parallèle à l'axe -> passe par F'2 (haut? non, lentille convergente, rayon parallèle -> F'2).
% Donc faisceau émergent est PARALLÈLE, incliné vers le HAUT (angle +theta').
% Image finale à l'infini, renversée (haut en bas).

% Reprenons le tracé proprement pour le bord supérieur (rayons descendants).
% Rayon A : Passe par O1 (0,0). Direction descendante (angle -alpha).
\draw[popgreen, very thick, ->] (-1.5, {1.5*\tanTheta}) -- (0,0) -- (\Ltube, {-1.5*\tanTheta}) -- (\Ltube+2, {-3.5*\tanTheta});
\node[popgreen, anchor=west] at (-1.5, {1.5*\tanTheta}) {Bord sup. Lune};

% Rayon B : Parallèle à l'axe, hauteur y = 1.5 (au-dessus de l'axe). 
% Il converge vers F'1 (10, 0).
\draw[popgreen, very thick, ->] (-1.5, 1.5) -- (0, 1.5) -- (\focObj, 0);
% De F'1 (qui est F2) vers oculaire : comme il passe par F2, il sort PARALLÈLE à l'axe.
\draw[popgreen, very thick, ->] (\focObj, 0) -- (\Ltube, 0) -- (\Ltube+2, 0);

% Rayon C : Parallèle au rayon A (donc même angle -alpha), passant par le haut de l'objectif (y=1.5 à x=0? Non, hauteur arbitraire).
% Prenons un rayon incident parallèle à A, passant par le haut de l'objectif (0, 1.5).
% Il converge dans le plan focal en B1 (10, yB).
\pgfmathsetmacro{\yInc}{1.5}
\pgfmathsetmacro{\yBcalc}{\yInc + \focObj * (-\tanTheta)} % y = y0 + f * tan(theta_ray). Theta_ray = -alpha.
% Si rayon incident a pente -tanTheta. A x=0, y=1.5. A x=10, y = 1.5 - 10*tanTheta.
\draw[popgreen, very thick, ->] (-1.5, {\yInc + 1.5*\tanTheta}) -- (0, \yInc) -- (\focObj, {\yInc - \focObj*\tanTheta});
% Ce point est B1.
\pgfmathsetmacro{\yBfinal}{\yInc - \focObj*\tanTheta}
\fill[popgreen] (\focObj, \yBfinal) circle (2pt) node[right] {$B_1$ (Image intermédiaire)};
% De B1 vers oculaire. 
% Rayon de B1 vers centre O2 (10.5, 0). Pente = (0 - yB) / (10.5 - 10) = -yB / 0.5.
% Comme yB est négatif, pente positive. Sort vers le HAUT.
\draw[popgreen, very thick, ->] (\focObj, \yBfinal) -- (\Ltube, 0) -- (\Ltube+2, {-2 * \yBfinal}); % Prolongé
% Rayon de B1 parallèle à l'axe (horizontal). Passe par F'2 (11, 0).
\draw[popgreen, very thick, ->] (\focObj, \yBfinal) -- (\Ltube, \yBfinal) -- (\Ltube+\focOcu, 0) -- (\Ltube+2, {-2 * \yBfinal}); % Pente après F'2 : (0 - yB)/0.5 = -2*yB. Même pente que l'autre. Bien.

% Faisceau émergent : deux rayons parallèles (sortant de O2 et de F'2) avec même pente positive.
% Angle theta' = -yB / f'_2 = (f'_1 * tanTheta) / f'_2 = G * tanTheta.
% Flèche angle theta'
\draw[poppurple, <->] (\Ltube+0.2, 0) arc (0:{\angleInc*20/5}:0.5) node[midway, right] {$\theta'$};
% Flèche angle theta (incident)
\draw[poppurple, <->] (-0.5, 0) arc (0:{-\angleInc}:0.5) node[midway, left] {$\theta$};

% Légendes
\node[below] at (0, -1.8) {$O_1$};
\node[below] at (\Ltube, -1.3) {$O_2$};
\node[below] at (\focObj, -1.8) {$F'_1=F_2$};
\node[below] at (\Ltube+\focOcu, -1.3) {$F'_2$};
\draw[<->, popmuted] (0, -2.5) -- (\Ltube, -2.5) node[midway, below] {$L = \SI{10,5}{cm}$ (échelle 1/10)};
\end{tikzpicture}
\legende{Marche d'un faisceau lumineux issu du bord supérieur de la Lune (échelle 1/10). L'image intermédiaire $B_1$ est réelle, renversée, formée au foyer commun $F'_1=F_2$. Le faisceau émergent est parallèle, incliné de $\theta' = G\theta$ de part et d'autre de l'axe : l'image finale à l'infini est renversée.}
\end{popfigure}

\textbf{5. Argumentaire.} \emph{Monsieur le Proviseur, pour la Lune, notre lunette est idéale : son grossissement de 20 révèle les détails du relief (cratères, mers), et le renversement de l'image n'a aucune importance pour un astre, qui n'a ni haut ni bas. En revanche, pour suivre le match au stade, nous déconseillons la lunette : l'image y est renversée, ce qui rendrait le jeu incompréhensible ; son champ de vision est très étroit (environ $\SI{10}{\degree}/20 = \SI{0,5}{\degree}$ apparent), car un fort grossissement ne montre qu'une petite portion du stade ; enfin, le tube de \SI{1,05}{m} est encombrant et doit rester sur trépied. Des jumelles, qui redressent l'image grâce à des prismes de Porro ou de toit, offrent un champ large (typiquement $\SI{5}{\degree}$ à $\SI{8}{\degree}$), un grossissement modéré ($\times 7$ à $\times 10$) et une excellente maniabilité : elles sont l'instrument pertinent pour le match.}
\end{exoresolu}

\courssub{Bilan de l'activité}

\begin{aretenir}[Bilan synthétique]
Cette activité a permis de mettre en évidence plusieurs idées essentielles de la leçon :
\begin{itemize}
\item le \cle{choix des lentilles} n'est pas arbitraire : la condition $f'_1 \gg f'_2$ (objectif de grande focale, oculaire de courte focale) découle directement de la formule du grossissement $G = f'_1/f'_2$ ;
\item la \cle{configuration afocale} fixe entièrement la géométrie de l'instrument : $L = f'_1 + f'_2$, et la vérification de la faisabilité ($\SI{105}{cm}$ dans un tube réglable jusqu'à \SI{150}{cm}) fait partie du travail de l'ingénieur autant que du physicien ;
\item le grossissement relie deux \cle{angles} et non des tailles : c'est l'angle apparent de la Lune qui passe de $\SI{0,5}{\degree}$ à $\SI{10}{\degree}$ ;
\item le tracé de la marche du faisceau a concrétisé le rôle de l'\cle{image intermédiaire} dans le plan focal commun et expliqué le \cle{renversement} de l'image finale ;
\item enfin, aucun instrument n'est universel : le \cle{choix raisonné} d'un instrument d'optique dépend du besoin (grossissement, sens de l'image, champ de vision, maniabilité). La lunette astronomique excelle pour les astres ; les jumelles, à image redressée et champ large, sont préférables pour les scènes terrestres.
\end{itemize}
Tu es désormais capable de concevoir, dimensionner et justifier un instrument d'optique complet : une compétence directement mobilisable au Probatoire comme dans la vie scientifique quotidienne.
\end{aretenir}

\newpage

\coursec{Méthodes et exercices résolus}

\courssub{Expliquer le fonctionnement d'une lunette astronomique}

\begin{methode}[Décrire et expliquer le fonctionnement d'une lunette astronomique]
Pour répondre à une question du type « Expliquer le fonctionnement d'une lunette astronomique », suis la démarche suivante :
\begin{enumerate}
\item \textbf{Identifier les éléments} : la lunette comprend un \cle{objectif} (lentille convergente $L_1$ de grande distance focale $f'_1$, du côté de l'objet observé) et un \cle{oculaire} (lentille convergente $L_2$ de petite distance focale $f'_2$, du côté de l'œil).
\item \textbf{Décrire la chaîne des images} : l'objet $AB$ (très éloigné, quasi à l'infini) donne à travers $L_1$ une image \emph{intermédiaire} $A_1B_1$, réelle, renversée, située dans le plan focal image de $L_1$. Cette image $A_1B_1$ joue le rôle d'\emph{objet} pour l'oculaire $L_2$.
\item \textbf{Préciser le rôle de l'oculaire} : il fonctionne comme une \cle{loupe} ; il donne de $A_1B_1$ une image définitive $A'B'$ que l'œil peut observer sans fatigue si elle est rejetée à l'infini.
\item \textbf{Conclure sur la condition d'afocalité} : pour une vision à l'infini (œil emmétrope au repos, sans accommodation), il faut que $A_1B_1$ soit dans le plan focal objet de $L_2$, c'est-à-dire que le foyer image $F'_1$ de l'objectif coïncide avec le foyer objet $F_2$ de l'oculaire : la lunette est dite \cle{afocale}.
\item \textbf{Schématiser} : toujours accompagner l'explication d'un schéma de principe avec les foyers $F_1$, $F'_1 \equiv F_2$, $F'_2$ et les centres optiques $O_1$, $O_2$ alignés sur l'axe optique principal.
\end{enumerate}
\textbf{Réflexe} : la distance entre les deux lentilles d'une lunette afocale vaut $O_1O_2 = f'_1 + f'_2$.
\end{methode}

\begin{aretenir}[L'essentiel sur le fonctionnement]
Une lunette astronomique est l'association de deux lentilles minces convergentes de même axe optique : l'\cle{objectif} $L_1$ (grande focale $f'_1$) et l'\cle{oculaire} $L_2$ (petite focale $f'_2$). Dans le réglage \cle{afocal} (vision à l'infini sans fatigue) :
\[ F'_1 \equiv F_2 \qquad \text{et} \qquad O_1O_2 = f'_1 + f'_2. \]
Un faisceau parallèle entrant ressort parallèle : le système ne possède ni foyer objet ni foyer image, d'où le nom « afocal ».
\end{aretenir}

\begin{exoresolu}[Lunette d'observation du ciel de Yaoundé]
Un élève du lycée de Nkol-Eton fabrique une lunette astronomique avec deux lentilles minces convergentes : un objectif $L_1$ de vergence $C_1 = 1,0\ \delta$ et un oculaire $L_2$ de vergence $C_2 = 20\ \delta$.
\begin{enumerate}
\item Calculer les distances focales des deux lentilles.
\item Quelle doit être la distance entre les centres optiques des deux lentilles pour que la lunette soit afocale ? Justifier.
\item Expliquer en quelques lignes le fonctionnement de cette lunette.
\end{enumerate}
\tcblower
\textbf{1. Distances focales.} La distance focale image est l'inverse de la vergence : $f' = \dfrac{1}{C}$ (avec $C$ en dioptries et $f'$ en mètres).
\begin{align*}
f'_1 &= \dfrac{1}{C_1} = \dfrac{1}{1,0} = 1,0\ \text{m} ;\\
f'_2 &= \dfrac{1}{C_2} = \dfrac{1}{20} = 0,050\ \text{m} = 5,0\ \text{cm}.
\end{align*}
L'objectif a bien une grande distance focale et l'oculaire une petite distance focale : la répartition est conforme à la constitution d'une lunette astronomique.

\textbf{2. Distance entre les lentilles.} Pour une vision à l'infini sans accommodation, la lunette doit être \cle{afocale} : le foyer image $F'_1$ de l'objectif doit coïncider avec le foyer objet $F_2$ de l'oculaire. La distance entre les centres optiques est donc :
\[ O_1O_2 = O_1F'_1 + F_2O_2 = f'_1 + f'_2 = 1,0 + 0,050 = 1,05\ \text{m}. \]
Justification : l'objet céleste est à l'infini, donc son image $A_1B_1$ par $L_1$ se forme dans le plan focal image de $L_1$ ; pour que l'oculaire en donne une image définitive à l'infini, $A_1B_1$ doit se trouver dans son plan focal objet, d'où $F'_1 \equiv F_2$.

\textbf{3. Fonctionnement.} L'objectif $L_1$ reçoit la lumière d'un astre très éloigné (faisceau de rayons parallèles entre eux, inclinés d'un angle $\theta$ sur l'axe) et en forme une image intermédiaire $A_1B_1$ réelle et renversée dans son plan focal image. L'oculaire $L_2$, utilisé en loupe, donne de $A_1B_1$ une image définitive virtuelle rejetée à l'infini, vue sous un angle $\theta'$ supérieur à $\theta$ : l'astre apparaît ainsi sous un diamètre apparent plus grand, ce qui permet de distinguer des détails invisibles à l'œil nu.

\textbf{Conclusion} : $f'_1 = 1,0\ \text{m}$, $f'_2 = 5,0\ \text{cm}$, et la lunette afocale a une longueur $O_1O_2 = 1,05\ \text{m}$.
\end{exoresolu}

\courssub{Faire la mise au point à l'aide d'une lunette astronomique}

\begin{methode}[Régler une lunette : la mise au point]
\begin{enumerate}
\item \textbf{Viser l'objet} : pointer la lunette vers l'objet (astre, clocher lointain, antenne relais) et repérer l'image dans l'oculaire.
\item \textbf{Régler le tirage} : déplacer l'oculaire par rapport à l'objectif (coulissement du tube oculaire) jusqu'à obtenir une image \emph{nette}.
\item \textbf{Adapter au réglage afocal} : pour un objet à l'infini et un œil normal, la netteté sans fatigue est obtenue lorsque $O_1O_2 = f'_1 + f'_2$ (réglage afocal).
\item \textbf{Objet à distance finie} : l'image intermédiaire $A_1B_1$ se forme alors \emph{au-delà} de $F'_1$ (car $\overline{O_1A_1} > f'_1$ d'après la relation de conjugaison) ; il faut donc \emph{allonger} légèrement la lunette en éloignant l'oculaire de l'objectif pour ramener $A_1B_1$ dans le plan focal objet de $L_2$.
\item \textbf{Vérifier la qualité} : l'image définitive doit être nette, renversée (c'est normal pour une lunette astronomique) et observable sans effort d'accommodation.
\end{enumerate}
\textbf{Réflexe} : « objet plus proche $\Rightarrow$ image intermédiaire plus loin $\Rightarrow$ on tire l'oculaire vers l'extérieur ».
\end{methode}

\begin{exoresolu}[Mise au point sur un objet rapproché]
Une lunette astronomique est constituée d'un objectif de distance focale $f'_1 = 80\ \text{cm}$ et d'un oculaire de distance focale $f'_2 = 2,0\ \text{cm}$, réglée de façon afocale pour observer la Lune.
\begin{enumerate}
\item Quelle est la longueur de la lunette dans ce réglage ?
\item On observe maintenant le sommet d'un pylône de la SONATREL situé à $D = 400\ \text{m}$ de l'objectif. Déterminer la position de l'image intermédiaire $A_1B_1$ donnée par l'objectif.
\item De quelle distance et dans quel sens faut-il déplacer l'oculaire pour que l'image définitive reste à l'infini ?
\end{enumerate}
\tcblower
\textbf{1. Longueur afocale.} Dans le réglage afocal, $F'_1 \equiv F_2$, donc :
\[ O_1O_2 = f'_1 + f'_2 = 80 + 2,0 = 82\ \text{cm}. \]

\textbf{2. Position de l'image intermédiaire.} L'objet est à distance finie : $\overline{O_1A} = -D = -400\ \text{m}$. Appliquons la relation de conjugaison des lentilles minces à l'objectif :
\[ \dfrac{1}{\overline{O_1A_1}} - \dfrac{1}{\overline{O_1A}} = \dfrac{1}{f'_1}
\quad\Longrightarrow\quad
\dfrac{1}{\overline{O_1A_1}} = \dfrac{1}{f'_1} + \dfrac{1}{\overline{O_1A}} = \dfrac{1}{0,80} - \dfrac{1}{400}. \]
\[ \dfrac{1}{\overline{O_1A_1}} = 1,25 - 0,0025 = 1,2475\ \text{m}^{-1}
\quad\Longrightarrow\quad
\overline{O_1A_1} = \dfrac{1}{1,2475} \approx 0,8016\ \text{m} \approx 80,16\ \text{cm}. \]
L'image intermédiaire se forme à environ $80,16\ \text{cm}$ de l'objectif, soit $1,6\ \text{mm}$ au-delà de $F'_1$ : l'écart est faible car l'objet, quoique à distance finie, reste très éloigné devant $f'_1$.

\textbf{3. Déplacement de l'oculaire.} Pour que l'image définitive soit à l'infini, $A_1B_1$ doit coïncider avec le foyer objet $F_2$ de l'oculaire. La nouvelle distance objectif--oculaire doit donc être :
\[ O_1O_2' = \overline{O_1A_1} + f'_2 = 80,16 + 2,0 = 82,16\ \text{cm}. \]
Il faut donc \emph{éloigner} l'oculaire de l'objectif de :
\[ \Delta = 82,16 - 82 = 0,16\ \text{cm} = 1,6\ \text{mm}. \]

\textbf{Conclusion} : la lunette afocale mesure $82\ \text{cm}$ ; pour observer le pylône à $400\ \text{m}$, on tire l'oculaire vers l'extérieur d'environ $1,6\ \text{mm}$. C'est exactement le rôle de la bague (ou de la crémaillère) de mise au point.
\end{exoresolu}

\courssub{Tracer la marche d'un faisceau lumineux à travers une lunette}

\begin{methode}[Tracer les rayons lumineux dans une lunette afocale]
\begin{enumerate}
\item \textbf{Placer les éléments} à l'échelle : axe optique, centres $O_1$ et $O_2$, foyers avec $F'_1 \equiv F_2$ (lunette afocale). Choisir une échelle adaptée (ex. $1\ \text{cm}$ sur le papier pour $10\ \text{cm}$ réels).
\item \textbf{Rayon incident incliné} : l'objet à l'infini envoie un faisceau de rayons \emph{parallèles entre eux}, faisant l'angle $\theta$ avec l'axe. Le rayon passant par $O_1$ n'est pas dévié ; tous les rayons du faisceau convergent dans le plan focal image de $L_1$, au point $B_1$ où le rayon non dévié coupe ce plan.
\item \textbf{Construire l'image intermédiaire} $A_1B_1$ : $A_1$ est sur l'axe (dans le plan focal image de $L_1$), $B_1$ est le point de convergence du faisceau. On a $A_1B_1 = f'_1 \tan\theta \approx f'_1\,\theta$.
\item \textbf{Traversée de l'oculaire} : $B_1$ étant dans le plan focal objet de $L_2$, tous les rayons issus de $B_1$ ressortent de $L_2$ \emph{parallèles entre eux}. Pour trouver leur direction, tracer le rayon auxiliaire issu de $B_1$ passant par $O_2$ (non dévié) : les rayons émergents lui sont parallèles. Ils font l'angle $\theta'$ avec l'axe.
\item \textbf{Vérifier} : le faisceau entrant parallèle ressort parallèle (système afocal), avec $|\theta'| > |\theta|$ : l'image à l'infini est vue sous un angle plus grand et est \emph{renversée} ($\theta'$ et $\theta$ de signes opposés).
\end{enumerate}
\textbf{Réflexe} : deux familles de rayons parallèles (avant $L_1$ et après $L_2$), un point de convergence intermédiaire $B_1$ au foyer commun $F'_1 \equiv F_2$.
\end{methode}

\begin{popfigure}
\begin{center}
\begin{tikzpicture}[scale=1.0,>=Stealth]
% Axe optique
\draw[->,popmuted] (-1.2,0) -- (9.6,0) node[below left] {\footnotesize axe optique};
% Lentilles (échelle : f'1 = 6 cm, f'2 = 1.5 cm sur le schéma)
\draw[popblue,very thick] (0,-1.6) -- (0,1.6);
\draw[popblue,->,thick] (0,1.6) -- (0.18,1.42);
\draw[popblue,->,thick] (0,-1.6) -- (0.18,-1.42);
\node[popblue,below] at (0,-1.7) {\footnotesize $L_1$ (objectif)};
\draw[poporange,very thick] (7.5,-1.1) -- (7.5,1.1);
\draw[poporange,->,thick] (7.5,1.1) -- (7.68,0.92);
\draw[poporange,->,thick] (7.5,-1.1) -- (7.68,-0.92);
\node[poporange,below] at (7.5,-1.7) {\footnotesize $L_2$ (oculaire)};
% Centres et foyers
\fill (0,0) circle (1pt) node[below=2pt] {\footnotesize $O_1$};
\fill (7.5,0) circle (1pt) node[below=2pt] {\footnotesize $O_2$};
\fill (6,0) circle (1pt);
\node[below=2pt] at (6,0) {\footnotesize $F'_1{\equiv}F_2$};
\fill (9,0) circle (1pt) node[below=2pt] {\footnotesize $F'_2$};
% Plan focal image de L1
\draw[popmuted,dashed] (6,-1.5) -- (6,1.5);
% Rayons incidents parallèles (angle theta, pente -0.1)
\draw[popgreen,thick,->] (-1.0,1.30) -- (0,1.20);
\draw[popgreen,thick] (0,1.20) -- (6,-0.60);
\draw[popgreen,thick,->] (-1.0,0.70) -- (0,0.60);
\draw[popgreen,thick] (0,0.60) -- (6,-0.60);
\draw[popgreen,thick,->] (-1.0,0.10) -- (0,0.00);
\draw[popgreen,thick] (0,0.00) -- (6,-0.60);
% Point B1
\fill[popink] (6,-0.60) circle (1.5pt) node[below right] {\footnotesize $B_1$};
\fill[popink] (6,0) circle (1pt);
\node[above=1pt] at (6.15,0.05) {\footnotesize $A_1$};
% Rayons de B1 vers L2
\draw[popgreen,thick] (6,-0.60) -- (7.5,0.60);
\draw[popgreen,thick] (6,-0.60) -- (7.5,0.00);
\draw[popgreen,thick] (6,-0.60) -- (7.5,-0.60);
% Rayons émergents parallèles (direction de B1O2 prolongée : pente +0.8)
\draw[popgreen,thick,->] (7.5,0.60) -- (9.4,2.12);
\draw[popgreen,thick,->] (7.5,0.00) -- (9.4,1.52);
\draw[popgreen,thick,->] (7.5,-0.60) -- (9.4,0.92);
% Angles
\draw[popink] (2.2,0) arc (0:-5.7:2.2);
\node[popink] at (2.6,-0.28) {\footnotesize $\theta$};
\draw[popink] (8.6,0) arc (0:38.7:1.1);
\node[popink] at (9.05,0.42) {\footnotesize $\theta'$};
\end{tikzpicture}
\end{center}
\legende{Marche d'un faisceau parallèle à travers une lunette afocale : convergence en $B_1$ dans le plan focal image de $L_1$, puis émergence en faisceau parallèle sous l'angle $\theta'$, de sens opposé à $\theta$ (image renversée). Les angles sont exagérés pour la lisibilité.}
\end{popfigure}

\begin{exoresolu}[Tracé et mesure sur un schéma à l'échelle]
Une lunette afocale est formée d'un objectif $L_1$ ($f'_1 = 50\ \text{cm}$) et d'un oculaire $L_2$ ($f'_2 = 5,0\ \text{cm}$). Elle est pointée vers l'étoile Polaire, dont les rayons arrivent sous un angle $\theta = 0,010\ \text{rad}$ par rapport à l'axe optique.
\begin{enumerate}
\item Calculer la taille de l'image intermédiaire $A_1B_1$.
\item Faire un schéma (décrit précisément) de la marche d'un faisceau de trois rayons parallèles à travers la lunette.
\item Déterminer l'angle émergent $\theta'$ et conclure sur le caractère de l'image définitive.
\end{enumerate}
\tcblower
\textbf{1. Taille de l'image intermédiaire.} L'image d'un objet à l'infini vu sous l'angle $\theta$ se forme dans le plan focal image de l'objectif, avec :
\[ A_1B_1 = f'_1 \tan\theta \approx f'_1\,\theta = 50 \times 0,010 = 0,50\ \text{cm} = 5,0\ \text{mm}. \]

\textbf{2. Schéma de la marche des rayons.} On adopte l'échelle $1/5$ : $O_1O_2 = f'_1 + f'_2 = 55\ \text{cm}$ réels, soit $11\ \text{cm}$ sur le schéma.
\begin{itemize}
\item trois rayons incidents parallèles entre eux, inclinés de $\theta$ vers le bas (par exemple), frappent $L_1$ ;
\item le rayon passant par $O_1$ traverse sans déviation et coupe le plan focal image de $L_1$ en $B_1$, à $5\ \text{mm}$ réels ($1\ \text{mm}$ sur le schéma) au-dessous de l'axe ; les deux autres rayons convergent aussi en $B_1$ ;
\item de $B_1$, situé au foyer objet $F_2$ de l'oculaire, les rayons divergent vers $L_2$ et ressortent \emph{parallèles entre eux}, dans la direction du rayon auxiliaire $B_1O_2$ prolongé, inclinée de $\theta'$ vers le haut.
\end{itemize}

\textbf{3. Angle émergent et conclusion.} Le rayon $B_1O_2$ fait avec l'axe un angle $\theta'$ tel que $\tan\theta' = \dfrac{A_1B_1}{f'_2}$, donc, aux petits angles :
\[ \theta' \approx \dfrac{A_1B_1}{f'_2} = \dfrac{0,50}{5,0} = 0,10\ \text{rad}. \]
L'angle $\theta'$ est de sens \emph{opposé} à $\theta$ : l'image définitive, rejetée à l'infini, est \emph{renversée} par rapport à l'objet, et elle est vue sous un angle dix fois plus grand ($\theta'/\theta = 10$).

\textbf{Conclusion} : $A_1B_1 = 5,0\ \text{mm}$ ; le faisceau parallèle incident ressort parallèle sous l'angle $\theta' = 0,10\ \text{rad}$ : image renversée et « grossie » angulairement d'un facteur $10$.
\end{exoresolu}

\courssub{Calculer le grossissement d'une lunette astronomique}

\begin{definition}[Grossissement et puissance]
Le \cle{grossissement} d'une lunette est le rapport :
\[ G = \dfrac{\theta'}{\theta} \]
où $\theta$ est l'angle sous lequel l'objet est vu à l'œil nu et $\theta'$ l'angle sous lequel l'image définitive est vue à travers l'instrument. $G$ est \emph{sans dimension} ; son signe négatif traduit le renversement de l'image.

La \cle{puissance} de l'instrument est le rapport :
\[ P = \dfrac{\theta'}{A_1B_1} \]
angle émergent par unité de taille de l'image intermédiaire ; elle s'exprime en dioptries ($\delta = \text{m}^{-1}$) si les longueurs sont en mètres et les angles en radians.
\end{definition}

\begin{propriete}[Grossissement et puissance de la lunette afocale]
Pour une lunette réglée de façon afocale, l'image intermédiaire $A_1B_1$ est commune aux deux lentilles, avec :
\[ \theta \approx \dfrac{A_1B_1}{f'_1} \qquad \text{et} \qquad \theta' \approx \dfrac{A_1B_1}{f'_2}. \]
On en déduit immédiatement :
\[ \boxed{\ G = -\dfrac{f'_1}{f'_2} = -\dfrac{C_2}{C_1}\ } \qquad \text{et} \qquad \boxed{\ P = \dfrac{\theta'}{A_1B_1} = \dfrac{1}{f'_2} = C_2\ } \]
d'où la relation de cohérence $|G| = P \times f'_1$. Le signe moins traduit le renversement de l'image ($\theta'$ et $\theta$ de sens opposés). Cette démonstration, simple application des définitions, est exigible au Probatoire.
\end{propriete}

\begin{methode}[Calculer grossissement et puissance]
\begin{enumerate}
\item \textbf{Identifier les focales} : $f'_1$ pour l'objectif (grande focale), $f'_2$ pour l'oculaire (petite focale), et \emph{convertir dans la même unité} avant tout rapport.
\item \textbf{Appliquer la formule de la lunette afocale} : $G = -\dfrac{f'_1}{f'_2}$. En valeur absolue, $|G| = \dfrac{f'_1}{f'_2}$ : le grossissement est d'autant plus grand que l'objectif a une grande focale et l'oculaire une petite focale.
\item \textbf{Retrouver la formule si besoin} : écrire $\theta \approx \dfrac{A_1B_1}{f'_1}$ et $\theta' \approx \dfrac{A_1B_1}{f'_2}$, puis faire le rapport en tenant compte du sens des angles.
\item \textbf{Calculer la puissance} : $P = \dfrac{\theta'}{A_1B_1} = \dfrac{1}{f'_2}$, avec $f'_2$ en \emph{mètres} pour obtenir $P$ en dioptries. Vérifier la cohérence par $|G| = P \times f'_1$.
\item \textbf{Vérifier l'homogénéité} : $G$ est un rapport de deux angles (sans unité) ; $P$ est l'inverse d'une longueur ($\delta = \text{m}^{-1}$).
\end{enumerate}
\textbf{Réflexe} : convertir toutes les focales dans la même unité \emph{avant} de faire le rapport, et en mètres pour la puissance.
\end{methode}

\begin{exoresolu}[Grossissement d'une lunette d'astronomie amateur]
Un club scientifique de Douala observe les cratères lunaires avec une lunette afocale dont l'objectif a pour distance focale $f'_1 = 900\ \text{mm}$ et l'oculaire $f'_2 = 15\ \text{mm}$.
\begin{enumerate}
\item Calculer le grossissement de la lunette. Interpréter son signe.
\item La Lune est vue à l'œil nu sous un diamètre apparent $\theta = 9,3\times 10^{-3}\ \text{rad}$. Sous quel angle apparent est-elle vue à travers la lunette ?
\item Calculer la puissance de la lunette, puis retrouver le grossissement à partir de $P$ et $f'_1$.
\item On dispose d'un second oculaire de focale $f''_2 = 25\ \text{mm}$. Lequel choisir pour observer un détail fin des cratères ? Justifier.
\end{enumerate}
\tcblower
\textbf{1. Grossissement.} Pour une lunette afocale :
\[ G = -\dfrac{f'_1}{f'_2} = -\dfrac{900}{15} = -60. \]
Le signe négatif signifie que l'image est \emph{renversée} par rapport à l'objet ; la valeur absolue $|G| = 60$ signifie que l'astre est vu sous un angle $60$ fois plus grand.

\textbf{2. Diamètre apparent à travers la lunette.} Par définition $G = \dfrac{\theta'}{\theta}$, donc en valeur absolue :
\[ \theta' = |G| \times \theta = 60 \times 9,3\times 10^{-3} = 0,56\ \text{rad} \approx 32^{\circ}. \]
La Lune apparaît sous un diamètre apparent d'environ $0,56\ \text{rad}$, très confortable pour l'observation.

\textbf{3. Puissance et vérification.} La puissance de la lunette afocale est :
\[ P = \dfrac{1}{f'_2} = \dfrac{1}{0,015} \approx 67\ \delta. \]
Le grossissement en valeur absolue s'en déduit par $|G| = P \times f'_1 = 66,7 \times 0,900 = 60$ : on retrouve exactement le résultat de la question 1, ce qui valide la cohérence des formules.

\textbf{4. Choix de l'oculaire.} Avec l'oculaire de $25\ \text{mm}$ : $|G''| = \dfrac{900}{25} = 36 < 60$. Pour observer un détail fin, il faut le grossissement le plus grand : on choisit donc l'oculaire de $f'_2 = 15\ \text{mm}$, car le grossissement est inversement proportionnel à la focale de l'oculaire.

\textbf{Conclusion} : $G = -60$ (image renversée, vue $60$ fois plus grande angulairement) ; $\theta' \approx 0,56\ \text{rad}$ ; $P \approx 67\ \delta$ ; l'oculaire de $15\ \text{mm}$ est le plus adapté aux détails fins.
\end{exoresolu}

\courssub{Choisir de manière pertinente un instrument d'optique}

\begin{methode}[Choisir et justifier un instrument d'optique]
\begin{enumerate}
\item \textbf{Analyser le besoin} : observer un objet \emph{très éloigné} (astre, paysage, animal) ou un objet \emph{proche et petit} (cellule, circuit électronique, insecte) ?
\item \textbf{Associer l'instrument} :
\begin{itemize}
\item objet lointain, besoin d'un fort grossissement angulaire $\Rightarrow$ \cle{lunette astronomique} (image renversée, acceptable en astronomie) ;
\item objet lointain terrestre où l'image doit être droite et l'encombrement réduit $\Rightarrow$ \cle{jumelles} (lunette terrestre à prismes redresseurs) ;
\item objet proche de petite taille $\Rightarrow$ \cle{loupe} (grossissement modeste) ou \cle{microscope} (fort grossissement).
\end{itemize}
\item \textbf{Justifier quantitativement} : comparer les grossissements possibles ($|G| = f'_1/f'_2$), l'encombrement ($O_1O_2 = f'_1 + f'_2$), la clarté (diamètre de l'objectif, qui fixe la quantité de lumière collectée) et le caractère droit ou renversé de l'image.
\item \textbf{Dimensionner si demandé} : choisir les focales pour atteindre un grossissement imposé, en vérifiant que l'encombrement reste raisonnable.
\item \textbf{Rédiger la justification} : une phrase liant le besoin, la grandeur pertinente et la valeur numérique obtenue.
\end{enumerate}
\end{methode}

\begin{exoresolu}[Équipement d'un poste d'observation du parc de Waza]
L'Office de tourisme veut équiper un poste d'observation du parc national de Waza. Deux besoins sont exprimés :
\begin{itemize}
\item[\textbullet] Besoin A : observer des éléphants à plus de $2\ \text{km}$, avec une image droite et un appareil compact ;
\item[\textbullet] Besoin B : initier les visiteurs à l'observation de Jupiter et de ses satellites, avec le plus fort grossissement possible.
\end{itemize}
On dispose de : une lunette astronomique ($f'_1 = 80\ \text{cm}$, $f'_2 = 2\ \text{cm}$), des jumelles ($G = 8$, longueur $15\ \text{cm}$, image redressée par des prismes), une loupe ($f' = 5\ \text{cm}$).
\begin{enumerate}
\item Attribuer à chaque besoin l'instrument le plus pertinent en justifiant.
\item Calculer le grossissement et la longueur de la lunette astronomique.
\end{enumerate}
\tcblower
\textbf{1. Attribution des instruments.}
\begin{itemize}
\item \emph{Besoin A} : les \textbf{jumelles} conviennent. L'objet est terrestre : l'image doit être \emph{droite} (grâce au redressement par prismes), l'appareil doit être \emph{compact} ($15\ \text{cm}$ contre $82\ \text{cm}$ pour la lunette) et un grossissement de $8$ suffit pour un animal de grande taille à $2\ \text{km}$.
\item \emph{Besoin B} : la \textbf{lunette astronomique} convient. En astronomie, le renversement de l'image est sans importance (Jupiter n'a pas de « haut » ni de « bas »), et son grossissement est bien supérieur à celui des jumelles, ce qui permet de distinguer les satellites galiléens. La loupe, prévue pour les objets proches, est inadaptée aux deux besoins.
\end{itemize}

\textbf{2. Caractéristiques de la lunette.}
\[ G = -\dfrac{f'_1}{f'_2} = -\dfrac{80}{2} = -40 \quad ; \quad O_1O_2 = f'_1 + f'_2 = 82\ \text{cm}. \]
Son grossissement ($40$ en valeur absolue) est cinq fois celui des jumelles : c'est bien l'instrument le plus performant pour le besoin B, au prix d'un encombrement plus grand et d'une image renversée.

\textbf{Conclusion} : jumelles pour l'observation terrestre (image droite, compacité, $G = 8$ suffisant) ; lunette astronomique ($G = -40$, longueur $82\ \text{cm}$) pour l'astronomie, où le fort grossissement prime sur le redressement de l'image.
\end{exoresolu}

\begin{savaistu}[Galilée et la lunette]
En 1609, Galilée pointe vers le ciel une lunette de sa fabrication (grossissement d'environ $20$) : il découvre les montagnes de la Lune, les quatre satellites principaux de Jupiter (dits « galiléens ») et les phases de Vénus, apportant des arguments décisifs en faveur du modèle héliocentrique. Quatre siècles plus tard, les grands télescopes modernes fonctionnent sur le même principe de grossissement angulaire, avec des miroirs en lieu et place de l'objectif à lentille. Au Cameroun, les clubs d'astronomie des universités de Yaoundé et de Douala organisent régulièrement des nuits d'observation avec des lunettes dont le grossissement se calcule exactement comme dans cette leçon.
\end{savaistu}

\begin{attention}[Les 5 erreurs qui coûtent des points au Probatoire]
\begin{enumerate}
\item \textbf{Inverser objectif et oculaire.} L'\emph{objectif} (grande focale, côté objet) et l'\emph{oculaire} (petite focale, côté œil) sont souvent confondus. Dans $G = -f'_1/f'_2$, $f'_1$ est la focale de l'objectif : intervertir donne un grossissement inférieur à $1$ en valeur absolue, absurde pour une lunette.
\item \textbf{Oublier la condition d'afocalité.} Écrire que la lunette est afocale exige de justifier : $F'_1 \equiv F_2$, donc $O_1O_2 = f'_1 + f'_2$. Beaucoup de copies utilisent cette formule sans dire \emph{pourquoi}, ou l'appliquent alors que l'énoncé impose une mise au point différente (objet à distance finie : la lunette est alors légèrement allongée).
\item \textbf{Mélanger les unités dans le rapport des focales.} Calculer $G = 900/1,5 = 600$ avec $f'_1$ en mm et $f'_2$ en cm est faux : convertir d'abord dans la même unité ($900/15 = 60$). La puissance $P = 1/f'_2$ exige $f'_2$ en \emph{mètres} pour obtenir des dioptries.
\item \textbf{Tracer des rayons faux.} Erreurs classiques : faire converger le faisceau parallèle incident au foyer $F'_1$ \emph{sur l'axe} (il converge en $B_1$, hors de l'axe, dans le plan focal image) ; oublier que les rayons ressortent de l'oculaire \emph{parallèles entre eux} ; oublier le rayon auxiliaire passant par $O_2$, non dévié, qui fixe la direction émergente.
\item \textbf{Négliger le signe et l'interprétation du grossissement.} Écrire $G = 40$ sans commenter le signe, ou affirmer que l'image est « agrandie » : le grossissement est un rapport d'\emph{angles}, pas de tailles ; son signe négatif traduit le \emph{renversement} de l'image. Toujours conclure par une phrase : image renversée, vue sous un angle $|G|$ fois plus grand.
\end{enumerate}
\end{attention}

\newpage

\coursec{Exercices}

\coursec{La lunette astronomique}

\courssub{Je vérifie mes connaissances}

\begin{exercice}[QCM : la lunette astronomique]
Pour chaque question, une seule réponse est exacte. Recopie la lettre correspondante sur ta copie.
\begin{enumerate}
\item Dans une lunette astronomique, l'objectif est une lentille :
\begin{enumerate}
\item[a)] divergente de grande distance focale ;
\item[b)] convergente de grande distance focale ;
\item[c)] convergente de courte distance focale ;
\item[d)] divergente de courte distance focale.
\end{enumerate}
\item Une lunette est dite \cle{afocale} lorsque :
\begin{enumerate}
\item[a)] l'image définitive est rejetée à l'infini ;
\item[b)] l'image définitive se forme sur la rétine sans accommodation ;
\item[c)] le foyer image de l'objectif coïncide avec le foyer objet de l'oculaire ;
\item[d)] les trois réponses a), b) et c) sont vraies.
\end{enumerate}
\item Le grossissement d'une lunette afocale est donné par :
\begin{enumerate}
\item[a)] $G = \dfrac{f'_2}{f'_1}$ ; \qquad b) $G = f'_1 \times f'_2$ ; \qquad c) $G = \dfrac{f'_1}{f'_2}$ ; \qquad d) $G = f'_1 + f'_2$.
\end{enumerate}
\item L'image d'un objet céleste donnée par une lunette astronomique est :
\begin{enumerate}
\item[a)] droite et réelle ;
\item[b)] renversée et virtuelle ;
\item[c)] droite et virtuelle ;
\item[d)] renversée et réelle.
\end{enumerate}
\end{enumerate}
\end{exercice}

\begin{exercice}[Vrai ou faux ? Justifier]
Réponds par \textbf{vrai} ou \textbf{faux}, puis justifie ta réponse en une ou deux phrases.
\begin{enumerate}
\item Pour augmenter le grossissement d'une lunette, on peut remplacer l'oculaire par un oculaire de distance focale plus grande.
\item Dans une lunette afocale, un faisceau incident parallèle à l'axe optique émerge sous forme d'un faisceau parallèle à l'axe optique.
\item La puissance d'une lunette s'exprime en dioptries ($\delta$).
\item La longueur d'une lunette afocale est égale à la somme des distances focales de l'objectif et de l'oculaire.
\end{enumerate}
\end{exercice}

\begin{exercice}[Définitions et conversions]
\begin{enumerate}
\item Définis les termes suivants : \cle{objectif}, \cle{oculaire}, \cle{cercle oculaire}, \cle{grossissement}.
\item Le diamètre apparent de la Lune vue à l'œil nu est d'environ $30$ minutes d'arc. Exprime cet angle en secondes d'arc, puis en radians. On donne : $1° = 60'$ ; $1' = 60''$ ; $180° = \pi\ \text{rad}$.
\item Une lunette a un grossissement $G = 25$. Rappelle la définition du grossissement d'une lunette afocale et précise la signification physique de cette valeur.
\end{enumerate}
\end{exercice}

\begin{exercice}[Calculs directs]
Une lunette astronomique afocale est constituée d'un objectif de distance focale $f'_1 = \SI{80}{cm}$ et d'un oculaire de distance focale $f'_2 = \SI{2}{cm}$.
\begin{enumerate}
\item Calcule la longueur $L$ de la lunette (distance entre l'objectif et l'oculaire).
\item Calcule son grossissement $G$.
\item Calcule la vergence de l'objectif $C_1$ et celle de l'oculaire $C_2$ en dioptries. Rappelle pourquoi la vergence totale de la lunette afocale est nulle.
\item Jupiter est vue à l'œil nu sous un diamètre apparent de $40$ secondes d'arc. Sous quel angle $\theta'$ la voit-on à travers cette lunette ? Exprime le résultat en secondes d'arc puis en minutes d'arc.
\end{enumerate}
\end{exercice}

\courssub{Je m'entraîne}

\begin{exercice}[Construction graphique d'une lunette afocale]
On dispose d'une lunette afocale dont l'objectif a pour distance focale $f'_1 = \SI{50}{cm}$ et l'oculaire $f'_2 = \SI{5}{cm}$.
\begin{enumerate}
\item Calcule la longueur de la lunette et son grossissement.
\item Sur une feuille de papier millimétré, trace à l'échelle $1/5$ l'axe optique, l'objectif $L_1$ et l'oculaire $L_2$ en plaçant correctement les foyers $F'_1$ et $F_2$.
\item Un faisceau de trois rayons parallèles entre eux, incliné de $\theta = 2°$ sur l'axe optique, arrive sur l'objectif. Trace la marche de ce faisceau à travers la lunette. Justifie chaque étape du tracé.
\item Mesure l'angle $\theta'$ du faisceau émergent par rapport à l'axe optique. Retrouve graphiquement la valeur du grossissement et compare au résultat de la question 1.
\end{enumerate}
\end{exercice}

\begin{exercice}[Le choix de l'oculaire]
Un astronome amateur de Yaoundé veut observer les anneaux de Saturne avec un grossissement $G = 50$. Il possède un objectif de distance focale $f'_1 = \SI{1}{m}$.
\begin{enumerate}
\item Quelle distance focale $f'_2$ d'oculaire doit-il choisir ?
\item Quelle sera la longueur de sa lunette une fois réglée en afocal ?
\item Il dispose en réalité de trois oculaires de distances focales $\SI{25}{mm}$, $\SI{20}{mm}$ et $\SI{10}{mm}$. Quel grossissement obtient-il avec chacun ? Lequel choisir pour approcher au mieux son objectif ?
\item Cite un inconvénient pratique d'un grossissement trop élevé pour l'observation astronomique.
\end{enumerate}
\end{exercice}

\begin{exercice}[Démonstration du grossissement]
On considère une lunette astronomique réglée en afocal, d'objectif de distance focale $f'_1$ et d'oculaire de distance focale $f'_2$. Un objet $AB$ à l'infini, situé dans la direction faisant l'angle $\theta$ avec l'axe optique, envoie un faisceau parallèle sur l'objectif.
\begin{enumerate}
\item Fais un schéma de principe de la lunette afocale et représente l'image intermédiaire $A_1B_1$ donnée par l'objectif. Où se forme-t-elle ?
\item Exprime la taille $A_1B_1$ de l'image intermédiaire en fonction de $f'_1$ et de $\theta$, en précisant l'approximation utilisée.
\item L'image intermédiaire joue le rôle d'objet pour l'oculaire. Exprime l'angle $\theta'$ sous lequel l'observateur voit l'image définitive, en fonction de $A_1B_1$ et de $f'_2$.
\item Déduis-en la relation $G = \dfrac{f'_1}{f'_2}$.
\item Vérifie numériquement cette relation avec $f'_1 = \SI{900}{mm}$ et $f'_2 = \SI{20}{mm}$, puis calcule l'angle sous lequel on voit la Lune ($\theta = 0,5°$) à travers cette lunette.
\end{enumerate}
\end{exercice}

\begin{exercice}[Mise au point : de la Lune au clocher]
Une lunette dont l'objectif a une distance focale $f'_1 = \SI{60}{cm}$ est d'abord réglée en afocal pour observer la Lune (objet à l'infini) sans fatigue : l'image définitive est rejetée à l'infini. L'oculaire a une distance focale $f'_2 = \SI{3}{cm}$.
\begin{enumerate}
\item Quelle est alors la distance entre l'objectif et l'oculaire ?
\item On observe maintenant le clocher d'une église situé à $D = \SI{500}{m}$ de l'objectif. En appliquant la relation de conjugaison des lentilles minces à l'objectif, calcule la position $\overline{O_1A_1}$ de l'image intermédiaire. On rappelle : $\dfrac{1}{\overline{O_1A_1}} - \dfrac{1}{\overline{O_1A}} = \dfrac{1}{f'_1}$.
\item Pour garder une image définitive à l'infini, l'image intermédiaire doit rester au foyer objet de l'oculaire. Dans quel sens et de quelle distance $\Delta$ doit-on déplacer l'oculaire ? Cette distance est appelée \emph{variation du tirage}.
\item Explique, schéma à l'appui, pourquoi la mise au point est nécessaire lorsqu'on passe d'un objet à l'infini à un objet rapproché.
\end{enumerate}
\end{exercice}

\courssub{Je me prépare au Probatoire}

\begin{exercice}[Observation de Jupiter au club d'astronomie de Douala]
Le club de sciences du lycée de Douala-Bonabéri construit une lunette astronomique à l'aide de deux lentilles minces convergentes : un objectif $L_1$ de vergence $C_1 = 1{,}25\ \delta$ et un oculaire $L_2$ de distance focale $f'_2 = \SI{4}{cm}$. La lunette est réglée en afocal.
\begin{enumerate}
\item Calcule la distance focale $f'_1$ de l'objectif.
\item Définis l'expression « lunette afocale » et précise la position relative des foyers $F'_1$ et $F_2$. Fais le schéma de principe de la lunette en indiquant les foyers des deux lentilles.
\item Calcule la longueur $L$ de la lunette et son grossissement $G$.
\item Jupiter est observée à l'œil nu sous un diamètre apparent $\theta = 46''$ (secondes d'arc).
\begin{enumerate}
\item[a)] Convertis $\theta$ en radians. On donne : $1° = 3600''$ et $180° = \pi\ \text{rad}$.
\item[b)] Calcule le diamètre apparent $\theta'$ de Jupiter vue à travers la lunette, en secondes d'arc.
\item[c)] Calcule la taille $A_1B_1$ de l'image intermédiaire de Jupiter dans le plan focal image de l'objectif. On utilisera l'approximation des petits angles : $\tan\theta \approx \theta$ (rad).
\end{enumerate}
\item Un élève propose de remplacer l'oculaire par un autre de distance focale $f''_2 = \SI{2}{cm}$. Calcule le nouveau grossissement. La longueur de la lunette est-elle modifiée ? Si oui, dans quel sens et de combien ?
\item Le professeur rappelle que l'image donnée par la lunette est renversée. Pourquoi ce défaut est-il sans gravité pour l'observation astronomique, mais gênant pour observer un paysage terrestre ?
\end{enumerate}
\end{exercice}

\begin{exercice}[La lunette du collège de Nkongsamba]
Pour la semaine des sciences, les élèves de Première C du collège de Nkongsamba disposent d'un banc d'optique, d'un objectif de distance focale $f'_1 = \SI{75}{cm}$, d'un oculaire de distance focale $f'_2 = \SI{3}{cm}$ et d'un écran.
\begin{enumerate}
\item Les élèves veulent d'abord vérifier la distance focale de l'objectif. Ils visent un bâtiment éloigné (considéré à l'infini) et recueillent son image nette sur l'écran. À quelle distance de l'objectif se trouve l'écran ? Justifie à l'aide de la relation de conjugaison.
\item Ils réalisent ensuite une lunette afocale.
\begin{enumerate}
\item[a)] Donne la distance objectif--oculaire à respecter et calcule le grossissement $G$ de la lunette.
\item[b)] Trace, à l'échelle $1/5$, la marche d'un faisceau de deux rayons parallèles arrivant sur l'objectif sous une inclinaison $\theta = 3°$ par rapport à l'axe optique. Décris précisément ton protocole de tracé.
\item[c)] Mesure sur ta figure l'angle émergent $\theta'$ et déduis-en une détermination graphique du grossissement. Compare à la valeur calculée.
\end{enumerate}
\item La pleine Lune est vue à l'œil nu sous un diamètre apparent de $31'$. Calcule le diamètre apparent de la Lune observée à travers cette lunette, en degrés.
\item Les élèves veulent ensuite photographier l'image intermédiaire en plaçant l'écran au foyer image de l'objectif, l'oculaire étant retiré. Calcule le diamètre de l'image de la Lune sur l'écran.
\item Pourquoi dit-on qu'un œil emmétrope (normal) placé derrière l'oculaire d'une lunette afocale observe « sans accommoder » ? Quel avantage cela présente-t-il ?
\end{enumerate}
\end{exercice}

\begin{exercice}[Safari photo dans la réserve de Waza : choisir son instrument]
Un touriste en visite dans la réserve de Waza (Extrême-Nord du Cameroun) souhaite observer des oiseaux et des antilopes. Il hésite entre deux instruments :
\begin{itemize}
\item une \textbf{lunette astronomique} : grossissement $G_1 = 60$, image renversée, longueur $\SI{1}{m}$, champ de vision étroit ;
\item des \textbf{jumelles} : grossissement $G_2 = 8$, image redressée grâce à des prismes, encombrement réduit, large champ de vision.
\end{itemize}
\begin{enumerate}
\item Rappelle la constitution d'une lunette astronomique et le rôle de chacun de ses deux éléments optiques.
\item La lunette astronomique du touriste a un objectif de distance focale $f'_1 = \SI{96}{cm}$. Calcule la distance focale $f'_2$ de son oculaire et vérifie que sa longueur en configuration afocale est bien voisine de $\SI{1}{m}$.
\item Une antilope de hauteur $h = \SI{1,2}{m}$ se trouve à $d = \SI{600}{m}$ de l'observateur.
\begin{enumerate}
\item[a)] Calcule le diamètre apparent $\theta$ (en radians puis en minutes d'arc) sous lequel l'antilope est vue à l'œil nu. On utilisera $\tan\theta \approx \theta$.
\item[b)] Calcule le diamètre apparent de l'antilope vue à travers la lunette, puis à travers les jumelles.
\end{enumerate}
\item En t'appuyant sur les résultats précédents et sur les caractéristiques des deux instruments (grossissement, redressement de l'image, champ de vision, encombrement, stabilité de la visée), rédige un texte argumenté de dix à quinze lignes conseillant au touriste l'instrument le plus pertinent pour son safari. Tu justifieras ton choix par au moins trois arguments physiques ou pratiques.
\item Le touriste possède aussi un appareil photo muni d'un téléobjectif assimilable à une lentille convergente de distance focale $f' = \SI{300}{mm}$. Calcule la taille de l'image de l'antilope sur le capteur placé dans le plan focal image. Ce résultat te semble-t-il exploitable ? Commente.
\end{enumerate}
\end{exercice}

\begin{aretenir}[Formules clés de la lunette astronomique afocale]
\begin{itemize}
\item \textbf{Condition d'afocalité} : $F'_1$ (foyer image de l'objectif) coïncide avec $F_2$ (foyer objet de l'oculaire).
\item \textbf{Longueur} : $L = f'_1 + f'_2$.
\item \textbf{Grossissement (grandissement angulaire)} : $G = \dfrac{\theta'}{\theta} = \dfrac{f'_1}{f'_2}$.
\item \textbf{Nature de l'image finale} : Virtuelle, renversée, à l'infini.
\item \textbf{Vergences} : $C_1 = \dfrac{1}{f'_1}$, $C_2 = \dfrac{1}{f'_2}$ (en $\delta$ si $f'$ en m). Vergence totale $C = C_1 + C_2 - d C_1 C_2 = 0$ pour $d = f'_1 + f'_2$.
\end{itemize}
\end{aretenir}

\begin{attention}[Pièges et erreurs fréquentes]
\begin{itemize}
\item \textbf{Puissance vs Grossissement} : En optique physique, la \emph{puissance} (ou vergence) s'exprime en dioptries ($\delta$). Une lunette afocale a une puissance \textbf{nulle}. Le terme « puissance » dans le commerce (ex: « puissance 20x ») désigne en réalité le \textbf{grossissement $G$} (sans unité). Ne confonds pas les deux.
\item \textbf{Signe du grossissement} : $G = f'_1/f'_2 > 0$ car les deux lentilles sont convergentes. L'image est renversée (le signe $+$ indique que le sens de rotation du faisceau émergent est le même que l'incident par rapport à l'axe, mais l'image est inversée haut/bas). Au Probatoire, on donne la valeur absolue $|G|$.
\item \textbf{Mise au point} : Pour un objet rapproché, l'image intermédiaire se forme \textbf{au-delà} de $F'_1$. Il faut \textbf{écarter} l'oculaire de l'objectif (augmenter le tirage) pour maintenir $F_2$ sur cette image.
\item \textbf{Unités} : Vérifie toujours l'homogénéité : $f'_1$ et $f'_2$ dans la \textbf{même unité} (m ou cm ou mm) pour calculer $G$ et $L$. Les angles $\theta, \theta'$ doivent être en \textbf{radian} pour les formules $\tan\theta \approx \theta$ et $A_1B_1 = f'_1 \theta$.
\end{itemize}
\end{attention}

\begin{methode}[Tracer la marche d'un faisceau dans une lunette afocale]
\begin{enumerate}
\item Trace l'axe optique, les lentilles $L_1$ (objectif) et $L_2$ (oculaire) espacées de $L = f'_1 + f'_2$. Place $F'_1$ et $F_2$ confondus au centre.
\item Dessine le faisceau incident : rayons parallèles inclinés d'un angle $\theta$ (souvent 3 rayons : un par le centre optique $O_1$, un passant par $F_1$, un parallèle à l'axe).
\item \textbf{Objectif $L_1$} : Le rayon passant par $O_1$ n'est pas dévié. Le rayon parallèle à l'axe passe par $F'_1$. Le rayon passant par $F_1$ sort parallèle à l'axe. Les trois rayons se coupent en $A_1$ (image intermédiaire) dans le plan focal de $L_1$.
\item \textbf{Oculaire $L_2$} : L'image intermédiaire $A_1$ est au foyer objet $F_2$ de $L_2$. Traite $A_1$ comme un objet pour $L_2$ : le rayon passant par $O_2$ n'est pas dévié ; le rayon passant par $F_2$ (donc par $A_1$ et $O_2$ ? Non, le rayon issu de $A_1$ passant par $O_2$ continue tout droit) — en fait, tout rayon issu de $A_1$ (au foyer) émerge \textbf{parallèle} à la droite $(O_2 A_1)$.
\item Le faisceau émergent est parallèle, incliné de $\theta'$ tel que $\tan\theta' \approx \theta' = \dfrac{A_1B_1}{f'_2}$.
\end{enumerate}
\end{methode}

\begin{savaistu}[La lunette au Cameroun : de l'astronomie à la télécommunication]
Au Cameroun, l'optique instrumentale ne se limite pas aux salles de classe. Le radiotélescope de \textbf{Nkolbisson} (Yaoundé), bien qu'utilisant des ondes radio (grande longueur d'onde), repose sur le même principe de collecte et de focalisation qu'une lunette optique : une grande « ouverture » (antenne parabolique) joue le rôle de l'objectif pour capter le signal d'un satellite (ex: \textbf{Sat-1} ou liaisons \textbf{Intelsat}) et un « oculaire » électronique (récepteur LNB) placé au foyer traite le signal. Pour l'observation visuelle, les clubs d'astronomie de Douala, Yaoundé ou Ngaoundéré utilisent des lunettes de type \emph{réfracteur} (objectif lentille) ou \emph{réflecteur} (objectif miroir) pour suivre les éclipses, les conjonctions planétaires ou la Station Spatiale Internationale (ISS). Le choix de l'oculaire (distance focale, champ apparent, dégagement oculaire) est crucial : un oculaire de $\SI{20}{mm}$ (grossissement modéré, grand champ) est idéal pour la Lune ou les amas ouverts (Pléiades), tandis qu'un $\SI{6}{mm}$ (fort grossissement) sert aux détails planétaires (bandes de Jupiter, anneaux de Saturne) si la turbulence atmosphérique (\emph{seeing}) le permet.
\end{savaistu}

\newpage

\coursec{Corrigés des exercices}

\begin{corrige}[Exercice 1 — QCM : la lunette astronomique]
\begin{enumerate}
\item \textbf{Réponse b)} : l'objectif est une lentille \textbf{convergente de grande distance focale}. Il doit collecter la lumière d'un objet très éloigné et former une image intermédiaire la plus grande possible ($A_1B_1 \approx f'_1\theta$), d'où une grande distance focale.
\item \textbf{Réponse d)} : les trois affirmations sont vraies. Une lunette afocale est caractérisée par la coïncidence de $F'_1$ (foyer image de l'objectif) avec $F_2$ (foyer objet de l'oculaire) ; l'image intermédiaire étant au foyer objet de l'oculaire, l'image définitive est rejetée à l'infini, et l'œil emmétrope l'observe sans accommoder.
\item \textbf{Réponse c)} : $G = \dfrac{f'_1}{f'_2}$. Pour grossir davantage, il faut un objectif de grande focale et un oculaire de courte focale.
\item \textbf{Réponse b)} : l'image définitive est \textbf{renversée et virtuelle}. L'objectif donne une image réelle renversée $A_1B_1$ ; l'oculaire, utilisé comme loupe, donne de cette image intermédiaire une image virtuelle (à l'infini en réglage afocal), toujours renversée par rapport à l'objet.
\end{enumerate}
\end{corrige}

\begin{corrige}[Exercice 2 — Vrai ou faux ? Justifier]
\begin{enumerate}
\item \textbf{Faux.} Comme $G = \dfrac{f'_1}{f'_2}$, le grossissement \emph{augmente} lorsque la distance focale de l'oculaire \emph{diminue}. Pour augmenter $G$, il faut donc un oculaire de distance focale \textbf{plus petite} (ou un objectif de plus grande focale).
\item \textbf{Vrai.} C'est la définition même d'un système afocal : tout faisceau incident parallèle (objet à l'infini) émerge en faisceau parallèle (image à l'infini). Attention toutefois : si le faisceau incident est incliné de $\theta$ sur l'axe, le faisceau émergent est parallèle mais incliné de $\theta' = G\theta \neq 0$ ; il n'est parallèle \emph{à l'axe} que si le faisceau incident l'est.
\item \textbf{Faux (piège de vocabulaire).} La \emph{vergence} s'exprime en dioptries. Une lunette \textbf{afocale} a une vergence totale \textbf{nulle} : $C = C_1 + C_2 - d\,C_1C_2 = 0$ lorsque $d = f'_1 + f'_2$. Ce qui caractérise une lunette, c'est son \textbf{grossissement} $G$, grandeur sans unité.
\item \textbf{Vrai.} En configuration afocale, $F'_1$ et $F_2$ sont confondus, donc la distance objectif--oculaire vaut :
\[ L = \overline{O_1F'_1} + \overline{F_2O_2} = f'_1 + f'_2. \]
\end{enumerate}
\end{corrige}

\begin{corrige}[Exercice 3 — Définitions et conversions]
\begin{enumerate}
\item \textbf{Définitions.}
\begin{itemize}
\item \cle{Objectif} : lentille convergente de grande distance focale, placée du côté de l'objet ; elle forme de l'objet éloigné une image réelle, renversée, dans son plan focal image.
\item \cle{Oculaire} : lentille convergente de courte distance focale, placée du côté de l'œil ; il joue le rôle de loupe pour observer l'image intermédiaire.
\item \cle{Cercle oculaire} : image de la monture de l'objectif (la pupille d'entrée) donnée par l'oculaire ; c'est le cercle de plus petit diamètre traversé par toute la lumière émergente, où l'on place idéalement la pupille de l'œil.
\item \cle{Grossissement} : rapport $G = \dfrac{\theta'}{\theta}$ entre l'angle sous lequel on voit l'image à travers l'instrument et l'angle sous lequel on voit l'objet à l'œil nu.
\end{itemize}
\item Diamètre apparent de la Lune : $\theta = 30'$.
\[ \theta = 30 \times 60'' = 1800''. \]
En degrés : $\theta = \dfrac{30}{60} = 0{,}5°$. En radians :
\[ \theta = 0{,}5 \times \dfrac{\pi}{180} = \dfrac{\pi}{360} \approx \SI{8,73e-3}{rad}. \]
\item Pour une lunette afocale : $G = \dfrac{\theta'}{\theta} = \dfrac{f'_1}{f'_2}$. La valeur $G = 25$ signifie que l'angle apparent de l'astre observé est \textbf{multiplié par 25} : la Lune, vue à l'œil nu sous $0{,}5°$, apparaît sous $\theta' = 25 \times 0{,}5 = 12{,}5°$ à travers la lunette.
\end{enumerate}
\end{corrige}

\begin{corrige}[Exercice 4 — Calculs directs]
Données : $f'_1 = \SI{80}{cm}$, $f'_2 = \SI{2}{cm}$, lunette afocale.
\begin{enumerate}
\item Longueur de la lunette :
\[ L = f'_1 + f'_2 = 80 + 2 = \SI{82}{cm}. \]
\item Grossissement :
\[ G = \dfrac{f'_1}{f'_2} = \dfrac{80}{2} = 40. \]
\item Vergences (distances focales en mètres : $f'_1 = \SI{0,80}{m}$, $f'_2 = \SI{0,02}{m}$) :
\[ C_1 = \dfrac{1}{f'_1} = \dfrac{1}{0{,}80} = 1{,}25\ \delta \qquad ; \qquad C_2 = \dfrac{1}{f'_2} = \dfrac{1}{0{,}02} = 50\ \delta. \]
La vergence totale d'un système de deux lentilles séparées de $d$ est $C = C_1 + C_2 - d\,C_1C_2$. Avec $d = L = f'_1 + f'_2 = \SI{0,82}{m}$ :
\[ C = 1{,}25 + 50 - 0{,}82 \times 1{,}25 \times 50 = 51{,}25 - 51{,}25 = 0\ \delta. \]
Une vergence nulle signifie que le système est afocal : il ne converge ni ne diverge un faisceau parallèle, il le transforme en un autre faisceau parallèle.
\item Angle observé : $\theta' = G\,\theta = 40 \times 40'' = 1600''$.
En minutes d'arc : $\theta' = \dfrac{1600}{60} \approx 26{,}7'$, soit environ $27'$ d'arc (comparable au diamètre apparent de la Lune à l'œil nu !).
\end{enumerate}
\end{corrige}

\begin{corrige}[Exercice 5 — Construction graphique d'une lunette afocale]
\begin{enumerate}
\item Longueur : $L = f'_1 + f'_2 = 50 + 5 = \SI{55}{cm}$. Grossissement : $G = \dfrac{f'_1}{f'_2} = \dfrac{50}{5} = 10$.
\item À l'échelle $1/5$ : $f'_1$ est représentée par $\SI{10}{cm}$, $f'_2$ par $\SI{1}{cm}$, $L$ par $\SI{11}{cm}$. On trace l'axe optique horizontal, on place $O_1$, puis $O_2$ à $\SI{11}{cm}$ à droite. On place $F'_1$ à $\SI{10}{cm}$ de $O_1$ et $F_2$ à $\SI{1}{cm}$ à gauche de $O_2$ : les deux points sont \textbf{confondus}, ce qui traduit l'afocalité.
\item Protocole de tracé du faisceau incliné de $\theta = 2°$ :
\begin{itemize}
\item Le rayon passant par le centre optique $O_1$ n'est \textbf{pas dévié} : il coupe le plan focal image de $L_1$ en un point $B_1$.
\item Les deux autres rayons, parallèles au premier, émergent de $L_1$ en passant par $B_1$ (propriété du plan focal image : un faisceau parallèle converge en un point du plan focal image). Pour tracer l'un d'eux commodément, on utilise le rayon auxiliaire passant par $O_1$.
\item L'image intermédiaire $B_1$ est située dans le plan contenant $F'_1 = F_2$, donc au foyer objet de l'oculaire : tout rayon issu de $B_1$ émerge de $L_2$ \textbf{parallèle} à la direction $(O_2B_1)$ (le rayon $B_1O_2$ n'est pas dévié et fixe la direction d'émergence).
\item Le faisceau émergent est donc parallèle, incliné de $\theta'$ sur l'axe.
\end{itemize}
\item Mesure graphique : on mesure $\theta' \approx 20°$ (au rapporteur d'angles). D'où :
\[ G_{\text{graphique}} = \dfrac{\theta'}{\theta} \approx \dfrac{20°}{2°} = 10, \]
en accord avec le calcul de la question 1 (écarts possibles de quelques pourcents dus à la précision du tracé).
\end{enumerate}
\textbf{Point de vigilance :} les angles étant ici « grands » ($20°$), la relation exacte fait intervenir les tangentes ; la mesure graphique reste néanmoins cohérente car le tracé géométrique est exact.
\end{corrige}

\begin{corrige}[Exercice 6 — Le choix de l'oculaire]
\begin{enumerate}
\item On veut $G = 50$ avec $f'_1 = \SI{1}{m} = \SI{1000}{mm}$ :
\[ G = \dfrac{f'_1}{f'_2} \quad \Longrightarrow \quad f'_2 = \dfrac{f'_1}{G} = \dfrac{1000}{50} = \SI{20}{mm} = \SI{2}{cm}. \]
\item Longueur en afocal : $L = f'_1 + f'_2 = 1000 + 20 = \SI{1020}{mm} = \SI{1,02}{m}$.
\item Grossissements obtenus avec chaque oculaire :
\begin{center}
\begin{tabular}{|c|c|c|}
\hline
\rowcolor{popblueL} Oculaire $f'_2$ & $G = f'_1/f'_2$ & Commentaire \\
\hline
$\SI{25}{mm}$ & $1000/25 = 40$ & grossissement insuffisant \\
\hline
$\SI{20}{mm}$ & $1000/20 = 50$ & exactement l'objectif visé \\
\hline
$\SI{10}{mm}$ & $1000/10 = 100$ & trop fort (voir question 4) \\
\hline
\end{tabular}
\end{center}
Il choisit l'oculaire de $\SI{20}{mm}$, qui donne exactement $G = 50$.
\item Inconvénients d'un grossissement trop élevé (au choix) : le champ de vision devient très étroit (difficile de pointer et de suivre l'astre) ; l'image est plus sombre (la lumière est « étalée ») ; les tremblements de l'instrument et la turbulence atmosphérique sont amplifiés, rendant l'image instable et floue.
\end{enumerate}
\end{corrige}

\begin{corrige}[Exercice 7 — Démonstration du grossissement]
\begin{enumerate}
\item Schéma de principe :
\begin{popfigure}
\begin{tikzpicture}[scale=0.9]
\draw[->,popmuted] (-0.5,0) -- (11.5,0) node[below left] {axe optique};
% Objectif L1
\draw[popblue,thick,<->,>=Stealth] (2,-1.6) -- (2,1.6);
\node[popblue,above] at (2,1.6) {$L_1$ (objectif)};
\node[below] at (2,0) {$O_1$};
% Oculaire L2
\draw[poporange,thick,<->,>=Stealth] (9,-1.2) -- (9,1.2);
\node[poporange,above] at (9,1.2) {$L_2$ (oculaire)};
\node[below] at (9,0) {$O_2$};
% Foyers
\fill (7,0) circle (1.5pt) node[below] {$F'_1 \equiv F_2$};
\fill (8,0) circle (0pt);
\draw[popmuted,dashed] (7,-1.4) -- (7,1.4) node[above,black] {\footnotesize plan focal};
% Image intermédiaire
\draw[->,popgreen,thick] (7,0) -- (7,-0.9) node[midway,right] {$B_1$};
\node[popgreen] at (7.3,0.15) {$A_1$};
% Rayons incidents
\draw[popink,thick,->] (-0.3,1.35) -- (2,0.75);
\draw[popink,thick,->] (-0.3,0.85) -- (2,0.25);
\node[popink] at (0.2,1.6) {$\theta$};
\draw[popink] (1.1,0.62) arc[start angle=150,end angle=180,radius=0.5];
% Rayons entre les lentilles
\draw[popink,thick] (2,0.75) -- (7,-0.9);
\draw[popink,thick] (2,0.25) -- (7,-0.9);
% Rayons émergents parallèles
\draw[popink,thick,->] (9,-0.5) -- (11.3,0.35);
\draw[popink,thick,->] (9,-0.9) -- (11.3,-0.05);
\node[popink] at (10.6,0.75) {$\theta'$};
\draw[popink] (10.2,0.14) arc[start angle=20,end angle=0,radius=1.1];
\end{tikzpicture}
\legende{Lunette afocale : l'image intermédiaire $A_1B_1$ se forme dans le plan focal image de l'objectif, confondu avec le plan focal objet de l'oculaire.}
\end{popfigure}
L'objet $AB$ étant à l'infini, l'image intermédiaire $A_1B_1$ se forme dans le \textbf{plan focal image de l'objectif} (plan passant par $F'_1$, perpendiculaire à l'axe).
\item Le rayon issu de $B$ (à l'infini, direction $\theta$) passant par $O_1$ n'est pas dévié ; il coupe le plan focal image en $B_1$. Dans le triangle $O_1F'_1B_1$ :
\[ \tan\theta = \dfrac{A_1B_1}{f'_1}. \]
Comme $\theta$ est petit (objets célestes), on utilise l'\textbf{approximation des petits angles} $\tan\theta \approx \theta$ (en radians) :
\[ A_1B_1 \approx f'_1\,\theta. \]
\item L'image intermédiaire $A_1B_1$ est au foyer objet $F_2$ de l'oculaire : l'image définitive est rejetée à l'infini et vue sous l'angle $\theta'$ tel que, dans le triangle $O_2F_2B_1$ :
\[ \tan\theta' = \dfrac{A_1B_1}{f'_2} \quad \Longrightarrow \quad \theta' \approx \dfrac{A_1B_1}{f'_2} \quad (\text{petits angles}). \]
\item En reportant $A_1B_1 = f'_1\theta$ :
\[ \theta' = \dfrac{f'_1\,\theta}{f'_2} \quad \Longrightarrow \quad G = \dfrac{\theta'}{\theta} = \dfrac{f'_1}{f'_2}. \qquad \text{CQFD.} \]
\item Application numérique : $G = \dfrac{900}{20} = 45$. Pour la Lune ($\theta = 0{,}5°$) :
\[ \theta' = G\,\theta = 45 \times 0{,}5 = 22{,}5°. \]
La Lune apparaît sous un angle 45 fois plus grand, soit $22{,}5°$ : elle remplit largement le champ de l'oculaire.
\end{enumerate}
\textbf{Point de vigilance (Probatoire) :} cette démonstration est exigible. Il faut citer explicitement l'approximation $\tan\theta \approx \theta$ avec $\theta$ en \textbf{radians}, et justifier la position de l'image intermédiaire (objet à l'infini $\Rightarrow$ image au foyer image).
\end{corrige}

\begin{corrige}[Exercice 8 — Mise au point : de la Lune au clocher]
\begin{enumerate}
\item Réglage afocal (Lune à l'infini) : $F'_1 \equiv F_2$, donc
\[ L = f'_1 + f'_2 = 60 + 3 = \SI{63}{cm}. \]
\item Le clocher est à $D = \SI{500}{m}$ : $\overline{O_1A} = -\SI{500}{m}$ (objet réel, sens de la lumière vers la lentille). Relation de conjugaison :
\[ \dfrac{1}{\overline{O_1A_1}} = \dfrac{1}{f'_1} + \dfrac{1}{\overline{O_1A}} = \dfrac{1}{0{,}60} + \dfrac{1}{-500} = 1{,}6667 - 0{,}0020 = 1{,}6647\ \text{m}^{-1}. \]
\[ \overline{O_1A_1} = \dfrac{1}{1{,}6647} \approx \SI{0,6007}{m} = \SI{60,07}{cm}. \]
L'image intermédiaire se forme à $\SI{60,07}{cm}$ de l'objectif, soit $\SI{0,07}{cm} = \SI{0,7}{mm}$ \textbf{au-delà} de $F'_1$.
\item Pour que l'image définitive reste à l'infini, $A_1$ doit coïncider avec le foyer objet $F_2$ de l'oculaire. Il faut donc \textbf{éloigner l'oculaire de l'objectif} de :
\[ \Delta = \overline{O_1A_1} - f'_1 = 60{,}07 - 60 = \SI{0,07}{cm} = \SI{0,7}{mm}. \]
La nouvelle longueur de la lunette est $L' = 60{,}07 + 3 = \SI{63,07}{cm}$.
\item Explication : pour un objet à l'infini, l'image intermédiaire se forme exactement en $F'_1$ ; quand l'objet se rapproche, la relation de conjugaison montre que $\overline{O_1A_1}$ \emph{augmente} : l'image intermédiaire recule au-delà de $F'_1$. Si l'oculaire restait fixe, $A_1B_1$ ne serait plus en $F_2$ et l'image finale ne serait plus à l'infini : l'œil devrait accommoder, voire verrait une image floue. On translate donc l'oculaire vers l'extérieur (tirage allongé) pour ramener $F_2$ sur $A_1$.
\begin{popfigure}
\begin{tikzpicture}[scale=0.85]
\draw[->,popmuted] (-0.5,0) -- (12,0);
\draw[popblue,thick,<->,>=Stealth] (2,-1.4) -- (2,1.4);
\node[popblue,above] at (2,1.4) {$L_1$};
\fill (8,0) circle (1.5pt) node[below] {$F'_1$};
\fill (8.07,0) circle (1.5pt) node[above right,popgreen] {$A_1$};
\draw[poporange,thick,<->,>=Stealth] (11.07,-1.1) -- (11.07,1.1);
\node[poporange,above] at (11.07,1.1) {$L_2$ déplacé};
\draw[poporange,dashed,thick,<->,>=Stealth] (11,-1.1) -- (11,1.1);
\draw[<->,>=Stealth,popink] (8,-0.5) -- (8.07,-0.5) node[midway,below] {\footnotesize $\Delta$};
\end{tikzpicture}
\legende{Objet rapproché : l'image intermédiaire $A_1$ recule au-delà de $F'_1$ ; on écarte l'oculaire de $\Delta$ (position pleine) pour que $F_2$ coïncide avec $A_1$.}
\end{popfigure}
\end{enumerate}
\textbf{Point de vigilance :} respecter les conventions de signes dans la relation de conjugaison ($\overline{O_1A} < 0$ pour un objet réel) et convertir toutes les distances dans la même unité (ici le mètre) avant le calcul.
\end{corrige}

\begin{corrige}[Exercice 9 — Observation de Jupiter au club d'astronomie de Douala]
\textbf{Barème indicatif (sur 10) :} 1) 1 pt ; 2) 2 pts ; 3) 1,5 pt ; 4) 3 pts ; 5) 1,5 pt ; 6) 1 pt.
\begin{enumerate}
\item Distance focale de l'objectif :
\[ f'_1 = \dfrac{1}{C_1} = \dfrac{1}{1{,}25} = \SI{0,80}{m} = \SI{80}{cm}. \]
\item Une lunette est dite \cle{afocale} lorsqu'elle donne d'un objet à l'infini une image définitive également à l'infini : le foyer image $F'_1$ de l'objectif coïncide avec le foyer objet $F_2$ de l'oculaire.
\begin{popfigure}
\begin{tikzpicture}[scale=0.9]
\draw[->,popmuted] (-0.5,0) -- (11,0);
\draw[popblue,thick,<->,>=Stealth] (1.5,-1.5) -- (1.5,1.5);
\node[popblue,above] at (1.5,1.5) {$L_1$};
\node[below] at (1.5,0) {$O_1$};
\draw[poporange,thick,<->,>=Stealth] (9.9,-1) -- (9.9,1);
\node[poporange,above] at (9.9,1) {$L_2$};
\node[below] at (9.9,0) {$O_2$};
\fill (9.5,0) circle (1.5pt) node[below] {$F'_1 \equiv F_2$};
\fill (-6.5,0) circle (1.5pt) node[below] {$F_1$};
\fill (10.3,0) circle (1.5pt) node[below right] {$F'_2$};
\draw[<->,>=Stealth,popink] (1.5,-1.9) -- (9.9,-1.9) node[midway,below] {$L = f'_1 + f'_2$};
\end{tikzpicture}
\legende{Schéma de principe de la lunette afocale (échelle non respectée).}
\end{popfigure}
\item Longueur : $L = f'_1 + f'_2 = 80 + 4 = \SI{84}{cm}$. Grossissement :
\[ G = \dfrac{f'_1}{f'_2} = \dfrac{80}{4} = 20. \]
\item 
\begin{enumerate}
\item[a)] Conversion : $\theta = 46'' = \dfrac{46}{3600}° = \dfrac{46}{3600}\times\dfrac{\pi}{180} \approx \SI{2,23e-4}{rad}$.
\item[b)] Diamètre apparent à travers la lunette :
\[ \theta' = G\,\theta = 20 \times 46'' = 920'' \approx 15{,}3'. \]
\item[c)] Taille de l'image intermédiaire :
\[ A_1B_1 = f'_1\,\theta = 0{,}80 \times 2{,}23\times10^{-4} \approx \SI{1,78e-4}{m} \approx \SI{0,18}{mm}. \]
\end{enumerate}
\item Nouvel oculaire $f''_2 = \SI{2}{cm}$ :
\[ G' = \dfrac{f'_1}{f''_2} = \dfrac{80}{2} = 40. \]
La longueur devient $L' = f'_1 + f''_2 = 80 + 2 = \SI{82}{cm}$ : la lunette est \textbf{raccourcie de $\SI{2}{cm}$} (l'oculaire est rapproché de l'objectif).
\item Pour l'astronomie, le renversement de l'image est sans gravité : un astre (planète, étoile, Lune) n'a pas de « haut » ni de « bas » privilégié, et on peut tourner la carte du ciel en conséquence. Pour un paysage terrestre (clocher, animal, bateau), une image renversée est inutilisable : il faut alors un système redresseur (prismes ou lentille supplémentaire), comme dans les jumelles ou les lunettes terrestres.
\end{enumerate}
\textbf{Points de vigilance :} convertir les secondes d'arc en radians \emph{avant} d'appliquer $\tan\theta \approx \theta$ ; exprimer $f'_1$ et $f'_2$ dans la même unité pour $G$ ; citer la condition d'afocalité dans la définition.
\end{corrige}

\begin{corrige}[Exercice 10 — La lunette du collège de Nkongsamba]
\textbf{Barème indicatif (sur 10) :} 1) 1,5 pt ; 2a) 1,5 pt ; 2b) 2 pts ; 2c) 1 pt ; 3) 1,5 pt ; 4) 1,5 pt ; 5) 1 pt.
\begin{enumerate}
\item Le bâtiment est à l'infini : $\dfrac{1}{\overline{O_1A}} \to 0$, donc la relation de conjugaison donne
\[ \dfrac{1}{\overline{O_1A_1}} = \dfrac{1}{f'_1} \quad \Longrightarrow \quad \overline{O_1A_1} = f'_1 = \SI{75}{cm}. \]
L'écran est placé à $\SI{75}{cm}$ de l'objectif, dans son plan focal image. C'est la méthode de l'objet à l'infini pour mesurer une distance focale.
\item 
\begin{enumerate}
\item[a)] Condition d'afocalité : distance objectif--oculaire $L = f'_1 + f'_2 = 75 + 3 = \SI{78}{cm}$. Grossissement : $G = \dfrac{75}{3} = 25$.
\item[b)] Protocole de tracé (échelle $1/5$ : $f'_1 \to \SI{15}{cm}$, $f'_2 \to \SI{0,6}{cm}$, $L \to \SI{15,6}{cm}$) :
\begin{itemize}
\item Tracer l'axe optique ; placer $O_1$, $O_2$ à $\SI{15,6}{cm}$ ; placer $F'_1$ à $\SI{15}{cm}$ de $O_1$ et $F_2$ à $\SI{0,6}{cm}$ à gauche de $O_2$ (points confondus).
\item Tracer deux rayons incidents parallèles inclinés de $\theta = 3°$ ; le rayon passant par $O_1$ n'est pas dévié et coupe le plan focal image en $B_1$ ; le second rayon émerge de $L_1$ en passant par $B_1$.
\item De $B_1$ (situé en $F_2$), tracer le rayon passant par $O_2$ non dévié ; l'autre rayon émerge de $L_2$ parallèle à $(O_2B_1)$ : le faisceau émergent est parallèle.
\end{itemize}
\item[c)] On mesure $\theta' \approx 75°$ ? Non — plus rigoureusement, pour $G = 25$ et $\theta = 3°$, la relation exacte est $\tan\theta' = G\tan\theta$, soit $\tan\theta' = 25 \times 0{,}0524 = 1{,}31$, d'où $\theta' \approx 52{,}6°$. La mesure au rapporteur donne une valeur voisine de $53°$, et $G_{\text{graphique}} = \dfrac{\tan\theta'}{\tan\theta} \approx 25$, en accord avec le calcul. (L'approximation des petits angles $\theta' \approx G\theta$ n'est plus valable ici car $\theta'$ est grand : il faut comparer les \emph{tangentes}.)
\end{enumerate}
\item Diamètre apparent de la Lune : $\theta = 31'$. À travers la lunette :
\[ \theta' = G\,\theta = 25 \times 31' = 775' = \dfrac{775}{60}° \approx 12{,}9°. \]
\item Diamètre de l'image de la Lune sur l'écran placé au foyer image de l'objectif : $\theta = 31' = \dfrac{31}{60}\times\dfrac{\pi}{180} \approx \SI{9,02e-3}{rad}$, donc
\[ A_1B_1 = f'_1\,\theta = 0{,}75 \times 9{,}02\times10^{-3} \approx \SI{6,8e-3}{m} = \SI{6,8}{mm}. \]
\item Un œil emmétrope au repos (sans accommodation) voit net un objet à l'infini. Or la lunette afocale rejette l'image définitive à l'infini : l'œil la reçoit donc comme un objet à l'infini et n'a pas besoin d'accommoder. Avantage : l'observation se fait \textbf{sans fatigue musculaire}, même prolongée, et l'image reste nette quelle que soit la durée d'observation.
\end{enumerate}
\textbf{Points de vigilance :} au 2c), ne pas appliquer $\theta' = G\theta$ lorsque $\theta'$ dépasse une dizaine de degrés ; utiliser les tangentes. Au 4), convertir les minutes d'arc en radians avant d'utiliser $A_1B_1 = f'_1\theta$.
\end{corrige}

\begin{corrige}[Exercice 11 — Safari photo dans la réserve de Waza : choisir son instrument]
\textbf{Barème indicatif (sur 10) :} 1) 1,5 pt ; 2) 1,5 pt ; 3a) 1,5 pt ; 3b) 1,5 pt ; 4) 3 pts ; 5) 1 pt.
\begin{enumerate}
\item Une lunette astronomique comprend :
\begin{itemize}
\item un \cle{objectif} : lentille convergente de grande distance focale, tournée vers l'objet ; il forme de l'objet éloigné une image réelle renversée dans son plan focal image ;
\item un \cle{oculaire} : lentille convergente de courte distance focale, placée du côté de l'œil ; il joue le rôle de loupe et donne de l'image intermédiaire une image définitive virtuelle, renversée, rejetée à l'infini (réglage afocal).
\end{itemize}
\item Distance focale de l'oculaire :
\[ f'_2 = \dfrac{f'_1}{G_1} = \dfrac{96}{60} = \SI{1,6}{cm} = \SI{16}{mm}. \]
Longueur afocale : $L = f'_1 + f'_2 = 96 + 1{,}6 = \SI{97,6}{cm} \approx \SI{1}{m}$. La vérification est conforme à l'énoncé.
\item 
\begin{enumerate}
\item[a)] Diamètre apparent de l'antilope à l'œil nu :
\[ \theta \approx \tan\theta = \dfrac{h}{d} = \dfrac{1{,}2}{600} = \SI{2,0e-3}{rad}. \]
En minutes d'arc : $\theta = 2{,}0\times10^{-3} \times \dfrac{180}{\pi} \approx 0{,}1146°$, soit $\theta \approx 0{,}1146 \times 60 \approx 6{,}9'$.
\item[b)] À travers la lunette : $\theta'_1 = G_1\theta = 60 \times 6{,}9' \approx 413' \approx 6{,}9°$. À travers les jumelles : $\theta'_2 = G_2\theta = 8 \times 6{,}9' \approx 55' \approx 0{,}92°$.
\end{enumerate}
\item \textbf{Texte argumenté (corrigé-type).} Pour ce safari dans la réserve de Waza, je conseille au touriste de choisir les \textbf{jumelles} plutôt que la lunette astronomique, et ce pour au moins trois raisons. D'abord, les jumelles donnent une image \textbf{redressée} grâce à leur système de prismes : observer une antilope la tête en bas, comme le ferait la lunette astronomique, serait à la fois déroutant et inutilisable pour identifier l'animal. Ensuite, leur \textbf{champ de vision est large} : dans la savane, les animaux se déplacent et il faut pouvoir les localiser rapidement puis les suivre ; le champ étroit de la lunette ($G_1 = 60$) rendrait la visée presque impossible sur un animal en mouvement. Troisièmement, les jumelles sont \textbf{légères, compactes et utilisables à main levée}, alors qu'une lunette d'un mètre de long exige un trépied stable : à main levée, un grossissement de 60 amplifie énormément les tremblements et l'image devient inobservable. Certes, le grossissement des jumelles ($G_2 = 8$) est plus faible, mais il est largement suffisant : l'antilope, invisible en détail à l'œil nu ($\theta \approx 7'$), apparaît sous près d'un degré, angle très confortable pour l'œil. Enfin, la luminosité des jumelles à fort diamètre d'objectif est bien adaptée aux observations au lever ou au coucher du soleil, moments où la faune est la plus active. La lunette astronomique resterait pertinente uniquement pour l'observation du ciel nocturne de Waza, exceptionnellement pur.
\item Taille de l'image sur le capteur (antilope quasi à l'infini, image au foyer) :
\[ A'B' = f'\,\theta = 0{,}300 \times 2{,}0\times10^{-3} = \SI{6,0e-4}{m} = \SI{0,6}{mm}. \]
Sur un capteur de format standard (environ $24 \times \SI{36}{mm}$), une image de $\SI{0,6}{mm}$ est minuscule : l'antilope n'occupe qu'une centaine de pixels. Le résultat est \textbf{exploitable mais très limite} : pour un sujet aussi éloigné, il faudrait un téléobjectif de focale bien plus grande (ou se rapprocher de l'animal), ce qui confirme qu'en safari, la distance au sujet compte autant que l'instrument.
\end{enumerate}
\textbf{Points de vigilance :} au 3a), justifier $\tan\theta \approx \theta$ par la petitesse de l'angle ; au 4), la rédaction doit mobiliser au moins trois arguments distincts (redressement, champ, stabilité/encombrement) appuyés sur les valeurs calculées.
\end{corrige}

\newpage

\coursec{Fiche bilan}

\begin{popfigure}
\centering
\begin{center}\tcbox[colback=popblue,colframe=popblue,coltext=white,arc=9pt,boxsep=4pt,left=6pt,right=6pt]{\bfseries\small La lunette astronomique}\end{center}
\vspace{-2pt}
\begin{tcbraster}[raster columns=3,raster equal height=rows,raster column skip=6pt,raster row skip=6pt,enhanced,blankest]
\begin{tcolorbox}[enhanced,colback=poporange,colframe=poporange,coltext=white,boxrule=0.9pt,arc=3pt,left=4pt,right=4pt,top=3pt,bottom=3pt,boxsep=1pt,halign=left]\bfseries\scriptsize Constitution : objectif $L_1$ ($f'_1$ grand) + oculaire $L_2$ ($f'_2$ petit)\end{tcolorbox}
\begin{tcolorbox}[enhanced,colback=popgreen,colframe=popgreen,coltext=white,boxrule=0.9pt,arc=3pt,left=4pt,right=4pt,top=3pt,bottom=3pt,boxsep=1pt,halign=left]\bfseries\scriptsize Lunette afocale : $F'_1 \equiv F_2$, objet et image à l'infini\end{tcolorbox}
\begin{tcolorbox}[enhanced,colback=poppurple,colframe=poppurple,coltext=white,boxrule=0.9pt,arc=3pt,left=4pt,right=4pt,top=3pt,bottom=3pt,boxsep=1pt,halign=left]\bfseries\scriptsize Tracé des rayons : faisceau parallèle $\to$ faisceau parallèle\end{tcolorbox}
\begin{tcolorbox}[enhanced,colback=poppink,colframe=poppink,coltext=white,boxrule=0.9pt,arc=3pt,left=4pt,right=4pt,top=3pt,bottom=3pt,boxsep=1pt,halign=left]\bfseries\scriptsize Grossissement $G=\dfrac{\alpha'}{\alpha}=\dfrac{f'_1}{f'_2}$\end{tcolorbox}
\begin{tcolorbox}[enhanced,colback=popgold,colframe=popgold,coltext=white,boxrule=0.9pt,arc=3pt,left=4pt,right=4pt,top=3pt,bottom=3pt,boxsep=1pt,halign=left]\bfseries\scriptsize Puissance $P=\dfrac{\alpha'}{D}$, en dioptries\end{tcolorbox}
\begin{tcolorbox}[enhanced,colback=popblue!70,colframe=popblue!70,coltext=white,boxrule=0.9pt,arc=3pt,left=4pt,right=4pt,top=3pt,bottom=3pt,boxsep=1pt,halign=left]\bfseries\scriptsize Mise au point : bague de l'oculaire, image nette sans fatigue\end{tcolorbox}
\end{tcbraster}

\legende{Carte mentale de la leçon : la lunette astronomique.}
\end{popfigure}

\begin{bilan}
\textbf{Définitions clés}
\begin{itemize}
  \item \cle{Lunette astronomique} : instrument d'optique formé de deux lentilles convergentes de même axe optique, destiné à observer des objets très éloignés (astres).
  \item \cle{Objectif} $L_1$ : lentille tournée vers l'objet, de grande distance focale $f'_1$ (et de grand diamètre).
  \item \cle{Oculaire} $L_2$ : lentille placée du côté de l'\oe il, de petite distance focale $f'_2$.
  \item \cle{Lunette afocale} : réglage où le foyer image $F'_1$ de l'objectif coïncide avec le foyer objet $F_2$ de l'oculaire ; l'image d'un objet à l'infini est alors rejetée à l'infini (vision sans accommodation).
  \item \cle{Grossissement} $G$ : rapport du diamètre apparent de l'image à celui de l'objet vu à l'\oe il nu.
  \item \cle{Puissance} $P$ : quotient de l'angle $\alpha'$ sous lequel on voit l'image par le diamètre $D$ de l'objet ; $P$ s'exprime en dioptries ($\delta$).
\end{itemize}

\textbf{Formules à connaître par c\oe ur}
\begin{center}
\begin{tabularx}{\linewidth}{|l|X|X|}\hline
\rowcolor{popblueL}
\textbf{Grandeur} & \textbf{Expression} & \textbf{Conditions} \\ \hline
Grossissement & $G=\dfrac{\alpha'}{\alpha}=\dfrac{f'_1}{f'_2}$ & lunette afocale, angles petits (approximation de Gauss) \\ \hline
Puissance & $P=\dfrac{\alpha'}{D}$ & $D$ diamètre (ou taille) de l'objet ; $P$ en $\delta$ si $\alpha'$ en rad et $D$ en m \\ \hline
Distance objectif--oculaire & $O_1O_2 = f'_1 + f'_2$ & lunette afocale (réglage normal) \\ \hline
Angle objet & $\alpha \approx \tan\alpha = \dfrac{D}{d}$ & objet de diamètre $D$ à la distance $d$ \\ \hline
\end{tabularx}
\end{center}

\textbf{Propriétés essentielles}
\begin{itemize}
  \item L'objectif forme de l'objet à l'infini une image réelle, renversée, dans son plan focal image ; l'oculaire joue le rôle de loupe sur cette image intermédiaire.
  \item L'image définitive est \textbf{virtuelle, renversée et grossie} ; l'inversion ne gêne pas l'observation astronomique.
  \item Un faisceau parallèle incident émerge de la lunette en faisceau parallèle (système afocal) : l'\oe il emmétrope observe sans fatigue.
  \item Pour augmenter $G$, on augmente $f'_1$ ou on diminue $f'_2$ (au prix d'un champ plus petit et d'une image moins lumineuse).
\end{itemize}

\textbf{Méthodes express}
\begin{enumerate}
  \item \emph{Calculer $G$} : identifier $f'_1$ (objectif, la plus grande focale) et $f'_2$ (oculaire), puis $G=f'_1/f'_2$ (sans unité).
  \item \emph{Tracer la marche d'un faisceau} : faisceau parallèle incliné $\to$ convergence au plan focal image de $L_1$ $\to$ émergence parallèle après $L_2$ (rayon passant par $O_2$ non dévié).
  \item \emph{Faire la mise au point} : viser un objet lointain, déplacer l'oculaire jusqu'à une image nette ; pour l'infini, $O_1O_2=f'_1+f'_2$.
  \item \emph{Choisir un instrument} : objet lointain $\to$ lunette ($G$ grand) ; objet proche et petit $\to$ loupe ou microscope ; justifier par $G$, la clarté et le champ.
\end{enumerate}
\end{bilan}

\begin{aretenir}[Checklist avant le Probatoire]
\begin{itemize}
  \item[$\square$] Je sais nommer et placer l'objectif et l'oculaire sur un schéma, avec $f'_1 \gg f'_2$.
  \item[$\square$] Je sais définir une lunette afocale ($F'_1 \equiv F_2$) et expliquer son intérêt (image à l'infini, \oe il au repos).
  \item[$\square$] Je sais donner la longueur d'une lunette afocale : $O_1O_2 = f'_1 + f'_2$.
  \item[$\square$] Je sais tracer la marche d'un faisceau parallèle à travers les deux lentilles.
  \item[$\square$] Je sais définir le grossissement $G=\alpha'/\alpha$ et le calculer par $G=f'_1/f'_2$.
  \item[$\square$] Je sais définir la puissance $P=\alpha'/D$ et donner son unité (la dioptrie).
  \item[$\square$] Je sais que l'image définitive est virtuelle, renversée et grossie.
  \item[$\square$] Je sais expliquer comment faire la mise au point avec la bague de l'oculaire.
  \item[$\square$] Je sais choisir et justifier un instrument d'optique adapté à une situation d'observation.
  \item[$\square$] Je vérifie toujours les unités (focales en mètres, angles en radians) et l'homogénéité de mes formules.
\end{itemize}
\end{aretenir}

\findecours

\end{document}
